\documentclass[11pt, a4paper, oneside]{book}

% --- PACCHETTI BASE ---
% --- PACCHETTI BASE ---
\usepackage[italian]{babel}
\usepackage[a4paper, top=2.5cm, bottom=2.5cm, left=2.5cm, right=2.5cm]{geometry}
\usepackage{fontspec} 
\usepackage{amsmath}
\usepackage{amssymb}
\usepackage{xcolor}
\usepackage{hyperref}
\hypersetup{colorlinks=true, linkcolor=blue!50!black, urlcolor=blue!50!black, citecolor=blue!50!black}
\usepackage{float}
\usepackage{graphicx}
\usepackage{booktabs}
\usepackage{array}
\usepackage{tabularx}
\newcolumntype{L}{>{\raggedright\arraybackslash}X}
\renewcommand{\arraystretch}{1.3}
\setlength{\tabcolsep}{4pt}
\setlength{\heavyrulewidth}{0.08em}
\setlength{\lightrulewidth}{0.05em}

% --- PACCHETTI PER GRAFICI ---
\usepackage{tikz}
\usetikzlibrary{arrows.meta}

% --- CONFIGURAZIONE GRAFICA ESEMPI ---
\usepackage[export]{adjustbox}
\usepackage{needspace}

% --- TCOLORBOX ---
\usepackage[many]{tcolorbox}

\tcbset{
	stilebase/.style={
		enhanced,
		breakable, 
		colback=white,
		coltitle=white,
		toptitle=1.2mm,
		bottomtitle=1.2mm,
		boxrule=1pt, 
		separator sign={:},
		before skip=12pt, 
		after skip=12pt,
		fonttitle=\bfseries\large\raggedright
	}
}

% \label nel corpo di un box rimanda al numero del box (non a quello della sezione)
\makeatletter
\newcommand{\tcbsetcurrentlabel}{\protected@edef\@currentlabel{\thetcbcounter}}
\makeatother

% Fabbrica di box per OCP (Open/Closed Principle)
\newcommand{\newmanualbox}[3]{%
    % #1 = environment name
    % #2 = title prefix
    % #3 = additional tcb keys
    \newtcolorbox[auto counter,number within=chapter,list inside=#1]{#1}[1][]{%
        stilebase,
        adjusted title={#2 \thetcbcounter\if\relax\detokenize{##1}\relax\else: ##1\fi},
        before upper={\tcbsetcurrentlabel},
        #3
    }%
}

% Ambiente per definizioni
\newmanualbox{definizione}{Definizione}{colframe=black!80!white}

% Ambiente per esempi
\newmanualbox{esempio}{Esempio}{colframe=black!60!white, boxrule=0.6pt, before upper app={\interlinepenalty=10000\relax}}

% Ambiente per regole
\newmanualbox{regola}{Regola}{colframe=black!80!white}

% Ambiente per errori comuni
\newmanualbox{errore}{Errore}{colframe=black!95!white}

% Ambiente per fisica
\newmanualbox{fisica}{Con la fisica}{colback=black!8!white,colframe=black!50!white,boxrule=0.4pt,fontupper=\small,sharp corners}

% Ambiente per matematica
\newmanualbox{matematica}{Con la matematica}{colback=black!8!white,colframe=black!50!white,boxrule=0.4pt,fontupper=\small,sharp corners}

% Ambiente per osservazioni
\newmanualbox{osservazione}{Osservazione}{colback=black!2!white,colframe=black!40!white,boxrule=0.4pt,fontupper=\small\itshape,sharp corners}

% --- MACRO MUSICALI E SEMANTICHE ---

% Macro native per alterazioni nel testo LaTeX
\newcommand{\diesis}{\ensuremath{\sharp}}
\newcommand{\bemolle}{\ensuremath{\flat}}
\newcommand{\nat}{\ensuremath{\natural}}

% Macro combinate per le note più comuni
\newcommand{\Dodiesis}{Do\diesis{}}
\newcommand{\Dobemolle}{Do\bemolle{}}
\newcommand{\Rediesis}{Re\diesis{}}
\newcommand{\Rebemolle}{Re\bemolle{}}
\newcommand{\Midiesis}{Mi\diesis{}}
\newcommand{\Mibemolle}{Mi\bemolle{}}
\newcommand{\Fadiesis}{Fa\diesis{}}
\newcommand{\Fabemolle}{Fa\bemolle{}}
\newcommand{\Soldiesis}{Sol\diesis{}}
\newcommand{\Solbemolle}{Sol\bemolle{}}
\newcommand{\Ladiesis}{La\diesis{}}
\newcommand{\Labemolle}{La\bemolle{}}
\newcommand{\Sidiesis}{Si\diesis{}}
\newcommand{\Sibemolle}{Si\bemolle{}}

% Macro integrata per l'inclusione di blocchi LilyPond.
\newcommand{\musica}[1]{%
	\begin{center}
		\lily[nofragment]{
            \include "definizioni.ly"
			#1
		}%
	\end{center}
}


\begin{document}
	
	
	\title{Appunti sparsi di Armonia}
	\author{Giulio Leuci\\e vari Large Language Model}
	\date{\today}
	\maketitle
	
	\frontmatter
	\tableofcontents
	
	\mainmatter
	
	% ============================================================
	% CAPITOLI
	% ============================================================
	% ============================================================
% PREFAZIONE (CAPITOLO NON NUMERATO)
% ============================================================
\chapter*{Istruzioni per l'uso}

\markboth{Prefazione}{}

Quello che hai davanti non è un manuale di armonia nel senso tradizionale del termine. Non è un testo di consultazione, né pretende di essere esaustivo. Queste pagine sono, molto più semplicemente, i miei appunti sparsi: il tentativo di mettere ordine nel meraviglioso e a volte caotico mondo della sintassi tonale, così come l'ho appreso io stesso.

Da studente di musica, ma con un passato da fisico e un presente da insegnante di matematica, ho affrontato questa disciplina con una certa "deformazione professionale". Ho cercato schemi, logiche e principi primi, convinto che l'armonia, prima di essere arte, sia una scienza di straordinaria coerenza. Questi appunti, quindi, non seguono un ordine tematico o enciclopedico, ma un percorso di apprendimento progressivo, pensato per costruire la conoscenza un mattone alla volta.

Gli appunti qui trascritti derivano dalle lezioni del \textbf{prof. M° Davide Vallini}, dalla manualistica classica e, lo ammetto, dal pesante contributo dei Grandi Modelli di Linguaggio (LLM)\footnote{In particolare Google Gemini e Claude di Anthropic.}. Proprio per questo, non troverai qui un eserciziario: per quello, ti rimando con convinzione ai manuali "ufficiali" citati in bibliografia. Loro sono le fonti autorevoli; questi sono solo gli appunti di un compagno di viaggio.

Per navigare in queste pagine, ho trovato utile una struttura che tradisce la mia formazione. Ogni concetto è affrontato seguendo tre passi:
Per navigare in queste pagine, ho trovato utile una struttura che tradisce la mia formazione. Ogni concetto è affrontato seguendo questi passi:
\begin{itemize}
	\item \textbf{Definizione.} Come in matematica, ci serve un linguaggio comune per sapere esattamente di cosa stiamo parlando. Le definizioni non si discutono, si imparano.
	\item \textbf{Regola.} Una volta compreso il concetto, la regola ci dice come applicarlo nella pratica della scrittura armonica.
	\item \textbf{Esempio.} Qui vediamo la regola in azione.
	\item \textbf{Errore.} Non un divieto fine a se stesso, ma l'analisi del \textit{perché} qualcosa non funziona. Capire la logica di un errore è il primo passo per evitarlo, e un giorno, forse, per infrangerlo con consapevolezza.
\end{itemize}

A questa struttura si aggiungono, in alcuni punti del manuale, degli approfondimenti racchiusi in box dalla cornice squadrata e dallo sfondo grigio chiaro, intitolati \textbf{"Con la fisica"} e \textbf{"Con la matematica"}. Questi non sono essenziali per lo studio dell'armonia in senso stretto, ma offrono una prospettiva complementare a chi, come me, trova che la comprensione e l'apprezzamento di un fenomeno passi anche attraverso la sua formalizzazione quantitativa mediante le equazioni, i numeri e le formule.

Perché questa scelta? Perché l'armonia tonale non è solo una convenzione culturale o un insieme di regole stilistiche tramandate per tradizione. Essa è profondamente radicata nelle leggi della fisica (la serie degli armonici, i rapporti di frequenza, i battimenti) e può essere descritta con straordinaria precisione attraverso gli strumenti della matematica (l'aritmetica modulare, la teoria degli insiemi, le funzioni ricorsive). Comprendere \textit{perché} una quinta giusta è consonante, \textit{perché} il tritono genera tensione, o \textit{perché} le progressioni armoniche seguono pattern prevedibili, non toglie nulla alla bellezza della musica. Al contrario, ne rivela la natura di linguaggio universale, ponte tra arte e scienza, tra intuizione estetica e rigore logico.

Se questi box ti sembrano superflui o troppo tecnici, saltali pure senza remore: il percorso principale procede in modo del tutto autonomo. Ma se, come me, trovi che "sapere il perché" arricchisca il "sapere il come", allora potresti trovarvi qualche spunto interessante per guardare la musica con occhi diversi.

Ora, lasciami essere brutalmente onesto: leggere questi appunti non ti insegnerà l'armonia. L'armonia non si impara leggendo, si impara scrivendo. Svolgere esercizi è l'unica via. Falli. E, ancora più importante, \textbf{falli correggere dal tuo insegnante}. Un errore che non viene visto è un'occasione di apprendimento sprecata. L'occhio esterno di un maestro vede le incertezze che a noi sfuggono e trasforma uno scarabocchio in una lezione.

Questa centralità dell'errore non è una mia fissazione, ma la convinzione che sia il vero motore della conoscenza. Nessuno lo ha espresso meglio del filosofo Karl Popper:

\begin{quote}
\emph{	Evitare errori è un ideale meschino: se non osiamo affrontare problemi che siano così difficili da rendere l'errore quasi inevitabile, non vi sarà allora sviluppo della conoscenza. In effetti, è dalle nostre teorie più ardite, incluse quelle che sono erronee, che noi impariamo di più.} (K. Popper, 1972)
\end{quote}

Quindi, ti auguro un buono studio, un buon lavoro e, soprattutto, buoni errori.

	% ============================================================
% CAPITOLO A - FONDAMENTI
% ============================================================
\chapter{Fondamenti}
\label{chap:fondamenti}

L'armonia tonale, prima ancora di essere un insieme di regole stilistiche, è un sistema logico basato su fenomeni acustici naturali. Per comprendere la sua grammatica, è indispensabile partire dalla fisica del suono e dalle relazioni matematiche che legano le note tra loro. Questo capitolo esplora i mattoni fondamentali del nostro sistema musicale: gli armonici, gli intervalli e la struttura gerarchica della scala.

\section{La natura del suono}
\label{sec:natura_suono}

Ogni suono che percepiamo come "musicale", ovvero di altezza determinata, non è un'entità singola, ma un fenomeno complesso. Quando una corda di chitarra o la colonna d'aria in un flauto vibrano, non lo fanno solo per la loro intera lunghezza, ma anche per frazioni intere di essa (metà, un terzo, un quarto, e così via). Ognuna di queste vibrazioni secondarie produce un suono più acuto e meno intenso, che si fonde con il suono principale.

\begin{definizione}[La serie degli armonici]
	\label{def:serie_armonici}
	La \textbf{serie degli armonici naturali} è la successione di suoni secondari che accompagnano un suono fondamentale. Questi suoni, detti \textbf{armonici} (o ipertoni), hanno frequenze (percepite dal nostro orecchio come altezza del suono) che sono multipli interi della frequenza del suono fondamentale.
\end{definizione}

L'ordine in cui gli armonici appaiono è universale e immutabile. Se prendiamo come fondamentale un Do, i primi suoni che emergono dalla sua vibrazione sono, in ordine: un Do più acuto, un Sol, un altro Do, un Mi, un altro Sol e un \Sibemolle  (leggermente calante rispetto al sistema musicale moderno).

\begin{esempio}[I primi 7 suoni armonici di Do]
	\label{ex:serie_armonici_do}
	% Questo esempio non è SATB, quindi usa la sintassi standard di LilyPond
	\musica{es_cap01_01}
\end{esempio}

Come si può osservare nell'Esempio \ref{ex:serie_armonici_do}, la natura stessa ci fornisce gli intervalli fondamentali dell'armonia. Dopo l'ottava (il raddoppio della fondamentale), il primo nuovo suono a comparire è la \textbf{quinta giusta} (Do-Sol), seguito dalla \textbf{terza maggiore} (Do-Mi). Questi tre suoni (Do, Mi, Sol), presi insieme, formano l'accordo perfetto maggiore (come scopriremo nel Capitolo \ref{chap:triadi}), la cellula base di tutta la musica tonale. La serie armonica, quindi, non è un'astrazione teorica, ma la giustificazione fisica dell'esistenza e della stabilità dell'accordo maggiore.

È interessante notare che il settimo armonico è un \Sibemolle  "stonato" rispetto al nostro sistema. Questo ci introduce a un concetto importante: il \textbf{temperamento}. Per poter suonare in tutte le tonalità su strumenti a intonazione fissa come il pianoforte o la chitarra, è stato necessario "temperare" gli intervalli, ovvero alterarli leggermente rispetto ai loro rapporti fisici puri. Il sistema oggi in uso, il \textbf{temperamento equabile}, divide l'ottava in dodici semitoni perfettamente uguali, sacrificando la purezza acustica in favore della flessibilità.

Più avanti analizzeremo la sequenza degli armonici prodotti da alcuni particolari gradi della scala musicale, per cui ne diamo adesso un elenco:

\begin{itemize}
	\label{itemize:elenco_armonici}
	
	\item \textbf{Do}: Do, Do, Sol, Do, Mi, Sol, La/\Ladiesis , Do
	
	\item \textbf{\Dodiesis }: \Dodiesis , \Dodiesis , \Soldiesis , \Dodiesis , Fa, \Soldiesis , \Ladiesis /Si, \Dodiesis 
	
	\item \textbf{Re}: Re, Re, La, Re, \Fadiesis , La, Si/Do, Re
	
	\item \textbf{\Rediesis }: \Rediesis , \Rediesis , \Ladiesis , \Rediesis , Sol, \Ladiesis , Do/\Dodiesis , \Rediesis 
	
	\item \textbf{Mi}: Mi, Mi, Si, Mi, \Soldiesis , Si, \Dodiesis /Re, Mi
	
	\item \textbf{Fa}: Fa, Fa, Do, Fa, La, Do, Re/\Rediesis , Fa
	
	\item \textbf{\Fadiesis }: \Fadiesis , \Fadiesis , \Dodiesis , \Fadiesis , \Ladiesis , \Dodiesis , \Rediesis /Mi, \Fadiesis 
	
	\item \textbf{Sol}: Sol, Sol, Re, Sol, Si, Re, Mi/Fa, Sol
	
	\item \textbf{\Soldiesis }: \Soldiesis , \Soldiesis , \Rediesis , \Soldiesis , Do, \Rediesis , Fa/\Fadiesis , \Soldiesis 
	
	\item \textbf{La}: La, La, Mi, La, \Dodiesis , Mi, \Fadiesis /Sol, La
	
	\item \textbf{\Ladiesis }: \Ladiesis , \Ladiesis , Fa, \Ladiesis , Re, Fa, Sol/\Soldiesis , \Ladiesis 
	
	\item \textbf{Si}: Si, Si, \Fadiesis , Si, \Rediesis , \Fadiesis , \Soldiesis /La, Si
	
\end{itemize}


\section{Gli intervalli}
\label{sec:intervalli}

Un intervallo è la distanza tra due suoni. Per misurarlo, usiamo due criteri: uno quantitativo (il numero) e uno qualitativo (la specie).

\begin{definizione}[Nomenclatura degli intervalli]
	\label{def:intervalli}
	Un intervallo si definisce contando il numero di gradi della scala compresi tra le due note (estremi inclusi). La sua specie può essere:
	\begin{itemize}
		\item \textbf{Maggiore (M)} o \textbf{Minore (m)}: per intervalli di seconda, terza, sesta e settima.
		\item \textbf{Giusto (G)}: per intervalli di unisono, quarta, quinta e ottava.
		\item \textbf{Aumentato (Eccedente)} o \textbf{Diminuito}: quando un intervallo maggiore, minore o giusto viene ampliato o ristretto cromaticamente di un semitono.
	\end{itemize}
\end{definizione}

Gli intervalli si classificano anche in base alla loro sonorità: le \textbf{consonanze} producono un effetto di stabilità e riposo, mentre le \textbf{dissonanze} generano tensione e un bisogno di risoluzione.

\subsection{I rivolti e la regola del 9}

Ogni intervallo può essere trasformato invertendo la posizione delle sue due note. Questa operazione, fondamentale sia nella teoria che nella pratica musicale, prende il nome di \textbf{rivolto} (o inversione).

\begin{definizione}[Rivolto di un intervallo]
	\label{def:rivolto_intervallo}
	Il \textbf{rivolto} (o \textbf{inversione}) di un intervallo si ottiene spostando la nota più grave un'ottava sopra (o equivalentemente, la nota più acuta un'ottava sotto). In questo modo, la nota che era al grave diventa acuta e viceversa.
\end{definizione}

Il rivolto non è una semplice trasposizione: cambia la natura stessa dell'intervallo. Una terza maggiore, ad esempio, diventa una sesta minore; una quinta giusta diventa una quarta giusta. Esiste tuttavia una relazione aritmetica precisa che governa questa trasformazione.

\begin{regola}[Aritmetica dei rivolti: la Regola del 9]
	\label{reg:regola_del_9}
	La somma di un intervallo e del suo rivolto è sempre 9. La qualità si trasforma secondo il seguente schema:
	\begin{itemize}
		\item \textbf{Maggiore} diventa \textbf{Minore} (e viceversa).
		\item \textbf{Giusto} rimane \textbf{Giusto}.
		\item \textbf{Aumentato} diventa \textbf{Diminuito} (e viceversa).
	\end{itemize}
\end{regola}

\begin{esempio}[Inversione di una Terza Maggiore]
	\label{ex:inversione_intervallo}
	Una terza maggiore (Do-Mi) si rivolta in una sesta minore (Mi-Do). Aritmeticamente: $3+6=9$.
	\musica{es_cap01_02}
\end{esempio}

Questo principio è fondamentale non solo per la teoria, ma anche per la pratica, poiché la natura di un accordo e la sua funzione possono cambiare drasticamente a seconda di quale intervallo si forma con la nota più grave (il basso).

\subsection{Dalla melodia all'armonia: la simultaneità dei suoni}

Finora abbiamo analizzato gli intervalli come relazioni tra coppie di suoni. Tuttavia, la musica occidentale, almeno a partire dal tardo Medioevo, non si limita a far suonare note una dopo l'altra (dimensione \textbf{orizzontale} o \textbf{melodica}), né a far suonare simultaneamente solo due note. La sua vera ricchezza risiede nella capacità di sovrapporre \textit{simultaneamente} tre o più suoni (dimensione \textbf{verticale} o \textbf{armonica}), creando strutture sonore complesse che hanno una loro logica interna e una loro funzione espressiva.

Questa sovrapposizione non è arbitraria o casuale. Come vedremo, esistono modi privilegiati di combinare i suoni che generano stabilità o tensione, che creano senso di movimento o di riposo. L'elemento base di questa costruzione verticale è l'\textbf{accordo}.


\begin{definizione}[Accordo]
	\label{def:accordo}
	Un \textbf{accordo} è la sovrapposizione simultanea di tre o più suoni. Nella pratica tonale, gli accordi si costruiscono prevalentemente sovrapponendo intervalli di terza. Un accordo di tre note è detto \textbf{triade}, uno di quattro note è detto \textbf{quadriade} (o accordo di settima), e così via.
\end{definizione}

La costruzione per terze sovrapposte non è casuale: come vedremo nel Capitolo \ref{chap:triadi}, essa riflette la struttura naturale della serie armonica e genera combinazioni di suoni che il nostro orecchio percepisce come unitarie e coerenti. Un accordo non è una semplice giustapposizione di note, ma un'entità sonora con una propria identità, funzione e tendenza risolutiva all'interno del sistema tonale.

\section{Il sistema tonale}
\label{sec:sistema_tonale}

Il sistema tonale non è una semplice successione di note, ma una \textbf{struttura gerarchica}. Ogni grado della scala ha un nome e una funzione specifica, creando una rete di tensioni e risoluzioni che guidano il discorso musicale.

\begin{definizione}[Gradi e funzioni della scala]
	\label{def:gradi_scala}
	I gradi della scala hanno nomi che ne descrivono la funzione (l'argomento sarà approfondito nel Capitolo \ref{chap:funzioni_armoniche}) e vengono denotati con numeri romani:
	\begin{itemize}
		\item \textbf{I grado (tonica):} il centro di gravità, il punto di massimo riposo.
		\item \textbf{V grado (dominante):} il polo di massima tensione, genera una forte aspettativa verso la tonica.
		\item \textbf{VII grado (sensibile):} così chiamata perché dista solo un semitono dalla tonica e "tende" a risolvere su di essa. È il motore principale della tensione della dominante.
		\item \textbf{IV grado (sottodominante):} crea una tensione più debole, "plagale", e spesso precede la dominante.
	\end{itemize}
	Gli altri gradi (II sopratonica, III mediante, VI sopradominante) hanno funzioni secondarie, spesso di preparazione a quelli principali.
\end{definizione}

La differenza principale tra il \textbf{modo maggiore} e il \textbf{modo minore} risiede nella diversa sequenza di toni e semitoni. In particolare, la distanza tra il I e il III grado (l'intervallo di terza) determina il "colore" del modo. Un'altra asimmetria cruciale riguarda la sensibile: nel modo maggiore è naturalmente presente, mentre nel modo minore naturale il VII grado è un tono sotto la tonica (si chiama infatti sottotonica). Per creare la necessaria tensione, questo grado viene quasi sempre innalzato cromaticamente, dando origine alla \textbf{scala minore armonica}.

 L'innalzamento di un semitono cromatico della sottotonica crea sì una sensibile, ma introduce anche un intervallo "problematico" di seconda aumentata tra sopradominante e sensibile che la melodia tende a evitare.

\section{Aritmetica e geometria tonale}
\label{sec:aritmetica_geometria}

Le relazioni tra le tonalità non sono casuali, ma seguono uno schema geometrico e matematico preciso, visualizzabile attraverso il \textbf{circolo delle quinte}.

	\begin{figure}[H]
	\centering
	% SCALA RIDOTTA per adattarsi alla pagina
	\begin{tikzpicture}[scale=3.6] 
		% --- Cerchi concentrici ---
		\draw[thick, gray!80] (0,0) circle (1.5cm);
		\draw[thick, gray!80] (0,0) circle (1.1cm);
		
		% --- Tonalità maggiori (cerchio esterno) ---
		\node at (90:1.3cm) {\textbf{Do}};
		\node at (60:1.3cm) {\textbf{Sol}};
		\node at (30:1.3cm) {\textbf{Re}};
		\node at (0:1.3cm) {\textbf{La}};
		\node at (330:1.3cm) {\textbf{Mi}};
		\node at (300:1.3cm) {\textbf{Si}};
		\node at (270:1.3cm) {\textbf{\Fadiesis /\Solbemolle }};
		\node at (240:1.3cm) {\textbf{\Rebemolle }};
		\node at (210:1.3cm) {\textbf{\Labemolle }};
		\node at (180:1.3cm) {\textbf{\Mibemolle }};
		\node at (150:1.3cm) {\textbf{\Sibemolle }};
		\node at (120:1.3cm) {\textbf{Fa}};
		
		% --- Tonalità minori (cerchio interno) ---
		\node[scale=0.85] at (90:0.8cm) {la m};
		\node[scale=0.85] at (60:0.8cm) {mi m};
		\node[scale=0.85] at (30:0.8cm) {si m};
		\node[scale=0.85] at (0:0.8cm) {fa\sharp\ m};
		\node[scale=0.85] at (330:0.8cm) {do\sharp\ m};
		\node[scale=0.85] at (300:0.8cm) {sol\sharp\ m};
		\node[scale=0.85] at (270:0.8cm) {re\sharp/mi\flat\ m};
		\node[scale=0.85] at (240:0.8cm) {si\flat\ m};
		\node[scale=0.85] at (210:0.8cm) {fa m};
		\node[scale=0.85] at (180:0.8cm) {do m};
		\node[scale=0.85] at (150:0.8cm) {sol m};
		\node[scale=0.85] at (120:0.8cm) {re m};
		
		% --- Elenco esplicito delle alterazioni (anello più esterno) ---
		% Raggi leggermente ridotti per un layout più compatto
		\node[gray, scale=0.7] at (90:1.75cm) {Nessuna};
		\node[gray, scale=0.7] at (60:1.75cm) {\Fadiesis };
		\node[gray, scale=0.7, align=center] at (30:1.75cm) {\Fadiesis , \Dodiesis };
		\node[gray, scale=0.7, align=center] at (0:1.75cm) {\Fadiesis , \Dodiesis , \Soldiesis };
		\node[gray, scale=0.7, align=center] at (330:1.75cm) {\Fadiesis , \Dodiesis , \Soldiesis , \Rediesis };
		\node[gray, scale=0.7, align=center] at (300:1.75cm) {\Fadiesis , \Dodiesis , \Soldiesis , \Rediesis , \Ladiesis };
		\node[gray, scale=0.65, align=center] at (270:1.85cm) {%
			$\substack{\text{\Fadiesis , \Dodiesis , \Soldiesis , \Rediesis , \Ladiesis , \Midiesis } \\ \text{\Sibemolle , \Mibemolle , \Labemolle , \Rebemolle , \Solbemolle , \Dobemolle }}$};
		\node[gray, scale=0.7, align=center] at (240:1.75cm) {\Sibemolle , \Mibemolle , \Labemolle , \Rebemolle , \Solbemolle };
		\node[gray, scale=0.7, align=center] at (210:1.75cm) {\Sibemolle , \Mibemolle , \Labemolle , \Rebemolle };
		\node[gray, scale=0.7, align=center] at (180:1.75cm) {\Sibemolle , \Mibemolle , \Labemolle };
		\node[gray, scale=0.7, align=center] at (150:1.75cm) {\Sibemolle , \Mibemolle };
		\node[gray, scale=0.7] at (120:1.75cm) {\Sibemolle };
		
		% --- Frecce indicatrici (posizionate più vicino al cerchio) ---
		\draw[->, very thick, red!70, line width=1.2pt] (100:2.1cm) arc (100:20:2.1cm);
		\node[red!70] at (60:2.35cm) {\textbf{QUINTE} $\rightarrow$};
		
		\draw[<-, very thick, blue!70, line width=1.2pt] (140:2.1cm) arc (140:220:2.1cm);
		\node[blue!70] at (180:2.35cm) {$\leftarrow$ \textbf{QUARTE}};
		
	\end{tikzpicture}
	\caption{Il circolo delle quinte e delle quarte con elenco esplicito delle alterazioni.}
	\label{fig:circolo_quinte_migliorato_compatto}
\end{figure}

Ogni passo sul circolo corrisponde a un salto di quinta giusta ascendente (o quarta giusta discendente). Questo schema non è solo una rappresentazione grafica, ma riflette le leggi fisiche e acustiche che abbiamo discusso: le tonalità più vicine sul circolo condividono più note e hanno più armonici in comune, e sono quindi percepite come più "vicine" anche dall'orecchio.

% ============================================================
% INTEGRAZIONE SCIENTIFICA - FISICA
% ============================================================
\begin{fisica}[La serie degli armonici come fondamento fisico della tonalità]
	\label{phys:serie_armonici_fondamento}
	
	Per comprendere perché un suono musicale non è una singola frequenza, dobbiamo partire dall'equazione delle \textbf{onde stazionarie}. Quando una corda vibrante o una colonna d'aria sono vincolate ai loro estremi, possono vibrare solo secondo particolari modi, chiamati \textbf{modi normali di vibrazione}.
	
	La condizione fisica per questi modi è che la lunghezza della corda $L$ deve contenere un numero intero di mezze lunghezze d'onda. Questo si esprime matematicamente come:
	\[
	L = n \cdot \frac{\lambda_n}{2} \quad \text{con } n = 1, 2, 3, \ldots
	\]
	
	Usando la relazione fondamentale tra velocità, frequenza e lunghezza d'onda ($v = \lambda \cdot f$), otteniamo le frequenze dei modi di vibrazione possibili:
	\[
	f_n = n \cdot f_1 = n \cdot \frac{v}{2L}
	\]
	
	dove $f_1$ è la \textbf{frequenza fondamentale} (primo armonico) e $f_n$ sono gli \textbf{armonici superiori}. Questa equazione ci dice che:
	\begin{itemize}
		\item Il secondo armonico ha frequenza $2f_1$ (un'ottava sopra la fondamentale)
		\item Il terzo armonico ha frequenza $3f_1$ (una quinta giusta sopra l'ottava)
		\item Il quarto armonico ha frequenza $4f_1$ (due ottave sopra)
		\item Il quinto armonico ha frequenza $5f_1$ (una terza maggiore sopra le due ottave)
	\end{itemize}
	
	\textbf{Rapporti di frequenza e intervalli musicali}
	
	Analizziamo ora i rapporti di frequenza tra questi armonici e la fondamentale:
	\begin{itemize}
		\item \textbf{Ottava}: rapporto $\frac{f_2}{f_1} = \frac{2:1}{1} = 2:1$
		\item \textbf{Quinta giusta}: rapporto $\frac{f_3}{f_2} = \frac{3f_1}{2f_1} = 3:2$
		\item \textbf{Terza maggiore}: rapporto $\frac{f_5}{f_4} = \frac{5f_1}{4f_1} = 5:4$
	\end{itemize}
	
	Questi rapporti semplici tra numeri interi piccoli (2:1, 3:2, 5:4) sono precisamente quelli che il nostro orecchio percepisce come i più \textbf{consonanti} e \textbf{stabili}. La fisica delle onde stazionarie giustifica quindi perché l'accordo maggiore (costruito con questi intervalli) sia la base naturale del sistema tonale: \textit{è insito nella natura stessa della vibrazione}.
	
	\textbf{Conseguenze per la teoria armonica}
	
	La serie armonica spiega perché:
	\begin{enumerate}
		\item La quinta giusta (3:2) appare prima della terza maggiore (5:4) nella serie: la quinta è un intervallo più "fondamentale"
		\item L'accordo maggiore (Do-Mi-Sol) ha rapporti approssimativamente 4:5:6, molto vicini a multipli consecutivi
		\item Gli intervalli con rapporti più complessi (come il tritono) sono percepiti come più dissonanti
	\end{enumerate}
\end{fisica}

% ============================================================
% INTEGRAZIONE SCIENTIFICA - MATEMATICA
% ============================================================
\begin{matematica}[L'aritmetica modulare come "orologio" del sistema tonale]
	\label{math:aritmetica_modulare}
	
	Il sistema tonale occidentale divide l'ottava in 12 semitoni, creando una struttura che si ripete identica a ogni ottava. Questa natura \textbf{ciclica} può essere descritta in modo elegante usando l'\textbf{aritmetica modulare}, lo stesso tipo di matematica che usiamo per leggere l'orologio.
	
	\textbf{Definizione dell'insieme $\mathbb{Z}_{12}$}
	
	Associamo a ogni nota una classe di altezza rappresentata da un numero da 0 a 11:
	\[
	\begin{array}{cccccccccccc}
	\text{Do} & \text{\Dodiesis } & \text{Re} & \text{\Rediesis } & \text{Mi} & \text{Fa} & \text{\Fadiesis } & \text{Sol} & \text{\Soldiesis } & \text{La} & \text{\Ladiesis } & \text{Si} \\
	0 & 1 & 2 & 3 & 4 & 5 & 6 & 7 & 8 & 9 & 10 & 11
	\end{array}
	\]
	
	L'insieme $\mathbb{Z}_{12} = \{0, 1, 2, \ldots, 11\}$ con l'\textbf{operazione di addizione modulo 12} forma quello che in matematica si chiama un \textbf{gruppo ciclico}.
	
	\textbf{Operazioni in $\mathbb{Z}_{12}$}
	
	L'addizione modulo 12 si comporta come l'aritmetica dell'orologio: quando superiamo 11, "torniamo indietro" a 0. Formalmente:
	\[
	a +_{12} b = (a + b) \mod 12
	\]
	
	\textbf{Esempio pratico}: Calcolo di intervalli
	
	\begin{itemize}
		\item \textbf{Quinta giusta ascendente}: un intervallo di quinta corrisponde a 7 semitoni.\\
		Da Do (0) a Sol: $0 + 7 = 7$ (Sol)\\
		Da Mi (4) a Si: $4 + 7 = 11$ (Si)\\
		Da Sol (7) a Re: $7 + 7 = 14 \equiv 2 \pmod{12}$ (Re)
		
		\item \textbf{Trasposizione}: Trasporre un pezzo di una terza maggiore (4 semitoni) significa sommare 4 a ogni nota.\\
		Accordo di Do maggiore: $(0, 4, 7)$ diventa $(4, 8, 11)$ = Mi maggiore
	\end{itemize}
	
	\textbf{Il circolo delle quinte come sequenza algoritmica}
	
	Il circolo delle quinte può essere generato da una semplice \textbf{funzione ricorsiva}:
	\[
	x_{n+1} = (x_n + 7) \mod 12
	\]
	
	Partendo da Do ($x_0 = 0$), generiamo:
	\begin{align*}
	x_1 &= (0 + 7) \mod 12 = 7 \quad \text{(Sol)} \\
	x_2 &= (7 + 7) \mod 12 = 2 \quad \text{(Re)} \\
	x_3 &= (2 + 7) \mod 12 = 9 \quad \text{(La)} \\
	x_4 &= (9 + 7) \mod 12 = 4 \quad \text{(Mi)} \\
	&\vdots
	\end{align*}
	
	Dopo 12 iterazioni, torniamo a 0 (Do), completando il circolo. Questo dimostra che il circolo delle quinte non è solo una convenzione teorica, ma una proprietà \textbf{matematica intrinseca} del sistema a 12 semitoni.
	
	\textbf{Proprietà interessanti}
	
	\begin{itemize}
		\item L'operazione "+7" genera tutti i 12 elementi di $\mathbb{Z}_{12}$ perché 7 e 12 sono \textbf{coprimi} ($\text{MCD}(7, 12) = 1$)
		\item Il circolo delle quarte si ottiene con l'operazione inversa: $x_{n+1} = (x_n + 5) \mod 12$, equivalente a $x_{n+1} = (x_n - 7) \mod 12$
		\item La simmetria enarmonica (\Fadiesis  = \Solbemolle ) corrisponde al fatto che $6 + 6 = 12 \equiv 0 \pmod{12}$: il tritono divide esattamente l'ottava a metà
	\end{itemize}
	
	\textbf{Applicazione pratica}
	
	Questa formalizzazione matematica non è solo un esercizio teorico: è la base degli \textbf{algoritmi} usati nei software di composizione musicale, nei traspositori automatici e negli strumenti che analizzano progressioni armoniche. Comprendere il sistema tonale come una struttura algebrica ci permette di manipolarlo in modo sistematico e preciso.
\end{matematica}
  % Fondamenti
	% ============================================================
% CAPITOLO B - LE TRIADI
% ============================================================
\chapter{Le triadi}
\label{chap:triadi}

Dopo aver esaminato le fondamenta fisiche e gli intervalli nel Capitolo \ref{chap:fondamenti}, possiamo ora affrontare la costruzione degli accordi. L'elemento base dell'armonia tonale, la sua "molecola" fondamentale, è la \textbf{triade}. Quasi tutta la musica occidentale, dalla polifonia rinascimentale al pop moderno, si basa sulla successione e sulla manipolazione di questi accordi di tre note.

\section{Costruzione della triade}
\label{sec:costruzione_triade}

Come suggerisce il nome, una triade è un accordo composto da tre suoni.

\begin{definizione}[La triade]
	\label{def:triade}
	Una \textbf{triade} è un accordo di tre note costruito sovrapponendo due intervalli di terza a partire da una nota base, detta \textbf{fondamentale}. Le altre due note prendono il nome di \textbf{terza} e \textbf{quinta} dell'accordo, in base all'intervallo che formano con la fondamentale.
\end{definizione}

La combinazione dei due tipi di terza (maggiore e minore) genera le diverse "specie" o "qualità" di triadi, ognuna con un colore e una funzione armonica distinti.

\begin{definizione}[Specie di triadi]
	\begin{itemize}
	\item \textbf{Triade maggiore:} terza maggiore + terza minore. È la triade "naturale", generata direttamente dalla serie armonica.
	\item \textbf{Triade minore:} terza minore + terza maggiore. Stabile, ma con un colore più scuro.
	\item \textbf{Triade diminuita:} terza minore + terza minore. Contiene il tritono (quinta diminuita), intervallo dissonante che genera forte tensione.
	\item \textbf{Triade eccedente (o aumentata):} terza maggiore + terza maggiore. L'intervallo di quinta aumentata conferisce un carattere ambiguo e instabile.
\end{itemize}
\end{definizione}

Di queste quattro specie, solo le prime due (maggiore e minore) sono considerate \textbf{consonanti} e possono fungere da punto di riposo. Le triadi diminuita ed eccedente sono \textbf{dissonanti} e richiedono risoluzione.

\begin{esempio}[Le quattro specie principali di triadi]
	\label{ex:quattro_triadi}
\musica{es_cap02_01}
\end{esempio}

% ============================================================
% INTEGRAZIONE SCIENTIFICA - FISICA
% ============================================================
\begin{fisica}[I rapporti di frequenza e la stabilità degli accordi]
	\label{phys:rapporti_frequenza_triadi}
	
	Per comprendere perché alcune triadi suonano più "stabili" di altre, dobbiamo analizzare i rapporti di frequenza tra le note che le compongono. Abbiamo visto nel Capitolo \ref{chap:fondamenti} che la serie armonica genera naturalmente rapporti semplici tra le frequenze.
	
	\textbf{La triade maggiore e i rapporti semplici}
	
	Consideriamo una triade di Do maggiore (Do-Mi-Sol). Se assumiamo che la fondamentale Do abbia frequenza $f_1$, le frequenze approssimative delle altre note nella scala temperata sono:
	\begin{itemize}
		\item Do (fondamentale): $f_1$
		\item Mi (terza maggiore): $\approx 1.26 \cdot f_1$ (rapporto ideale $5:4 = 1.25$)
		\item Sol (quinta giusta): $\approx 1.50 \cdot f_1$ (rapporto ideale $3:2 = 1.50$)
	\end{itemize}
	
	Nel sistema di \textbf{intonazione naturale} (o "giusta"), questi rapporti sono esatti: $4:5:6$. Questo significa che se Do vibra a 200 Hz, Mi vibra a 250 Hz e Sol a 300 Hz. Questi rapporti tra \textbf{numeri interi piccoli} sono ciò che il nostro orecchio percepisce come massimamente consonante.
	
	\textbf{La triade diminuita e i rapporti complessi}
	
	Confrontiamo ora con una triade diminuita (Do-\Mibemolle -\Solbemolle ). I rapporti approssimativi sono:
	\begin{itemize}
		\item Do: $f_1$
		\item \Mibemolle  (terza minore): $\approx 1.19 \cdot f_1$ (rapporto $6:5$)
		\item \Solbemolle  (quinta diminuita): $\approx 1.41 \cdot f_1$ (rapporto complesso, circa $\sqrt{2}:1$)
	\end{itemize}
	
	Il \textbf{tritono} (Do-\Solbemolle , rapporto $\sqrt{2}:1$ nel temperamento equabile) in particolare non può essere espresso come rapporto di numeri interi piccoli per la nota irrazionalità di $\sqrt{2}$. Se tentassimo di approssimarlo, otterremmo qualcosa come $45:32$ o $64:45$ — numeri molto più grandi e complessi.
	
	\textbf{Il fenomeno dei battimenti acustici}
	
	Quando due frequenze vicine suonano insieme, si genera un fenomeno chiamato \textbf{battimento}. La frequenza del battimento è pari alla differenza tra le due frequenze:
	\[
	f_{\text{battimento}} = |f_2 - f_1|
	\]
	
	Con rapporti semplici (come 3:2), gli armonici delle due note tendono a coincidere o a essere molto distanziati, minimizzando i battimenti. Con rapporti complessi, invece, molti armonici cadono a frequenze vicine ma non coincidenti, generando battimenti multipli che l'orecchio percepisce come "ruvidità" o \textbf{dissonanza}.
	
	\textbf{Tavola comparativa}
	
	\begin{center}
	\begin{tabular}{|l|c|l|}
	\hline
	\textbf{Triade} & \textbf{Rapporti (ideali)} & \textbf{Percezione} \\
	\hline
	Maggiore & $4:5:6$ & Stabile, luminosa \\
	Minore & $10:12:15$ & Stabile, malinconica \\
	Diminuita & include tritono & Instabile, tesa \\
	Eccedente & include 5a aumentata & Ambigua, irrequieta \\
	\hline
	\end{tabular}
	\end{center}
	
	\textbf{Conclusione}
	
	La stabilità di un accordo ha quindi una base fisica oggettiva: triadi con rapporti di frequenza semplici generano meno interferenze acustiche e vengono percepite come più consonanti. Questo spiega perché la triade maggiore e minore sono i pilastri dell'armonia tonale, mentre la triade diminuita richiede quasi sempre risoluzione verso accordi più stabili. L'armonia non è solo convenzione culturale: è profondamente radicata nelle leggi della fisica.
\end{fisica}


\section{Costruzione delle triadi sui gradi della scala}
\label{sec:triadi_sui_gradi}

Ogni grado della scala può essere la fondamentale di una triade. La qualità di questa triade è determinata dalla struttura della scala stessa. Analizziamo il modo maggiore.

\begin{definizione}[Triadi diatoniche nel modo maggiore]
	\label{def:triadi_diatoniche_maggiore}
	Le triadi costruite sui sette gradi della scala maggiore sono:
	\begin{itemize}
		\item \textbf{I, IV, V:} Triadi maggiori (funzione di stabilità)
		\item \textbf{II, III, VI:} Triadi minori (funzione secondaria)
		\item \textbf{VII:} Triade diminuita (funzione di tensione)
	\end{itemize}
\end{definizione}

\begin{esempio}[Le sette triadi diatoniche di Do Maggiore]
	\label{ex:triadi_diatoniche_do}
	\musica{es_cap02_02}
\end{esempio}

La nomenclatura dell'esempio \ref{ex:triadi_diatoniche_do} usa numeri romani maiuscoli per tutti gli accordi (I, II, III, IV, V, VI, VII), specificando tra parentesi la natura dell'accordo (M per maggiore, m per minore, dim per diminuito).

\section{I rivolti della triade}
\label{sec:rivolti_triade}

Finora abbiamo considerato le triadi nel loro \textbf{stato fondamentale}, ovvero con la fondamentale dell'accordo posizionata al basso. Tuttavia, è possibile disporre le note dell'accordo in modo diverso, ponendo al basso la terza o la quinta anziché la fondamentale. Queste diverse disposizioni sono chiamate \textbf{rivolti}.

\subsection{Basso fondamentale e basso reale}
\label{subsec:basso_fondamentale_reale}

Prima di descrivere i rivolti, è necessario chiarire una distinzione terminologica fondamentale.

\begin{definizione}[Basso fondamentale e basso reale]
	\label{def:basso_fondamentale_reale}
	In un accordo si distinguono due concetti:
	\begin{itemize}
		\item \textbf{Basso fondamentale:} la nota che dà il nome all'accordo e sulla quale l'accordo è costruito per terze sovrapposte. È un concetto teorico-funzionale, indipendente dalla disposizione effettiva delle note.
		\item \textbf{Basso reale:} la nota più grave effettivamente suonata, qualunque essa sia.
	\end{itemize}
	
	Se basso reale e fondamentale coincidono, si dice che l'accordo è in stato fondamentale, altrimenti si dice che è in stato di rivolto.
	
\end{definizione}

\begin{esempio}[Basso fondamentale vs. basso reale in Do Maggiore]
	\label{ex:basso_fond_reale}
	Lo stesso accordo di Do maggiore (basso fondamentale: Do) con tre diversi bassi reali.
\musica{es_cap02_03}
	In tutti e tre i casi l'accordo è sempre "Do maggiore" (basso fondamentale = Do), ma il basso reale cambia.
\end{esempio}


\subsection{Il primo rivolto delle triadi (stato di sesta)}
\label{subsec:primo_rivolto}

Il primo rivolto si ottiene ponendo la \textbf{terza} dell'accordo al basso.

\begin{definizione}[Primo rivolto delle triadi]
	\label{def:primo_rivolto}
	Un accordo è in \textbf{primo rivolto} quando la sua terza si trova al basso. Viene indicato con la cifratura \textbf{6} (o talvolta 6/3), poiché rispetto al basso la fondamentale si trova a distanza di sesta e la quinta a distanza di terza.
\end{definizione}

\begin{esempio}[Una triade in primo rivolto]
	\label{ex:primo_rivolto_triadi}
	\musica{es_cap02_04}
\end{esempio}

\subsection{Il secondo rivolto delle triadi (stato di quarta e sesta)}
\label{subsec:secondo_rivolto}

Il secondo rivolto si ottiene ponendo la \textbf{quinta} dell'accordo al basso.

\begin{definizione}[Secondo rivolto delle triadi]
	\label{def:secondo_rivolto}
	Un accordo è in \textbf{secondo rivolto} quando la sua quinta si trova al basso. Viene indicato con la cifratura \textbf{6/4}, poiché rispetto al basso la fondamentale si trova a distanza di quarta e la terza a distanza di sesta.
\end{definizione}

\begin{esempio}[Una triade in secondo rivolto]
	\label{ex:secondo_rivolto_triadi}
	\musica{es_cap02_05}
\end{esempio}

\subsection{Il problema del secondo rivolto}
\label{subsec:problema_secondo_rivolto}

Il secondo rivolto presenta una peculiarità acustica che lo rende profondamente diverso dallo stato fondamentale e dal primo rivolto: pur essendo costituito dalle stesse note di una triade consonante, viene percepito come \textbf{dissonante}. Il suo uso sarà quindi molto limitato.

\begin{fisica}[Il conflitto acustico dell'accordo di quarta e sesta]
	\label{phys:conflitto_6_4}
	
	Quando ascoltiamo un accordo, il nostro orecchio tende a identificare la nota più grave (il basso reale) come fondamentale "percepita". Questo avviene perché la nota al basso genera la propria serie armonica, che domina la percezione complessiva.
	
	\textbf{Il conflitto della quarta}
	
	Nel secondo rivolto, la fondamentale dell'accordo si trova a distanza di \textbf{quarta} dal basso. Questo intervallo è problematico per ragioni fisiche:
	
	\begin{itemize}
		\item La quarta (rapporto $4:3$) non compare tra i primi armonici semplici della serie del basso
		\item Si crea un'interferenza: l'orecchio "sente" il basso come fondamentale, ma la struttura dell'accordo suggerisce una nota diversa
		\item Questa ambiguità genera instabilità percettiva
	\end{itemize}
	
	\textbf{Rappresentazione matematica}
	
	Se il basso suona a frequenza $f_b$, i suoi armonici principali sono $f_b$, $2f_b$, $3f_b$, $4f_b$, $5f_b$... La quarta sopra il basso ha frequenza $f_4 = \frac{4}{3}f_b$. Questo rapporto non appare tra i primi armonici semplici, mentre la quinta ($3:2$) è il secondo rapporto più consonante dopo l'ottava.
	
	\textbf{Conclusione}
	
	L'accordo di 6/4 non è dissonante perché contiene note "sbagliate", ma perché la disposizione delle note crea un conflitto tra ciò che l'orecchio percepisce (il basso come fondamentale) e ciò che la struttura dell'accordo implica (una nota diversa come fondamentale). Per questa ragione, l'accordo di quarta e sesta richiede sempre un trattamento particolare e non può essere usato liberamente come un rivolto qualsiasi.
\end{fisica}


  % Le triadi
	\chapter{Il secondo rivolto: quando e come usarlo}
\label{chap:quarta_sesta}

Abbiamo visto nel Capitolo \ref{chap:triadi} (\ref{subsec:problema_secondo_rivolto}) il problema acustico dell'accordo in secondo rivolto: la quarta rispetto al basso è un intervallo instabile che necessita di un trattamento specifico. In questo capitolo presentiamo i tre casi in cui il secondo rivolto è non solo ammesso, ma fondamentale nella pratica compositiva.

\section{Il 6/4 cadenziale}
\label{sec:cadenziale}

\begin{definizione}[6/4 cadenziale]
    È l'accordo di I6/4 che precede immediatamente la dominante nella formula cadenzale. Le note al soprano e contralto formano una quarta e una sesta rispetto al basso, che è fermo sulla nota della dominante.
\end{definizione}

\begin{regola}
    \begin{itemize}
        \item Le note in quarta e sesta devono essere \textbf{preparate} (tenute o per grado congiunto).
        \item La quarta risolve scendendo alla terza della dominante.
        \item La sesta risolve scendendo alla quinta della dominante.
        \item Metricamente: il 6/4 cade sul tempo \textbf{forte}, il V cade sul tempo successivo.
    \end{itemize}
\end{regola}

\begin{esempio}[Il 6/4 cadenziale in Do maggiore]
    Progressione IV - I6/4 - V - I.
    \musica{
        \partilate
        { \key do \major \time 4/4 }
        { la'2 sol'4 fa' | mi'2 r }
        { fa'2 mi'4 re' | do'2 r }
        { do'2 do'4 si | sol2 r }
        { fa2 sol2 | do2 r }
        \lyricmode { "IV"2 "I6/4 - V" | "I"2 }
        \figuremode { s2 <6 4>4 <5 3> | s2 }
    }

\end{esempio}

\section{Il 6/4 di passaggio}
\label{sec:64_passaggio}

\begin{definizione}[6/4 di passaggio]
    Si forma quando il basso si muove per grado congiunto attraverso la quinta dell'accordo in un passaggio scalare.
\end{definizione}

\begin{regola}
    \begin{itemize}
        \item Il basso deve muoversi per grado congiunto sia in entrata che in uscita.
        \item L'accordo è metricamente \textbf{debole}.
        \item Le voci superiori procedono per grado congiunto o rimangono ferme.
    \end{itemize}
\end{regola}

\begin{esempio}[6/4 di passaggio]
    Passaggio I - V6/4 - I6 su basso scalare Do-Re-Mi: il tenore resta fermo, il soprano procede per grado congiunto.
    \musica{
        \partilate
        { \key do \major \time 3/4 }
        { do''4 si' do'' }
        { mi'4 re' do' }
        { sol4 sol sol }
        { do4 re mi }
        \lyricmode { "I"4 "V6/4" "I6" }
        \figuremode { s4 <6 4>4 <6>4 }
    }
\end{esempio}

\section{Il 6/4 di pedale (o di ornamento)}
\label{sec:64_pedale}

\begin{definizione}[6/4 di pedale]
    Si forma quando le voci superiori si muovono sopra un basso fermo (pedale).
\end{definizione}

\begin{regola}
    Il basso deve essere consonante con l'accordo che precede e con quello che segue. Il 6/4 funge da ornamento transitorio sopra una nota fissa.
\end{regola}

\section{Sintesi degli usi del secondo rivolto}
\label{sec:sintesi_secondo_rivolto}

\begin{center}
\begin{tabular}{|l|l|l|l|}
\hline
Tipo & Posizione metrica & Movimento del basso & Ruolo \\
\hline
Cadenziale & Forte & Fermo sulla dominante & Preparazione della dominante \\
Passaggio & Debole & Grado congiunto & Collegamento melodico \\
Pedale & Indifferente & Nota fissa & Ornamento armonico \\
\hline
\end{tabular}
\end{center}
 % Il secondo rivolto: quando e come usarlo
	% ============================================================
% CAPITOLO D - LA SCRITTURA A QUATTRO PARTI E LA CONDOTTA DELLE VOCI
% (Versione revisionata con integrazione di errori melodici e
%  risoluzione della sensibile da fonti esterne)
% ============================================================
\chapter{La scrittura a quattro parti e la condotta delle voci}
\label{chap:condotta_voci}

Nei capitoli precedenti abbiamo analizzato gli accordi come entità verticali, statiche. La vera arte dell'armonia, tuttavia, risiede nel collegarli tra loro in modo fluido e logico. Questo processo, noto come \textbf{condotta delle voci} (o \textit{voice leading}), trasforma una sequenza di accordi statici in un discorso musicale dinamico e coerente. La scrittura a quattro parti in stile corale (Soprano, Contralto, Tenore, Basso) è il modello di studio per eccellenza per apprendere e padroneggiare questi principi.

\section{Disposizione delle voci}
\label{sec:disposizione_voci}

La scrittura a quattro parti simula un coro misto. A ogni voce è assegnata un'estensione tipica, che è bene rispettare per garantire una buona sonorità e cantabilità.

\begin{definizione}[Estensione delle voci corali (SATB)]
	\label{def:estensione_voci}
	Le estensioni standard approssimative per le quattro voci sono:
	\begin{itemize}
		\item \textbf{Soprano:} Dal Do centrale (Do') al La acuto (La'').
		\item \textbf{Contralto:} Dal Sol sotto il rigo (Sol) al Mi sul quarto spazio (Mi'').
		\item \textbf{Tenore:} Dal Do nel secondo spazio in chiave di basso (Do) al La sopra il rigo (La').
		\item \textbf{Basso:} Dal Fa sulla prima linea in chiave di basso (Fa,) al Mi sopra il rigo (Mi).
	\end{itemize}
\end{definizione}

Oltre all'estensione, è fondamentale considerare la \textbf{disposizione} delle tre voci superiori (Contralto, Tenore, Soprano). La loro distanza reciproca determina la ``geometria'' dell'accordo.

\begin{definizione}[Posizione Stretta vs. Posizione Lata]
	\label{def:posizione_stretta_lata}
	\begin{itemize}
		\item \textbf{Posizione Stretta:} La distanza tra il Soprano e il Tenore è inferiore a un'ottava. Le tre voci superiori sono il più possibile ravvicinate.
		\item \textbf{Posizione Lata (o Larga):} La distanza tra il Soprano e il Tenore è uguale o superiore a un'ottava. Le voci sono più distanziate.
	\end{itemize}
\end{definizione}

\begin{regola}[Disequazioni spaziali tra le voci]
	\label{reg:distanza_voci}
	Per preservare la coesione del tessuto corale, la distanza tra voci adiacenti non deve essere eccessiva.
	\begin{itemize}
		\item Distanza tra Soprano e Contralto $\le$ \textbf{un'ottava}.
		\item Distanza tra Contralto e Tenore $\le$ \textbf{un'ottava}.
		\item La distanza tra Tenore e Basso è più flessibile, ma si consiglia di non superare la dodicesima.
	\end{itemize}
\end{regola}

\section{Le posizioni delle triadi}
\label{sec:posizioni_triadi}

Oltre alla disposizione (stretta o lata), che riguarda la distanza tra le voci, esiste un'altra classificazione fondamentale nella scrittura a quattro parti: la \textbf{posizione} della triade.

\begin{definizione}[Posizione della Triade]
	\label{def:posizione_triade}
	La \textbf{posizione} di una triade è determinata dalla nota dell'accordo che si trova al Soprano (voce più acuta):
	\begin{itemize}
		\item \textbf{Posizione di Ottava (o Fondamentale):} La fondamentale dell'accordo è al Soprano.
		\item \textbf{Posizione di Terza:} La terza dell'accordo è al Soprano.
		\item \textbf{Posizione di Quinta:} La quinta dell'accordo è al Soprano.
	\end{itemize}
\end{definizione}

È fondamentale non confondere le \textit{posizioni} con i \textit{rivolti} (Capitolo \ref{chap:triadi}): i rivolti dipendono dalla nota al \textit{basso}, mentre le posizioni dipendono dalla nota al \textit{soprano}.

\section{Tipologie di moto}
\label{sec:tipologie_moto}

Quando due voci si muovono da un accordo al successivo, creano dei ``vettori melodici'' la cui relazione reciproca definisce il tipo di moto.

\begin{definizione}[Tipi di moto tra le parti]
	\label{def:tipi_moto}
	Esistono tre tipi di moto tra una qualsiasi coppia di voci:
	\begin{itemize}
		\item \textbf{Moto Retto (o Parallelo/Diretto):} Le due voci si muovono nella stessa direzione (entrambe salgono o entrambe scendono). Se mantengono lo stesso intervallo, si parla di \textit{moto parallelo}.
		\item \textbf{Moto Contrario:} Le due voci si muovono in direzioni opposte (una sale e l'altra scende).
		\item \textbf{Moto Obliquo:} Una voce sta ferma, mentre l'altra si muove.
	\end{itemize}
\end{definizione}

Una buona scrittura armonica mostra una netta preferenza per il \textbf{moto contrario}: esso garantisce la massima indipendenza e chiarezza melodica tra le parti.

\section{Il collegamento armonico}
\label{sec:collegamento_armonico}

Il principio guida per un collegamento armonico efficace è la \textbf{legge del minimo sforzo}: le voci devono muoversi il meno possibile per raggiungere l'accordo successivo.

\begin{regola}[Principio di minima azione]
	\label{reg:minima_azione}
	Per collegare due accordi, si seguono due principi gerarchici:
	\begin{enumerate}
		\item \textbf{Note comuni:} Se i due accordi hanno una o più note in comune, queste note dovrebbero, se possibile, rimanere ferme nella stessa voce (moto obliquo).
		\item \textbf{Grado congiunto:} Le altre voci si muovono per il percorso più breve possibile, preferibilmente per grado congiunto (intervalli di seconda).
	\end{enumerate}
\end{regola}

\subsection{Tipologie di collegamento in base alla distanza}

Il collegamento tra due triadi varia in base alla distanza tra le loro fondamentali, che determina il numero di note comuni.

\begin{itemize}
	\item \textbf{Relazione di quinta (o quarta):} Gli accordi hanno \textbf{una nota in comune}. (Esempio: I-IV, I-V, V-I).
	\item \textbf{Relazione di terza (o sesta):} Gli accordi hanno \textbf{due note in comune}. È il collegamento più fluido. (Esempio: I-VI, I-III).
	\item \textbf{Relazione di seconda (o settima):} Gli accordi \textbf{non hanno note in comune}. Richiede moto contrario obbligatorio per evitare parallelismi. (Esempio: IV-V, V-VI).
\end{itemize}

% ============================================================
% SEZIONE NUOVA: ERRORI MELODICI
% (Integrazione da armoniainrete.altervista.org/armonia-lezione-08)
% ============================================================
\section{Errori melodici}
\label{sec:errori_melodici}

Prima di affrontare gli errori armonici, è necessario considerare i vincoli che riguardano il \textbf{profilo melodico di ciascuna voce presa singolarmente}. Una linea vocale deve rispettare alcuni limiti di natura musicale ed esecutiva.

\begin{regola}[Limiti degli intervalli melodici]
	\label{reg:limiti_melodici}
	Nella condotta di una singola voce valgono le seguenti norme:
	\begin{enumerate}
		\item Sono \textbf{proibiti} intervalli melodici superiori alla \textbf{sesta minore} (ad eccezione dell'ottava giusta).
		\item Sono \textbf{proibiti} tutti gli \textbf{intervalli eccedenti e diminuiti}, con le eccezioni indicate di seguito.
		\item La \textbf{settima minore ascendente} è ammessa in due casi speciali: quando sostituisce il movimento di tono discendente tra il V e il IV grado, oppure in funzione modulante.
	\end{enumerate}
\end{regola}

\begin{osservazione}[Eccezione: seconda aumentata nel minore]
	\label{oss:seconda_aumentata}
	La \textbf{seconda aumentata} che si forma tra il VI e il VII grado della scala minore armonica è \textit{tollerata} (se impossibile da evitare), ma va evitata quando possibile. Essa nasce necessariamente dalla presenza simultanea della sopradominante naturale e della sensibile alzata.
\end{osservazione}

\begin{osservazione}[Eccezione: quarta e quinta diminuita]
	\label{oss:quarta_quinta_dim}
	Sono ammessi gli intervalli di \textbf{quarta e quinta diminuita} a condizione che il movimento sia \textbf{discendente} e che la nota successiva risalga di \textbf{semitono diatonico}. Questo schema è tipico del VII grado che funge da appoggiatura.
\end{osservazione}

\begin{errore}[Quarta eccedente melodica]
	\label{err:quarta_eccedente}
	È vietato che una voce si muova per salto di \textbf{quarta eccedente} (tritono). Questo intervallo, storicamente chiamato \textit{diabolus in musica}, risulta difficile da cantare e crea una forte discontinuità melodica.
\end{errore}

\begin{regola}[Eccezione: quarta eccedente con risoluzione]
	\label{reg:quarta_eccedente_eccezione}
	Il salto di quarta eccedente è ammesso quando la nota più acuta \textbf{o} più grave del salto si muove successivamente di \textbf{semitono diatonico} nella direzione opposta, ``risolvendo'' così la tensione del tritono melodico.
\end{regola}

% ============================================================
% FALSA RELAZIONE
% ============================================================
\subsection{La falsa relazione}
\label{subsec:falsa_relazione}

Un errore melodico di natura \textit{orizzontale}, ma che si manifesta tra due voci \textit{diverse}, è la falsa relazione.

\begin{definizione}[Falsa relazione]
	\label{def:falsa_relazione}
	Si definisce \textbf{falsa relazione} una successione armonica in cui la stessa nota appare alterata in una voce e non alterata in un'altra, in accordi consecutivi. Essa distrugge la coerenza cromatica del tessuto.
	\begin{itemize}
		\item \textbf{Falsa relazione melodica:} si verifica quando un suono viene alterato in una voce diversa da quella che l'ha cantato non alterato nell'accordo precedente.
		\item \textbf{Falsa relazione di tritono:} si ha quando due voci consecutive cantano le due note del tritono diatonico (IV e VII grado della scala) in modo tale da creare una relazione armonica ambigua tra un accordo e il successivo.
	\end{itemize}
\end{definizione}

\begin{regola}[Come evitare la falsa relazione]
	\label{reg:evitare_falsa_relazione}
	Se una nota viene alterata cromaticamente, la stessa voce che l'ha cantata non alterata deve eseguire anche la versione alterata nell'accordo successivo. Non si deve ``passare'' l'alterazione a un'altra voce.
\end{regola}

% ============================================================
% SEZIONE: RISOLUZIONE DELLA SENSIBILE
% (Integrazione da armoniainrete.altervista.org/armonia-lezione-09)
% ============================================================
\section{La risoluzione della sensibile}
\label{sec:risoluzione_sensibile}

La \textbf{sensibile} --- il VII grado della scala --- è il grado più ``direzionato'' dell'intero sistema tonale. Essa dista solo un semitono dalla tonica e, in virtù di questa vicinanza, possiede una forza attrattiva che non può essere ignorata nella condotta delle voci.

\begin{regola}[I tre movimenti possibili della sensibile]
	\label{reg:tre_movimenti_sensibile}
	La sensibile può muoversi in tre modi:
	\begin{enumerate}
		\item \textbf{Sale di semitono alla tonica} (risoluzione primaria e obbligatoria nel soprano). Questo è il movimento più naturale e il più atteso dall'orecchio.
		\item \textbf{Scende di terza alla dominante}, a condizione che si trovi in una \textbf{parte interna} dell'accordo (contralto o tenore, mai soprano) e che il basso proceda in \textbf{moto contrario} rispetto a essa. Questo espediente permette di ottenere un accordo di tonica completo, evitando di triplicare la fondamentale.
		\item \textbf{Scende di semitono cromatico} (abbassamento della sensibile): questo movimento è possibile solo in presenza di un \textbf{cambio di tonalità o di una tonicizzazione}. La sensibile abbassata non è più sensibile della tonalità d'impianto, ma diventa un altro grado funzionale nella nuova tonalità locale.
	\end{enumerate}
	Ogni altro movimento della sensibile è da evitare, a meno che essa si trovi in una voce intermedia e il movimento sia giustificato dal contesto armonico.
\end{regola}

\begin{esempio}[Le tre risoluzioni della sensibile in Do Maggiore]
	\label{ex:risoluzioni_sensibile}

	\textbf{A. Risoluzione primaria (sensibile sale alla tonica):}
	\musica{es_cap04_01}

	\vspace{0.5cm}
	\textbf{B. Sensibile che scende di terza (voce interna, moto contrario con il basso):}
	\musica{es_cap04_02}

	In questo caso il Si (sensibile) si trova al contralto e scende di terza al Sol, mentre il basso sale da Sol a Do. Il moto contrario e la posizione interna della sensibile rendono il collegamento corretto.
\end{esempio}

\begin{errore}[Sensibile non risolta nel soprano]
	\label{err:sensibile_non_risolta}
	Se la sensibile si trova al \textbf{soprano}, essa \textbf{deve} salire di semitono alla tonica. Non è mai ammesso che il soprano scenda di terza, anche se tecnicamente la nota di arrivo appartiene all'accordo. Questo errore produce un soprano non cantabile e armonicamente incoerente.
\end{errore}

% ============================================================
% SEZIONE: ERRORI ARMONICI
% ============================================================
\section{Errori armonici}
\label{sec:errori_armonici}

Gli errori armonici si producono nel \textbf{collegamento tra due voci} e riguardano la qualità degli intervalli che le voci formano tra loro in successione. Si dividono in errori gravi ed errori di media gravità.

\subsection{Errori armonici gravi: quinte e ottave parallele}

\begin{errore}[Quinte e Ottave parallele]
	\label{err:quinte_ottave_parallele}
	È severamente vietato muovere due voci per moto parallelo mantenendo tra loro un intervallo di \textbf{quinta giusta} o di \textbf{ottava giusta}. Questo errore distrugge l'illusione di indipendenza delle parti, facendole percepire come un'unica linea melodica raddoppiata.
	\musica{es_cap04_03}
\end{errore}

Due ottave o due quinte sono ammesse se sono \textbf{separate da un'armonia diversa} e non sono quindi in successione diretta.

\begin{regola}[Quinte parallele Giusta – Diminuita]
	\label{reg:quinte_g_dim}
	È ammesso il movimento di quinte parallele quando la \textbf{prima quinta è giusta} e la \textbf{seconda è diminuita}. Questo movimento è accettato perché la quinta diminuita, essendo dissonante, non viene percepita come una vera quinta parallela.
	\musica{es_cap04_04}
\end{regola}

\subsection{Errori armonici meno gravi: quinte e ottave nascoste}

\begin{errore}[Quinte e ottave nascoste (o dirette)]
	\label{err:quinte_nascoste}
	Si deve evitare di raggiungere un'ottava o una quinta giusta tra le parti estreme (Soprano e Basso) muovendo entrambe le voci nella stessa direzione per salto.
\end{errore}

\begin{regola}[Eccezioni alle quinte e ottave nascoste]
	\label{reg:eccezioni_nascoste}
	Le quinte e ottave nascoste tra parti estreme sono tollerate nei seguenti casi:
	\begin{itemize}
		\item Per le \textbf{ottave nascoste}: il basso sale di almeno una quarta giusta e il soprano sale di semitono diatonico (tipicamente la sensibile che sale alla tonica).
		\item Per le \textbf{quinte nascoste}: una delle voci si muove per grado congiunto, oppure il basso si muove verso la tonica o la dominante e la voce superiore procede per grado.
	\end{itemize}
\end{regola}

\subsection{Incrocio delle parti}

\begin{errore}[Incrocio delle voci]
	\label{err:incrocio}
	È proibito che una voce più grave canti una nota più acuta della voce adiacente superiore. Il Tenore non può cantare una nota più acuta del Contralto, né il Contralto una più acuta del Soprano.
\end{errore}

% ============================================================
% SEZIONE: SINCOPE ARMONICA
% (Integrazione da armoniainrete lezione-08)
% ============================================================
\section{La sincope armonica}
\label{sec:sincope_armonica}

Un caso particolare di condotta delle voci è la \textbf{sincope armonica}: una voce tiene una nota che nell'accordo successivo risulta dissonante rispetto alle altre voci. Questo fenomeno è strettamente legato all'uso dei ritardi (Capitolo \ref{chap:ritardi}) ed è uno dei pilastri espressivi dello stile corale.

\begin{definizione}[Sincope armonica]
	\label{def:sincope_armonica}
	La \textbf{sincope armonica} si verifica quando una nota \textit{consonante} rispetto all'accordo corrente viene \textit{tenuta} nel momento in cui il basso si muove verso un nuovo accordo, diventando così \textit{dissonante} rispetto alla nuova armonia.
\end{definizione}

\begin{regola}[Trattamento della sincope armonica]
	\label{reg:sincope_armonica}
	La nota tenuta per sincope deve:
	\begin{enumerate}
		\item Essere preparata come consonanza nell'accordo precedente.
		\item Cadere sul tempo forte (o sulla parte forte del tempo), rafforzando la tensione.
		\item Risolvere per grado congiunto \textbf{discendente} verso la nota reale dell'accordo (in questo modo si comporta come un ritardo).
	\end{enumerate}
\end{regola}

% ============================================================
% SEZIONE: RADDOPPI
% ============================================================
\section{I raddoppi}
\label{sec:raddoppi}

Poiché scriviamo a quattro parti ma gli accordi hanno tipicamente tre note, una di esse deve essere raddoppiata (o triplicata nella risoluzione della V7 completa).

\begin{regola}[Principio generale dei raddoppi]
	\label{reg:raddoppi}
	\begin{itemize}
		\item \textbf{Raddoppiare di norma la fondamentale}, specialmente negli stati fondamentali.
		\item \textbf{Mai} raddoppiare la \textbf{sensibile} (VII grado), perché porta inevitabilmente a ottave parallele nella risoluzione (entrambe le sensibili devono salire alla tonica).
		\item \textbf{Mai} raddoppiare la \textbf{settima} degli accordi di settima.
		\item Nei rivolti, è spesso utile raddoppiare la \textbf{terza} (specialmente nel collegamento II7-V per evitare quinte parallele, cfr. \ref{reg:ii7_raddoppio}).
		\item Il III grado maggiore è un grado caratteristico e va raddoppiato con parsimonia.
	\end{itemize}
\end{regola}

% ============================================================
% SEZIONE: CHECKLIST REVISIONE
% ============================================================
\section{Lista di controllo: revisione dell'elaborato}

\begin{regola}[Checklist errori]
\label{reg:checklist_errori}
Una volta completata l'armonizzazione, sottoporre l'esercizio a questa revisione sistematica.

\textbf{Blocco 1 — Errori verticali (accordo per accordo):}
\begin{enumerate}
	\item Ogni accordo contiene le note corrette per la sua funzione?
	\item Il raddoppio è corretto? (Raddoppiare la fondamentale di default; \textbf{mai} raddoppiare sensibile o settima).
	\item Le voci rispettano la loro estensione? (\ref{def:estensione_voci}).
	\item Le distanze tra voci adiacenti sono corrette? (\ref{reg:distanza_voci}).
\end{enumerate}

\textbf{Blocco 2 — Errori orizzontali (voce per voce):}
\begin{enumerate}
	\item Ci sono quinte parallele tra qualsiasi coppia di voci?
	\item Ci sono ottave parallele tra qualsiasi coppia di voci?
	\item Ci sono quinte o ottave nascoste tra le voci esterne?
	\item La sensibile nel soprano risolve salendo alla tonica? (\ref{reg:tre_movimenti_sensibile}).
	\item La sensibile in voce interna, se scende di terza, lo fa in moto contrario con il basso?
	\item La settima risolve scendendo di grado?
	\item Ci sono salti non giustificati (sesta, settima) in una voce?
	\item Ci sono movimenti di seconda aumentata (specialmente nel minore)?
	\item Ci sono false relazioni tra voci diverse? (\ref{def:falsa_relazione}).
	\item Ci sono tritoni armonici tra voci non adiacenti in successione?
\end{enumerate}

\textbf{Blocco 3 — Errori strutturali (visione d'insieme):}
\begin{enumerate}
	\item Le cadenze sono posizionate nei punti metricamente corretti?
	\item Le frasi hanno una struttura riconoscibile? (\ref{chap:struttura_frase}).
	\item Il basso dato e l'eventuale numerica sono stati rispettati?
\end{enumerate}
\end{regola}

Si suggerisce di rileggere ogni \textbf{coppia} di voci isolata (S+B, A+T, S+A, T+B ecc.) anziché leggere tutto l'elaborato insieme: gli errori di parallelismo si individuano molto più facilmente con questa tecnica.
  % La scrittura a quattro parti
	\chapter{Le funzioni armoniche}
\label{chap:funzioni_armoniche}

Ogni accordo all'interno di una tonalità non è semplicemente un insieme di note sovrapposte: possiede una \textbf{funzione armonica}, ovvero un ruolo specifico nel discorso musicale. Comprendere le funzioni armoniche significa comprendere la logica interna del sistema tonale, le forze di attrazione e repulsione che governano il movimento da un accordo all'altro.

\section{Le tre funzioni fondamentali}
\label{sec:tre_funzioni}

Nel sistema tonale, tutti gli accordi possono essere ricondotti a tre funzioni principali: \textbf{tonica}, \textbf{dominante} e \textbf{sottodominante}. Queste funzioni rappresentano rispettivamente riposo, tensione e movimento.

\begin{esempio}[Funzioni fondamentali]
\musica{es_cap05_01}
\end{esempio}

\begin{definizione}[Funzioni armoniche]
	
	\begin{itemize}
		\item \label{def:funzione_tonica}
		La \textbf{funzione di tonica} rappresenta il punto di riposo, la stabilità, la "casa" della tonalità. L'accordo principale che svolge questa funzione è la \textbf{triade sul I grado}. La tonica è il punto di partenza e di arrivo della maggior parte dei discorsi armonici.
		\item \label{def:funzione_dominante}
		La \textbf{funzione di dominante} rappresenta la massima tensione armonica, l'impulso irresistibile verso la risoluzione sulla tonica. L'accordo principale che svolge questa funzione è la \textbf{triade sul V grado}. La dominante contiene la \textit{sensibile} (VII grado della scala), che tende con forza verso la tonica.
		\item 	\label{def:funzione_sottodominante}
		La \textbf{funzione di sottodominante} rappresenta il movimento, l'allontanamento dalla tonica, la preparazione della dominante. L'accordo principale che svolge questa funzione è la \textbf{triade sul IV grado}. La sottodominante contiene il IV grado della scala, che tende a scendere verso il III grado.
	\end{itemize}
\end{definizione}

	La progressione, ovvero successione di accordi, I - IV - V - I rappresenta il modello base del movimento armonico tonale: partenza dalla tonica, allontanamento attraverso la sottodominante, tensione sulla dominante, ritorno alla tonica.

\section{Gli accordi secondari e le loro funzioni}
\label{sec:accordi_secondari}

Gli accordi costruiti sui gradi II, III, VI e VII, pur avendo una loro identità, possono essere pensati come "sostituti" delle tre funzioni principali. La chiave per capire questa relazione sta nelle \textbf{note comuni} che questi accordi condividono con gli accordi principali.

\begin{esempio}[II come sottodominante]
\musica{es_cap05_02}
\end{esempio}

Due accordi che condividono molte note tendono a svolgere funzioni simili. Osserviamo le triadi diatoniche di Do maggiore, che useremo per comprendere le prossime definizioni:

\begin{center}
\begin{tabular}{|c|c|c|c|c|}
\hline
\textbf{Grado} & \textbf{Note} & \textbf{Note comuni con I} & \textbf{Note comuni con IV} & \textbf{Note comuni con V} \\
\hline
I & Do Mi Sol & --- & Do & Sol \\
II & Re Fa La & --- & Fa La & Re \\
III & Mi Sol Si & Mi Sol & --- & Sol Si \\
IV & Fa La Do & Do & --- & --- \\
V & Sol Si Re & Sol & --- & --- \\
VI & La Do Mi & Do Mi & La & --- \\
VII & Si Re Fa & --- & Fa & Si Re \\
\hline
\end{tabular}
\end{center}

\subsection{Accordi in ogni grado della scala}
\label{subsec:accordi_funzioni}

\begin{definizione}[Accordi con funzione di tonica]
	\label{def:accordi_tonica}
	Oltre al I grado, possono svolgere funzione di tonica:
	\begin{itemize}
		\item \textbf{III grado:} condivide due note con la tonica. È la tonica "debole" o "relativa".
		\item \textbf{VI grado:} condivide due note con la tonica. È frequentemente usato nella \textit{cadenza d'inganno} come sostituto della tonica.
	\end{itemize}
\end{definizione}


\begin{definizione}[Accordi con funzione di sottodominante]
	\label{def:accordi_sottodominante}
	Oltre al IV grado, possono svolgere funzione di sottodominante:
	\begin{itemize}
		\item \textbf{II grado:} condivide due note con la sottodominante. È un sostituto molto comune del IV grado, specialmente nella progressione II - V - I.
		\item \textbf{VI grado:} condivide una nota con la sottodominante. In certi contesti può avere sfumatura sottodominante.
	\end{itemize}
\end{definizione}


\begin{definizione}[Accordi con funzione di dominante]
	\label{def:accordi_dominante}
	Oltre al V grado, può svolgere funzione di dominante:
	\begin{itemize}
		\item \textbf{VII grado:} condivide due note con la dominante e, soprattutto, contiene la sensibile. È essenzialmente una dominante senza fondamentale.
	\end{itemize}
\end{definizione}

\begin{matematica}[L'identità del III grado: tra Tonica e Dominante]
	Possiamo analizzare la parentela tra gli accordi utilizzando la teoria degli insiemi. Considerando le triadi in Do maggiore come insiemi di note, osserviamo:
	\begin{itemize}
		\item $I = \{Do, Mi, Sol\}$
		\item $III = \{Mi, Sol, Si\}$
		\item $V = \{Sol, Si, Re\}$
	\end{itemize}
	Si nota che la cardinalità dell'intersezione è identica in entrambi i casi: 
	$|I \cap III| = 2$ ($Mi, Sol$) e $|V \cap III| = 2$ ($Sol, Si$). 
	
	Perché allora il III grado è considerato una "tonica debole"  invece di una dominante? La risposta risiede nella \textbf{gerarchia pesata} degli elementi:
	\begin{enumerate}
		\item Il $Mi$ è la \textit{caratteristica modale} della tonica; la sua presenza nel III grado stabilisce un legame identitario fortissimo con la "casa" della tonalità.
		\item Sebbene il III grado contenga la sensibile ($Si$), manca della \textit{fondamentale} della dominante ($Sol$ al basso o come nota portante), che è necessaria per generare la spinta propulsiva verso la risoluzione.
	\end{enumerate}
	Pertanto, pur avendo lo stesso numero di note in comune con il V grado, la "forza gravitazionale" esercitata dalle note in comune con il I grado prevale, rendendolo un sostituto del riposo piuttosto che della tensione.
\end{matematica}


	Il seguente schema riassume le relazioni funzionali nel modo maggiore:
	
	\begin{center}
	\begin{tikzpicture}[
		box/.style={draw, minimum width=2cm, minimum height=1cm, align=center},
		arrow/.style={->, >=stealth, thick}
	]
	% Tonica
	\node[box, fill=white] (T) at (0,0) {\textbf{TONICA}\\I, III, VI};
	% Sottodominante  
	\node[box, fill=white] (S) at (-3,-2) {\textbf{SOTTOD.}\\IV, II};
	% Dominante
	\node[box, fill=white] (D) at (3,-2) {\textbf{DOMIN.}\\V, VII};
	
	% Frecce
	\draw[arrow] (T) -- (S);
	\draw[arrow] (S) -- (D);
	\draw[arrow] (D) -- (T);
	\draw[arrow, dashed] (T) -- (D);
	\end{tikzpicture}
	\end{center}
	
	Il percorso tipico è T → S → D → T, ma è possibile anche il collegamento diretto T → D → T.


\section{L'armonizzazione della scala: la regola dell'ottava}
\label{sec:regola_ottava_funzioni}

La \textbf{Regola dell'Ottava} è uno schema pratico e sistematico che insegna come armonizzare ogni grado della scala quando esso appare al basso. Questa regola, che ha origini storiche nella pratica del basso continuo del periodo barocco, rappresenta un potente strumento didattico per comprendere la logica dell'armonizzazione.

\begin{definizione}[Regola dell'ottava]
	\label{def:versioni_regola_ottava}
	Nel corso di questo manuale presenteremo la Regola dell'Ottava in \textbf{tre versioni progressive}, ciascuna più "completa" e "ricca" della precedente:
	
	\begin{enumerate}
		\item \textbf{Versione preliminare (questo capitolo):} utilizza esclusivamente \textbf{triadi} in stato fondamentale e rivolti.
		\item \textbf{Versione intermedia (Capitolo \ref{chap:settima_dominante}):} integra le \textbf{settime di dominante} e i loro rivolti nello schema.
		\item \textbf{Versione completa (Capitolo successivo):} include l'intero arsenale delle settime per una realizzazione più ricca e flessibile.
	\end{enumerate}
	
	Questo approccio progressivo permette di padroneggiare gradualmente la complessità della regola, evitando di essere sopraffatti dalle numerose possibilità armoniche fin dall'inizio.
\end{definizione}

\subsection{Movimenti del basso e assegnazione delle funzioni}
\label{subsec:principio_regola_ottava}

Quando il basso si muove per grado congiunto, la scelta dell'accordo dipende dalla posizione della nota all'interno della scala. Il principio guida di questa versione preliminare è quello di utilizzare i rivolti per mantenere il movimento melodico senza appesantire la struttura con troppi stati fondamentali.

\begin{regola}{Regola dell'Ottava (triadi) - Scala Maggiore Ascendente}
	\label{reg:regola_ottava_maggiore_asc}
	Armonizzazione dei gradi della scala maggiore in movimento ascendente:
	
	\begin{itemize}
		\item \textbf{I grado:} Accordo di \textbf{I} in stato fondamentale (5)
		\item \textbf{II grado:} Accordo di \textbf{VII} in primo rivolto (6)
		\item \textbf{III grado:} Accordo di \textbf{I} in primo rivolto (6)
		\item \textbf{IV grado:} Accordo di \textbf{II} in primo rivolto (6)
		\item \textbf{V grado:} Accordo di \textbf{V} in stato fondamentale (5) — oppure I$^6_4$ di passaggio
		\item \textbf{VI grado:} Accordo di \textbf{IV} in primo rivolto (6)
		\item \textbf{VII grado:} Accordo di \textbf{V} in primo rivolto (6)
		\item \textbf{VIII grado:} Accordo di \textbf{I} in stato fondamentale (5)
	\end{itemize}
\end{regola}



\begin{regola}{Regola dell'Ottava (triadi) - Scala Maggiore Discendente}
	\label{reg:regola_ottava_maggiore_disc}
	Armonizzazione dei gradi della scala maggiore in movimento discendente:
	
	\begin{itemize}
		\item \textbf{VIII grado:} Accordo di \textbf{I} in stato fondamentale (5)
		\item \textbf{VII grado:} Accordo di \textbf{V} in primo rivolto (6)
		\item \textbf{VI grado:} Accordo di \textbf{IV} in primo rivolto (6)
		\item \textbf{V grado:} Accordo di \textbf{V} in stato fondamentale (5)
		\item \textbf{IV grado:} Accordo di \textbf{II} in primo rivolto (6)
		\item \textbf{III grado:} Accordo di \textbf{I} in primo rivolto (6)
		\item \textbf{II grado:} Accordo di \textbf{VII} in primo rivolto (6)
		\item \textbf{I grado:} Accordo di \textbf{I} in stato fondamentale (5)
	\end{itemize}
\end{regola}



\begin{regola}{Regola dell'Ottava (triadi) - Scala Minore Armonica Ascendente}
	\label{reg:regola_ottava_minore_asc}
	Armonizzazione dei gradi della scala minore armonica in movimento ascendente:
	
	\begin{itemize}
		\item \textbf{I grado:} Accordo di \textbf{I} in stato fondamentale (5)
		\item \textbf{II grado:} Accordo di \textbf{VII°} in primo rivolto (6)
		\item \textbf{III grado:} Accordo di \textbf{I} in primo rivolto (6)
		\item \textbf{IV grado:} Accordo di \textbf{II°} in primo rivolto (6)
		\item \textbf{V grado:} Accordo di \textbf{V} in stato fondamentale (5) — oppure I$^6_4$ di passaggio
		\item \textbf{VI grado:} Accordo di \textbf{IV} in primo rivolto (6)
		\item \textbf{VII grado:} Accordo di \textbf{V} in primo rivolto (6)
		\item \textbf{VIII grado:} Accordo di \textbf{I} in stato fondamentale (5)
	\end{itemize}
\end{regola}



\begin{regola}{Regola dell'Ottava (triadi) - Scala Minore Armonica Discendente}
	\label{reg:regola_ottava_minore_disc}
	Armonizzazione dei gradi della scala minore armonica in movimento discendente:
	
	\begin{itemize}
		\item \textbf{VIII grado:} Accordo di \textbf{I} in stato fondamentale (5)
		\item \textbf{VII grado:} Accordo di \textbf{V} in primo rivolto (6) — può essere \textbf{III}
		\item \textbf{VI grado:} Accordo di \textbf{IV} in primo rivolto (6)
		\item \textbf{V grado:} Accordo di \textbf{V} in stato fondamentale (5)
		\item \textbf{IV grado:} Accordo di \textbf{II°} in primo rivolto (6)
		\item \textbf{III grado:} Accordo di \textbf{I} in primo rivolto (6)
		\item \textbf{II grado:} Accordo di \textbf{VII°} in primo rivolto (6)
		\item \textbf{I grado:} Accordo di \textbf{I} in stato fondamentale (5)
	\end{itemize}
\end{regola}



\begin{regola}{Regola dell'Ottava (triadi) - Scala Minore Melodica Ascendente}
	\label{reg:regola_ottava_melodica_asc}
	Armonizzazione dei gradi della scala minore melodica in movimento ascendente (con VI e VII gradi alterati):
	
	\begin{itemize}
		\item \textbf{I grado:} Accordo di \textbf{I} in stato fondamentale (5)
		\item \textbf{II grado:} Accordo di \textbf{VII°} in primo rivolto (6)
		\item \textbf{III grado:} Accordo di \textbf{I} in primo rivolto (6)
		\item \textbf{IV grado:} Accordo di \textbf{II} in primo rivolto (6)
		\item \textbf{V grado:} Accordo di \textbf{V} in stato fondamentale (5) — oppure I$^6_4$ di passaggio
		\item \textbf{VI grado:} Accordo di \textbf{IV} in primo rivolto (6)
		\item \textbf{VII grado:} Accordo di \textbf{V} in primo rivolto (6)
		\item \textbf{VIII grado:} Accordo di \textbf{I} in stato fondamentale (5)
	\end{itemize}
	
	\textit{Nota: Nella scala minore melodica ascendente, il VI e VII grado sono alterati ascendentemente, rendendo il IV grado maggiore (II) e il VII grado maggiore (VII), come nella scala maggiore.}
\end{regola}


\begin{regola}{Regola dell'Ottava (triadi) - Scala Minore Melodica Discendente}
	\label{reg:regola_ottava_melodica_disc}
	Armonizzazione dei gradi della scala minore melodica in movimento discendente (forma naturale):
	
	\begin{itemize}
		\item \textbf{VIII grado:} Accordo di \textbf{I} in stato fondamentale (5)
		\item \textbf{VII grado:} Accordo di \textbf{III} in primo rivolto (6) — oppure \textbf{V}$^6$
		\item \textbf{VI grado:} Accordo di \textbf{IV} in primo rivolto (6)
		\item \textbf{V grado:} Accordo di \textbf{V} in stato fondamentale (5)
		\item \textbf{IV grado:} Accordo di \textbf{II°} in primo rivolto (6)
		\item \textbf{III grado:} Accordo di \textbf{I} in primo rivolto (6)
		\item \textbf{II grado:} Accordo di \textbf{VII°} in primo rivolto (6)
		\item \textbf{I grado:} Accordo di \textbf{I} in stato fondamentale (5)
	\end{itemize}
	
	\textit{Nota: Nella scala minore melodica discendente si utilizzano i gradi naturali (VI e VII abbassati), coincidendo con la scala minore naturale.}
\end{regola}


\section{Basso che salta o che procede per grado congiunto}
\label{sec:basso_salti_gradi}

La distinzione tra movimento melodico (grado congiunto) e movimento armonico (salto) è fondamentale per decidere quale stato dell'accordo utilizzare.

\begin{regola}[Scelta dello stato dell'accordo in base al movimento del basso]
	\label{reg:stato_basso}
	\begin{itemize}
		\item \textbf{Basso per grado congiunto}: si applica la \textbf{Regola dell'Ottava}.
		\item \textbf{Basso che salta}: si predilige lo \textbf{stato fondamentale}, poiché il salto indica un cambio di pilastro armonico.
		\item \textbf{Eccezione del V grado}: l'accordo di dominante può essere sempre usato in stato fondamentale per sottolineare la forza delle cadenze, indipendentemente dal movimento del basso.
	\end{itemize}
\end{regola}

\begin{regola}[Salto d'ottava al basso]
	\label{reg:salto_ottava_basso}
	Quando il basso compie un salto d'ottava (sia verso l'alto che verso il basso), la seconda nota del salto si numera generalmente con il \textbf{6} (primo rivolto della triade), creando un arpeggio armonico più fluido e variato.
\end{regola}



\section{Il secondo rivolto e la funzione del V grado}
\label{sec:secondo_rivolto_regola}

In questa versione preliminare della regola, l'accordo sul \textbf{V grado} della scala ammette una duplice natura. Sebbene la triade di dominante in stato fondamentale sia la scelta più comune per garantire stabilità, è possibile utilizzare il secondo rivolto della tonica (\textbf{I6/4}).

L'uso del I6/4 sul quinto grado della scala può avvenire in due modalità:

\begin{itemize}
	\item \textbf{I6/4 di passaggio}: quando il basso tocca il V grado muovendosi tra il IV e il VI (o viceversa), creando un collegamento morbido.
	\item \textbf{I6/4 cadenzale}: quando il V grado al basso viene "appoggiato" da quarta e sesta che poi risolvono sulla terza e quinta dell'accordo di dominante (V). Questo uso sarà approfondito nei capitoli successivi.
\end{itemize}

Queste opzioni rendono il V grado della scala uno dei punti più flessibili e interessanti dell'intera Regola dell'Ottava, permettendo al compositore di scegliere tra la forza della dominante pura e l'eleganza della sospensione cadenzale.

\begin{fisica}[Perché il I6/4 sul V grado suona "sospeso"?]
	Quando suoniamo il secondo rivolto della tonica (I6/4) con il Sol al basso (V grado in Do maggiore), le note sono Sol-Do-Mi. Dal punto di vista acustico, questo accordo presenta una caratteristica particolare:
	
	\begin{itemize}
		\item La quarta Sol-Do crea un intervallo il cui rapporto di frequenza è 4:3
		\item Questo rapporto non appare tra i primi armonici semplici della nota al basso (Sol)
		\item Gli armonici di Sol includono: Sol, Re, Sol, Si, Re... ma non Do
		\item Il Do risulta quindi "alieno" rispetto alla serie armonica naturale del basso
	\end{itemize}
	
	Questa "estraneità acustica" del Do rispetto al basso Sol genera la sensazione di instabilità e sospensione che caratterizza il I6/4. Quando questo accordo risolve sul V (Sol-Si-Re-Fa), la quarta Sol-Do scende alla terza Sol-Si, che è invece presente come terzo armonico del basso, creando una risoluzione acusticamente "naturale" e soddisfacente.
	
	Questa è la ragione fisica per cui il secondo rivolto, pur essendo tecnicamente un accordo di tonica, funziona armonicamente come preparazione della dominante: la serie armonica stessa ci guida verso la risoluzione.
\end{fisica}

% ============================================================
% CAPITOLO 4
% ============================================================
  % Funzioni armoniche
		\chapter{Le cadenze e la sintassi musicale}
	\label{chap:cadenze}
	
	Se gli accordi sono le parole del linguaggio musicale, le \textbf{cadenze} ne sono la punteggiatura. 
		
\begin{definizione}[La Cadenza]
	\label{def:cadenza}
	Una \textbf{cadenza} è una formula melodico-armonica costituita da una successione di due o più accordi che crea un punto di articolazione, pausa o conclusione nel discorso musicale.
\end{definizione}

Le cadenze svolgono una funzione analoga alla punteggiatura nel linguaggio verbale: delimitano le frasi, creano momenti di sospensione o di chiusura, e strutturano il flusso temporale della composizione.

Le cadenze possono essere classificate in base al loro effetto percettivo e alla loro funzione formale. Alcune cadenze producono un senso di \textbf{conclusione} e \textbf{riposo}, altre creano \textbf{sospensione} e \textbf{attesa}. Questa capacità di guidare le aspettative dell'ascoltatore è uno degli strumenti più potenti a disposizione del compositore per dare forma e significato alla musica.

Storicamente, le cadenze sono state codificate nella pratica compositiva come formule ricorrenti, riconoscibili e immediatamente efficaci. La loro forza non risiede solo nella successione di accordi in sé, ma anche nella condotta delle voci, nel ritmo armonico e nella posizione metrica in cui vengono collocate.
	
	\section{Cadenze conclusive}
	\label{sec:cadenze_conclusive}
	
	Le cadenze conclusive, come suggerisce il nome, producono un effetto di chiusura e riposo. Sono utilizzate per terminare le frasi musicali in modo definitivo.
	
	\begin{definizione}[Cadenza Perfetta]
		\label{def:cadenza_perfetta}
		La \textbf{cadenza perfetta} (o autentica perfetta) è la formula conclusiva più forte e inequivocabile del sistema tonale. Consiste nella successione degli accordi \textbf{V - I}, entrambi in stato fondamentale. Per essere considerata "perfetta", deve soddisfare una condizione aggiuntiva: la nota finale al soprano deve essere la tonica.
	\end{definizione}
	
	La cadenza perfetta è il culmine della tensione tonale: la dominante, con la sua sensibile che spinge verso l'alto, risolve nel centro di gravità della tonica, portando il discorso a una quiete completa.
	
	\begin{esempio}[Cadenza Perfetta in Do Maggiore]
		\label{ex:cadenza_perfecta}
		\musica{
			\partilate
			{ \key do \major \time 4/4 }
			{ re'2 do'2 }
			{ si2 sol2 }
			{ sol2 mi2 }
			{ sol,2 do2 }
			\lyricmode { "V"2 "I"2 }
			\figuremode {}
		}
	\end{esempio}

	\begin{definizione}[Cadenza Imperfetta]
		\label{def:cadenza_imperfetta}
		La \textbf{cadenza imperfetta} (o autentica imperfetta) è una cadenza V-I (o V7-I) in cui almeno una delle seguenti condizioni non è soddisfatta:
		\begin{itemize}
			\item L'accordo di V non è in stato fondamentale (es. V6 o V6/5 invece di V).
			\item L'accordo di I non è in stato fondamentale (es. I6 invece di I).
			\item Il soprano all'accordo di arrivo non ha la tonica (es. ha la terza o la quinta).
		\end{itemize}
		Produce un senso di chiusura meno definitivo della cadenza perfetta, ed è spesso usata a metà periodo (\ref{chap:struttura_frase}).
	\end{definizione}

	\begin{esempio}[Esempi di Cadenza Imperfetta]
		Mostriamo due varianti: soprano sulla terza e V6 invece di V.
		\musica{
			\partilate
			{ \key do \major \time 4/4 }
			{ re'2 mi' | r2 re' | si2 do' }
			{ si2 do' | r2 si | sol2 sol }
			{ sol2 sol | r2 fa | re2 mi }
			{ sol,2 do | r2 sol, | sol,2 do }
			\lyricmode { "V-I (sop.3a)" | "V6-I" }
			\figuremode { s2 s | r <6> | s s }
		}
	\end{esempio}

    \begin{regola}[Gerarchia delle cadenze]
        \label{reg:gerarchia_cadenze}
        La forza conclusiva delle cadenze segue questa gerarchia decrescente:
        \begin{center}
            \textbf{C. Perfetta > C. Imperfetta > C. Plagale > C. Sospesa > C. d'Inganno}
        \end{center}
        Questa gerarchia guida il compositore nella scelta della cadenza in base alla collocazione nella struttura formale (\ref{chap:struttura_frase}).
    \end{regola}
	
	\begin{definizione}[Cadenza Plagale]
		\label{def:cadenza_plagale}
		La \textbf{cadenza plagale} consiste nella successione \textbf{IV - I}. Produce un senso di chiusura più debole e statico rispetto alla cadenza perfetta. È spesso usata come formula conclusiva aggiuntiva dopo una cadenza perfetta, ed è universalmente nota come la cadenza dell'"Amen" nella musica liturgica.
	\end{definizione}
	
	\begin{esempio}[Cadenza Plagale in Do Maggiore]
		\label{ex:cadenza_plagale}
		\musica{
			\partilate
			{ \key do \major \time 4/4 }
			{ fa'2 mi'2 }
			{ do'2 do'2 }
			{ la2 sol2 }
			{ fa,2 do2 }
			\lyricmode { "IV"2 "I"2 }
			\figuremode {}
		}
	\end{esempio}
	
	\section{Cadenze sospensive}
	\label{sec:cadenze_sospensive}
	
	Le cadenze sospensive creano un effetto di interruzione e attesa, lasciando il discorso musicale "aperto". Sono l'equivalente musicale di una virgola o di un punto interrogativo.
	
	\begin{definizione}[Cadenza Sospesa]
		\label{def:cadenza_sospesa}
		Una \textbf{cadenza sospesa} (o mezza cadenza) è una qualsiasi formula armonica che termina sull'accordo di \textbf{Dominante (V)}. L'accordo che la precede può essere la Tonica (I), la Sottodominante (IV), la Sopratonica (II) o altri. La sua funzione è quella di creare un'aspettativa per la risoluzione sulla tonica, che però viene posticipata.
	\end{definizione}
	
	\begin{esempio}[Cadenza Sospesa (IV - V)]
		\label{ex:cadenza_sospesa}
		\musica{
			\partilate
			{ \key do \major \time 4/4 }
			{ fa'2 sol'2 }
			{ do'2 re'2 }
			{ la2 si2 }
			{ fa,2 sol,2 }
			\lyricmode { "IV"2 "V"2 }
			\figuremode {}
		}
	\end{esempio}
	
	\begin{definizione}[Cadenza d'Inganno]
		\label{def:cadenza_inganno}
		La \textbf{cadenza d'inganno} è una delle figure più espressive dell'armonia. Consiste nella successione \textbf{V - VI}. Essa "inganna" l'ascoltatore, che si aspetterebbe la risoluzione sulla tonica (I) dopo la dominante. Al suo posto, appare l'accordo del VI grado, che condivide due note con la tonica ma ne sposta il centro tonale, creando una sorpresa e prolungando la frase musicale.
	\end{definizione}
	
	Nella cadenza d'inganno, la condotta delle voci segue regole precise: la sensibile (terza della dominante) sale comunque alla tonica, mentre il basso sale di grado congiunto. Per evitare quinte parallele, si raddoppia tipicamente la terza dell'accordo del VI grado.
	
	\begin{esempio}[Cadenza d'Inganno in Do Maggiore]
		\label{ex:cadenza_inganno}
		\musica{
			\partilate
			{ \key do \major \time 4/4 }
			{ si'2 do''2 } % La sensibile (Si) sale alla tonica (Do)
			{ re'2 mi'2 }
			{ sol2 la2 }
			{ sol,2 la,2 }   % Il basso sale di grado
			\lyricmode { "V"2 "VI"2 }
			\figuremode {}
		}
	\end{esempio}
	
	\section{La fraseologia}
	\label{sec:fraseologia}
	
	Le cadenze non sono eventi isolati, ma i punti di arrivo di un percorso logico governato dalle \textbf{funzioni armoniche}. La struttura più comune di una frase musicale può essere riassunta dal modello T-S-D-T.
	
	\begin{regola}[Logica Funzionale: T-S-D-T]
		\label{reg:tsdt}
		Il modello \textbf{T-S-D-T} descrive il ciclo di tensione e risoluzione che sta alla base della sintassi tonale:
		\begin{itemize}
			\item \textbf{T (Tonica):} Area di stabilità e riposo (I grado).
			\item \textbf{S (Sottodominante):} Area di allontanamento dalla tonica (IV, II grado).
			\item \textbf{D (Dominante):} Area di massima tensione, che prepara il ritorno alla tonica (V grado).
			\item \textbf{T (Tonica):} Ritorno al punto di riposo, chiusura del ciclo.
		\end{itemize}
		Questa progressione fondamentale (Stabilità $\rightarrow$ Allontanamento $\rightarrow$ Tensione $\rightarrow$ Risoluzione) costituisce l'ossatura di innumerevoli brani musicali.
	\end{regola}
	
	\begin{esempio}[Una frase completa basata sul modello T-S-D-T]
		\label{ex:tsdt}
		Una semplice realizzazione del modello I - IV - V - I.
		\musica{
			\partilate
			{ \key do \major \time 4/4 }
			{ mi'4 fa'4 sol'4 do'4 }
			{ do'4 do'4 re'4 mi'4 }
			{ sol4 la4 si4 sol4 }
			{ do4 fa,4 sol,4 do4 }
			\lyricmode { "I (T)"4 "IV (S)"4 "V (D)"4 "I (T)"4 }
			\figuremode {}
		}
	\end{esempio}
	
	
	% ============================================================
	% CAPITOLO 6
	% ============================================================
  % Le cadenze e la sintassi musicale
	\chapter{La struttura della frase e il periodo}
\label{chap:struttura_frase}

La musica non è una successione casuale di accordi, ma un discorso organizzato. Le cadenze fungono da punteggiatura, definendo unità strutturali chiamate frasi e periodi.

\section{La frase musicale}
\label{sec:frase_musicale}

\begin{definizione}[Frase musicale]
    È l'unità minima di senso compiuto nel discorso tonale. Tipicamente composta da 4 misure (ma con varianti), termina sempre con una cadenza che ne delimita il confine.
\end{definizione}

\begin{regola}
    In un basso dato, ogni frase deve concludersi con una cadenza riconoscibile. Proprio come una frase verbale ha un soggetto, uno sviluppo e un punto finale, la frase musicale ha un inizio, un percorso armonico e una conclusione cadenzale.
\end{regola}

\section{Il periodo}
\label{sec:periodo}

\begin{definizione}[Il periodo]
    È una struttura di norma di 8 misure, formata dalla successione di due frasi in relazione tra loro:
    \begin{itemize}
        \item \textbf{Antecedente} (mis. 1-4): termina con una cadenza non conclusiva (sospesa o imperfetta). Rimanda a \ref{def:cadenza_sospesa} e \ref{def:cadenza_imperfetta}.
        \item \textbf{Conseguente} (mis. 5-8): termina con una cadenza conclusiva (perfetta). Rimanda a \ref{def:cadenza_perfetta}.
    \end{itemize}
\end{definizione}

La relazione tra antecedente e conseguente è analoga a quella tra domanda e risposta: la prima frase crea un'aspettativa, la seconda la soddisfa.

\begin{esempio}[Struttura del periodo in un basso dato]
    Indichiamo con \textit{markup} i punti cadenzali di un basso di 8 misure.
    \musica{es_cap07_01}
\end{esempio}

\section{Applicazione al basso dato}
\label{sec:applicazione_frase}

\begin{regola}[Pianificazione cadenzale]
    Prima di armonizzare un basso, individuare:
    \begin{itemize}
        \item \textbf{Metà (mis. 4 di 8):} Deve presentare una cadenza sospesa (V) o imperfetta (I non fondamentale o soprano non in ottava).
        \item \textbf{Finale (mis. 8):} Deve presentare una cadenza perfetta (V-I in stato fondamentale).
    \end{itemize}
\end{regola}

\begin{regola}[Scelta della cadenza di metà]
    Preferire la cadenza sospesa su V se il basso si ferma sulla dominante; preferire la cadenza imperfetta se il basso approda al I ma senza la forza conclusiva necessaria per il finale.
\end{regola}

\begin{esempio}[Armonizzazione di un periodo di 8 misure]
    Mostriamo come la scelta delle cadenze definisce la struttura formale del brano.
    \musica{es_cap07_02}
\end{esempio}
 % La struttura della frase e il periodo
	\chapter{Teoria generale della dissonanza}
\label{chap:teoria_dissonanza}

Prima di addentrarci nello studio degli accordi di quattro suoni (le quadriadi), è indispensabile comprendere il comportamento della \textbf{dissonanza} nel sistema tonale. Se la consonanza rappresenta lo stato di quiete e stabilità, la dissonanza è l'elemento dinamico che genera tensione, movimento e la necessità di una risoluzione.

\section{Il ciclo della dissonanza}
\label{sec:ciclo_dissonanza}

Nello stile severo e nella teoria classica dell'armonia, la gestione di una nota dissonante non è arbitraria, ma segue un protocollo rigoroso diviso in tre fasi cronologiche.

\begin{regola}[Trattamento della dissonanza]
        \label{reg:trattamento_dissonanza}
        Ogni nota dissonante deve attraversare tre momenti distinti:
        \begin{enumerate}
                \item \textbf{Preparazione:} La nota che diventerà dissonante deve essere presente nell'accordo precedente come \textbf{consonanza} (solitamente mantenuta nella stessa voce). Questa fase "giustifica" l'ingresso della tensione.
                \item \textbf{Percussione:} La nota viene colpita (suonata) simultaneamente alle altre note dell'accordo, diventando effettivamente una \textbf{dissonanza} rispetto alla fondamentale o ad altre voci.
                \item \textbf{Risoluzione:} La nota dissonante deve muoversi per \textbf{grado congiunto discendente} verso una nota che sia consonante nell'accordo successivo.
        \end{enumerate}
\end{regola}

\begin{esempio}[Esempio di preparazione e risoluzione]
        \label{ex:ciclo_dissonanza}
        In un collegamento II-V, la settima del II grado deve essere preparata dalla tonica del I grado o da una nota del IV grado.
        \musica{
                \partilate
                { \key do \major \time 4/4 }
                { do''2 do''2 | si'1 }
                { sol'2 fa'2 | re'1 }
                { mi'2 la2 | sol1 }
                { do2 re2 | sol,1 }
                \lyricmode { "I"2 "II7"2 | "V"1 }
                \figuremode { s2 <7>2 | s1 }
        }
        In questo esempio, il Do al soprano è consonante nel primo accordo (ottava), diventa dissonante nel secondo (settima del Re) e risolve scendendo al Si (terza del Sol).
\end{esempio}

\section{Evoluzione del concetto di preparazione}
\label{sec:evoluzione_preparazione}

Mentre la \textbf{risoluzione discendente} rimane un pilastro quasi inviolabile dell'armonia tonale, la necessità della \textbf{preparazione} è andata scemando con l'evoluzione del linguaggio musicale.

Alcune dissonanze, per via della loro forte direzionalità e familiarità (come la settima di dominante), hanno iniziato a essere introdotte "di petto", ovvero senza preparazione. Tuttavia, anche quando la preparazione è omessa, la nota dissonante mantiene intatta la sua natura e la sua esigenza di risoluzione.

\begin{regola}[Dissonanze non preparate]
        \label{reg:dissonanze_non_preparate}
        Gli accordi di "dominante" (V7, VII7, e derivati) ammettono frequentemente la percussione della dissonanza senza preparazione. Al contrario, gli accordi di settima sulle altre specie (II7, IV7, ecc.) richiedono quasi sempre una preparazione rigorosa nello stile scolastico.
\end{regola}

Questo capitolo funge da fondamento teorico per tutti i successivi studi sulle quadriadi:
\begin{itemize}
        \item Per la \textbf{settima di dominante} (Capitolo \ref{chap:settima_dominante}), dove la risoluzione è il motore della cadenza.
        \item Per il \textbf{VII7} e il \textbf{II7} (Capitolo \ref{chap:vii7_ii7}), dove la preparazione diventa cruciale per la fluidità del discorso.
        \item Per tutte le altre \textbf{specie di settima} (Capitolo \ref{chap:sette_specie}), che analizzeremo come estensioni di questo principio.
\end{itemize}
 % Teoria generale della dissonanza
	\chapter{La settima di dominante}
\label{chap:settima_dominante}

Se la triade è la "molecola" base dell'armonia, la \textbf{settima di dominante} (V7) è il suo motore. Aggiungendo una quarta nota alla triade di dominante, si genera un accordo di straordinaria tensione e forza propulsiva. La sua corretta gestione è una delle competenze chiave nello studio dell'armonia tonale.

\section{Genesi e struttura}
\label{sec:genesi_v7}

L'accordo di settima di dominante si costruisce aggiungendo una terza minore sopra la quinta della triade di dominante. In Do Maggiore, la dominante è Sol (Sol-Si-Re). Aggiungendo una terza minore sopra il Re, otteniamo un Fa. L'accordo completo è Sol-Si-Re-Fa.

\begin{definizione}[Il Tritono]
	\label{def:tritono}
	Il \textbf{tritono} è un intervallo di tre toni interi, che può presentarsi come una quarta aumentata o una quinta diminuita. È l'intervallo più dissonante del sistema diatonico e, storicamente, veniva chiamato \textit{diabolus in musica}.
\end{definizione}

La caratteristica fondamentale della V7 è la presenza, al suo interno, di un tritono tra la sua terza (la sensibile) e la sua settima. In Do Maggiore, questo avviene tra il \textbf{Si} (sensibile) e il \textbf{Fa} (settima dell'accordo). Questa dissonanza interna è la vera fonte della sua instabilità e del suo bisogno di risoluzione.

\begin{esempio}[La V7 e il suo tritono in Do Maggiore]
	\label{ex:tritono_v7}
	\musica{
		\new Staff {
			\clef treble
			<sol' si' re'' fa''>1-\markup{V7}
			<si' fa''>1-\markup{"Tritono (5a dim.)"}
		}
	}
\end{esempio}

\begin{regola}[Preparazione e Risoluzione della Dissonanza]
        \label{reg:preparazione_risoluzione}
        Come esposto nel Capitolo \ref{chap:teoria_dissonanza}, ogni nota dissonante segue un ciclo di tre fasi: preparazione, percussione e risoluzione.
        Per la settima di dominante, la regola della preparazione è spesso omessa data la sua familiarità, ma la \textbf{risoluzione è obbligatoria}.
\end{regola}

\begin{fisica}[I battimenti e la tensione del tritono]
	\label{phys:battimenti_tritono}
	
	Il tritono è considerato l'intervallo più dissonante del sistema diatonico. Questa percezione di "ruvidità" e tensione non è solo una convenzione culturale, ma ha una base fisica oggettiva legata al fenomeno dei \textbf{battimenti acustici}.
	
	\textbf{Cosa sono i battimenti}
	
	Quando due onde sonore con frequenze vicine ma non identiche suonano simultaneamente, si verifica un fenomeno di \textbf{interferenza}. Le due onde si sommano, alternando momenti di interferenza \textit{costruttiva} (le onde si rafforzano a vicenda) e \textit{distruttiva} (si attenuano a vicenda). Questo causa una fluttuazione periodica dell'intensità sonora che percepiamo come un "tremolo" o \textbf{battimento}.
	
	La frequenza del battimento è data dalla differenza tra le due frequenze:
	\[
	f_{\text{battimento}} = |f_2 - f_1|
	\]
	
	\textbf{Il tritono e i battimenti multipli}
	
	Consideriamo il tritono Si-Fa. Nel temperamento equabile:
	\begin{itemize}
		\item Si: frequenza $f_1$
		\item Fa: frequenza $f_2 \approx 1.414 f_1$ (rapporto $\sqrt{2}:1$, cioè 6 semitoni)
	\end{itemize}
	
	Il rapporto $\sqrt{2}:1$ è \textbf{irrazionale}: non può essere espresso come frazione di numeri interi. Questo ha conseguenze cruciali per gli armonici:
	
	\begin{itemize}
		\item Gli armonici di Si: $f_1, 2f_1, 3f_1, 4f_1, 5f_1, \ldots$
		\item Gli armonici di Fa: $1.414f_1, 2.828f_1, 4.243f_1, 5.657f_1, \ldots$
	\end{itemize}
	
	Notiamo che gli armonici superiori di queste due note \textit{non coincidono mai} ma cadono sempre in posizioni intermedie, creando multiple coppie di frequenze vicine. Per esempio:
	\begin{itemize}
		\item $2f_1$ (secondo armonico di Si) vs. $2.828f_1$ (secondo armonico di Fa) $\rightarrow$ differenza $\approx 0.828f_1$
		\item $4f_1$ vs. $4.243f_1$ $\rightarrow$ differenza $\approx 0.243f_1$
	\end{itemize}
	
	Ogni coppia di armonici vicini ma non coincidenti genera il proprio battimento. Il risultato è una \textbf{ruvidità acustica complessa} con battimenti a diverse frequenze che si sovrappongono, creando quella sensazione di instabilità e tensione che identifichiamo con la dissonanza.
	
	\textbf{Confronto con intervalli consonanti}
	
	Confrontiamo con la quinta giusta (Sol-Re in Do maggiore):
	\begin{itemize}
		\item Rapporto $3:2$ (razionale, semplice)
		\item Gli armonici tendono a coincidere: il 3° armonico di Sol coincide con il 2° armonico di Re
		\item Pochi o nessun battimento $\rightarrow$ suono "pulito" e stabile
	\end{itemize}
	
	\textbf{La risoluzione come riduzione dei battimenti}
	
	Quando la V7 risolve sulla tonica, il tritono si muove verso intervalli più consonanti:
	\begin{itemize}
		\item Si $\rightarrow$ Do (sensibile che sale)
		\item Fa $\rightarrow$ Mi (settima che scende)
		\item Il tritono Si-Fa (dissonante) diventa la terza Do-Mi o la sesta Mi-Do (consonanti)
	\end{itemize}
	
	Questo movimento rappresenta una \textbf{transizione da rapporti complessi a rapporti semplici}, ovvero da molti battimenti a pochi o nessun battimento. L'orecchio percepisce questa transizione come un passaggio dalla tensione al riposo.
	
	\textbf{Conclusione}
	
	La tensione del tritono e, per estensione, della settima di dominante, ha una base fisica misurabile: i battimenti generati dall'interferenza tra armonici in rapporti complessi. La risoluzione armonica non è solo una convenzione stilistica, ma corrisponde a una transizione fisica da uno stato di alta interferenza acustica a uno di bassa interferenza.
\end{fisica}

\section{Risoluzione naturale (V7 - I)}
\label{sec:risoluzione_v7}

La risoluzione naturale della settima di dominante è sull'accordo di tonica. Le note che formano il tritono hanno dei movimenti melodici obbligati che non possono essere violati.

\begin{regola}[Vettori obbligati della V7]
	\label{reg:vettori_obbligati}
	Nella risoluzione V7 $\rightarrow$ I, le due note del tritono devono muoversi come segue:
	\begin{itemize}
		\item La \textbf{terza} dell'accordo (la sensibile) \textbf{sale} di semitono per risolvere sulla tonica. (Es. in Do: Si $\nearrow$ Do).
		\item La \textbf{settima} dell'accordo \textbf{scende} di grado per risolvere sulla terza della tonica. (Es. in Do: Fa $\searrow$ Mi).
	\end{itemize}
\end{regola}

Questo doppio movimento obbligato crea un problema pratico: come gestire quattro note che risolvono su una triade di tre? La soluzione risiede nella gestione delle note omesse.

\begin{definizione}[V7 Completa vs. Incompleta]
	\label{def:v7_completa_incompleta}
	Esistono due modi principali per risolvere la V7:
	\begin{enumerate}
		\item \textbf{V7 completa $\rightarrow$ I incompleta:} Si usano tutte e quattro le note della V7. Seguendo i vettori obbligati, l'accordo di tonica risultante sarà privo della quinta e avrà la fondamentale raddoppiata tre volte.
		\item \textbf{V7 incompleta $\rightarrow$ I completa:} Si omette la quinta dalla V7 (la nota meno caratterizzante) e si raddoppia la fondamentale. Questo permette di risolvere su un accordo di tonica completo, con fondamentale, terza e quinta.
	\end{enumerate}
\end{definizione}

\begin{esempio}[Le due risoluzioni della V7]
	\label{ex:due_risoluzioni_v7}
	\textbf{A. V7 completa $\rightarrow$ I incompleta}
	\musica{
		\partilate
		{ \key do \major \time 4/4 }
		{ \freccia fa'2 \glissando mi'2 }
		{ re'2 do'2 }
		{ \freccia si2 \glissando do'2 }
		{ sol,2 do2 }
		\lyricmode { "V7"2 "I"2 }
		\figuremode { <7>2 s2 }
	}
	
	\vspace{1cm}
	\textbf{B. V7 incompleta $\rightarrow$ I completa}
	\musica{
		\partilate
		{ \key do \major \time 4/4 }
		{ \freccia fa'2 \glissando mi'2 }
		{ si2 do'2 }
		{ sol2 sol2 }
		{ sol,2 do2 }
		\lyricmode { "V7 (senza 5a)"2 "I"2 }
		\figuremode { <7>2 s2 }
	}
\end{esempio}

\section{I quattro stati della V7}
\label{sec:stati_v7}

L'accordo di settima di dominante può presentarsi in quattro stati: uno fondamentale e tre rivolti. In ogni stato, le regole di risoluzione del tritono (sensibile $\nearrow$ tonica, settima $\searrow$ terza) rimangono il principio guida.

\subsection{V7 allo stato fondamentale (cifratura: 7)}
\begin{itemize}
    \item \textbf{Al basso:} fondamentale (Sol in Do M).
    \item \textbf{Risoluzione:} → I (fondamentale).
    \item \textbf{Condotta:} 
    \begin{itemize}
        \item Basso (Sol): salta a Do.
        \item Sensibile (Si): sale a Do.
        \item Settima (Fa): scende a Mi.
        \item Quinta (Re): scende a Do (I incompleto) o sale a Mi (I completo se V7 è incompleta).
    \end{itemize}
\end{itemize}
Rimandiamo a \ref{def:v7_completa_incompleta} per il dettaglio sulla completezza.

\subsection{V6/5 — Primo rivolto (cifratura: 6/5)}
\begin{itemize}
    \item \textbf{Al basso:} terza dell'accordo (Si in Do M), ovvero la \textbf{sensibile}.
    \item \textbf{Risoluzione:} → I6 (basso sale a Do).
    \item \textbf{Condotta:} 
    \begin{itemize}
        \item Basso (Si): sale obbligatoriamente a Do.
        \item Settima (Fa): scende obbligatoriamente a Mi.
        \item Quinta (Re) e Fondamentale (Sol): si muovono per grado congiunto verso note reali di I6.
    \end{itemize}
\end{itemize}

\subsection{V4/3 — Secondo rivolto (cifratura: 4/3)}
\begin{itemize}
    \item \textbf{Al basso:} quinta dell'accordo (Re in Do M).
    \item \textbf{Risoluzione:} → I6 (basso sale a Mi).
    \item \textbf{Condotta:} 
    \begin{itemize}
        \item Basso (Re): sale di grado a Mi.
        \item Sensibile (Si): sale a Do.
        \item Settima (Fa): scende a Mi (attenzione ai raddoppi con il basso).
    \end{itemize}
\end{itemize}

\subsection{V2 — Terzo rivolto (cifratura: 2 o 4/2)}
\begin{itemize}
    \item \textbf{Al basso:} settima dell'accordo (Fa in Do M).
    \item \textbf{Risoluzione:} → I6 (basso scende a Mi).
    \item \textbf{Condotta:} 
    \begin{itemize}
        \item Basso (Fa): scende obbligatoriamente a Mi.
        \item Sensibile (Si): sale a Do.
    \end{itemize}
\end{itemize}

\section{Sintesi dei rivolti della V7}
\label{sec:sintesi_rivolti_v7}

\begin{center}
\begin{tabular}{|l|l|l|l|l|}
\hline
Stato & Cifratura & Nota al basso & Risolve su & Basso obbligato \\
\hline
Fond. & 7 & Fondamentale & I & No \\
I rv. & 6/5 & Sensibile & I6 & Sì (sale) \\
II rv. & 4/3 & Quinta & I6 & No \\
III rv. & 2 & Settima & I6 & Sì (scende) \\
\hline
\end{tabular}
\end{center}

L'applicazione sistematica dei rivolti della V7 all'armonizzazione della scala è trattata nel Capitolo \ref{chap:regola_ottava_completa} sulla Regola dell'Ottava.

\begin{fisica}[La doppia tensione della V7 nel modo minore]
	Nel modo minore, la settima di dominante (ad esempio Mi-\Soldiesis -Si-Re in La minore) presenta una particolare intensità acustica dovuta alla sovrapposizione di due tipi di tensione:
	
	\begin{enumerate}
		\item \textbf{Tensione melodica:} La sensibile alterata (\Soldiesis  verso La) crea un semitono ascendente con forte attrazione gravitazionale
		\item \textbf{Tensione armonica:} La settima (Re verso Do) crea un'altra attrazione per semitono discendente
		\item \textbf{Tensione intervallare:} L'intervallo di quinta diminuita \Soldiesis -Re (tritono) è acusticamente instabile
	\end{enumerate}
	
	Dal punto di vista fisico, il tritono produce \textbf{battimenti} percepibili a causa della vicinanza dei suoi armonici superiori. Il rapporto di frequenza del tritono temperato ($2^{6/12} = \sqrt{2} \approx 1.414$) non corrisponde a nessun rapporto semplice tra numeri interi piccoli, generando continue interferenze tra le onde sonore.
	
	Questa "rugosità acustica" del tritono, combinata con le due tensioni melodiche opposte (\Soldiesis  ↑ e Re ↓), spiega perché la V7 nel modo minore suoni particolarmente drammatica e richieda imperiosamente la risoluzione sulla tonica.
\end{fisica}

\section{Collocazione sulle scale e casi particolari per il basso dato}
\label{sec:v7_basso_dato}

\begin{regola}[Collocazione della V7]
	\label{reg:collocazione_v7}
	La settima di dominante si costruisce sul \textbf{V grado} sia del modo maggiore sia del modo minore (armonico e melodico ascendente). È un accordo \textbf{naturale ed indipendente}: \textbf{non richiede la preparazione} della dissonanza (Capitolo \ref{chap:teoria_dissonanza}), perché la sua risoluzione è talmente direzionale da giustificarne l'introduzione "di petto".
\end{regola}

\begin{osservazione}[La settima nel modo minore]
	Nel modo maggiore la settima scende di \textbf{semitono} sulla tonica (es. in Do maggiore: Fa $\searrow$ Mi); nel modo minore scende invece di \textbf{tono} (es. in La minore: Re $\searrow$ Do). Nell'armonizzazione del basso dato in tonalità minore, ogni qualvolta la V7 risolve, la terza della tonica di arrivo si raggiunge scendendo di un tono.
\end{osservazione}

\subsection{Risoluzione d'inganno (V7 $\rightarrow$ VI)}
\label{sec:v7_inganno}

Oltre alla risoluzione naturale sulla tonica, la V7 può risolvere sul \textbf{VI grado} (cadenza d'inganno, Capitolo \ref{chap:cadenze}). Le note obbligate del tritono continuano a muoversi come nella risoluzione naturale: la sensibile sale alla tonica, la settima scende alla quinta del VI grado.

\begin{regola}[V7 $\rightarrow$ VI]
	\label{reg:v7_inganno}
	La risoluzione d'inganno della V7 sul VI grado è ammessa \textbf{purché la triade di destinazione sia allo stato completo}. La sensibile sale comunque alla tonica (terza del VI grado), la settima scende di grado sulla quinta del VI grado.
\end{regola}

\begin{esempio}[V7 completa $\rightarrow$ VI in Do Maggiore]
	\label{ex:v7_inganno}
	La sensibile (Si) sale alla tonica (Do), la settima (Fa) scende sulla quinta del VI grado (Mi).
	\musica{
		\partilate
		{ \key do \major \time 4/4 }
		{ re'2 mi'2 }
		{ si2 do'2 }
		{ fa2 mi2 }
		{ sol,2 la,2 }
		\lyricmode { "V7"2 "VI"2 }
		\figuremode { <7>2 s2 }
	}
\end{esempio}

\subsection{La V7 nelle cadenze}
\label{sec:v7_cadenze}

La V7 trova applicazione al posto della triade di dominante all'interno delle formule cadenzali (Capitolo \ref{chap:cadenze}):

\begin{itemize}
	\item \textbf{Cadenza composta consonante}: $I^6_4 - V^7$;
	\item \textbf{Cadenza composta dissonante}: $V^5_4 - V^7_3$;
	\item \textbf{Cadenza doppia}: $V - I^6_4 - V^5_4 - V^7_3$.
\end{itemize}

\begin{osservazione}[Eccezione di Bach]
	\label{oss:eccezione_bach}
	Se la sensibile si trova in una \textbf{voce interna} e si muove per \textbf{moto contrario} rispetto al basso, può eccezionalmente risolvere \textbf{saltando di terza discendente} sul I grado, per permettere alla voce superiore di raggiungere la tonica.
\end{osservazione}
  % La settima di dominante
	\chapter{I ritardi}
\label{chap:ritardi}

Oltre alle dissonanze "strutturali" contenute negli accordi di settima, l'armonia si arricchisce attraverso l'uso di dissonanze melodiche, ovvero note estranee all'accordo che creano una tensione momentanea prima di risolvere. Il \textbf{ritardo} è una delle figure più importanti ed espressive di questa categoria.

\section{Definizione e meccanismo del ritardo}
\label{sec:definizione_ritardo}

Il ritardo è una dissonanza che si manifesta sul tempo forte di una misura (o su una parte forte di un tempo), ritardando, appunto, l'arrivo di una nota reale dell'accordo. La sua efficacia espressiva deriva dal rispetto di un rigoroso processo temporale in tre fasi.

\begin{definizione}[Le tre fasi del Ritardo]
	\label{def:fasi_ritardo}
	Un ritardo si articola sempre in tre momenti distinti:
	\begin{enumerate}
		\item \textbf{Preparazione:} La nota che diventerà dissonante viene introdotta come nota consonante nell'accordo precedente, preferibilmente nella stessa voce e alla stessa altezza.
		\item \textbf{Percussione:} L'accordo cambia, ma la nota preparata viene tenuta (legata), diventando così dissonante rispetto al nuovo contesto armonico. Questa fase deve avvenire su un tempo forte.
		\item \textbf{Risoluzione:} La nota dissonante scende di grado congiunto su una nota reale dell'accordo.
	\end{enumerate}
\end{definizione}

\section{Il ritardo 4-3}
\label{sec:ritardo_4_3}

Il ritardo 4-3 è il più celebre e utilizzato. Ritarda l'arrivo della terza dell'accordo: sulla dominante ritarda la sensibile (Si in Do maggiore), sulla tonica ritarda la terza dell'accordo (Mi in Do maggiore).

\subsection{Il ritardo nella cadenza perfetta}
\label{subsec:ritardo_cadenza}

Il ritardo 4-3 trova la sua applicazione più efficace nella cadenza perfetta (V-I), dove può amplificare la tensione della dominante e ritardare la risoluzione sulla tonica. Come per la settima di dominante, esistono due forme principali di realizzazione.

\begin{esempio}[Forma A: V7 completo → I incompleto con ritardo 4-3 in Do Maggiore]
	\label{ex:ritardo_completo_maggiore}
	Questa progressione utilizza tutti e quattro i suoni sulla dominante (compresa la settima) e risolve su un accordo di tonica incompleto (senza quinta), con raddoppio della fondamentale e ritardo della terza.
	
	\musica{es_cap10_01}
	
	La preparazione del ritardo avviene nell'accordo di IV: il Fa (nel contralto) è nota reale dell'accordo. Con l'arrivo del V7 (Sol al basso) il Fa è tenuto (legato) e diventa la settima; il soprano scende alla sensibile Si e il tenore salta al Re (quinta). Nella seconda battuta, sul basso Do, il Fa è percosso come ritardo 4 e risolve su Mi (3), mentre la sensibile sale a Do e il Re scende a Do.
\end{esempio}

\begin{esempio}[Forma B: V7 incompleto → I completo con ritardo 4-3 in Do Maggiore]
	\label{ex:ritardo_incompleto_maggiore}
	Questa progressione omette la quinta dalla dominante, permettendo di risolvere su un accordo di tonica completo con ritardo della terza.
	
	\musica{es_cap10_02}
	
	In questa versione la dominante è incompleta (senza Re, la quinta) e la fondamentale è raddoppiata; la settima (Fa) è tenuta dal contralto e diventa il ritardo. L'accordo di tonica risultante è completo, con fondamentale, terza (ritardata) e quinta.
\end{esempio}

\subsection{Il ritardo nel modo minore}
\label{subsec:ritardo_minore}

Nel modo minore, il ritardo 4-3 segue gli stessi principi, ma si adatta alle caratteristiche della scala minore armonica, in particolare all'uso della sensibile innalzata nel V7.

\begin{esempio}[Forma A: V7 completo → I incompleto con ritardo 4-3 in La minore]
	\label{ex:ritardo_completo_minore}
	Progressione con V7 completo che risolve su I incompleto con ritardo.
	
	\musica{es_cap10_03}
\end{esempio}

\begin{esempio}[Forma B: V7 incompleto → I completo con ritardo 4-3 in La minore]
	\label{ex:ritardo_incompleto_minore}
	Progressione con V7 incompleto che risolve su I completo con ritardo.
	
	\musica{es_cap10_04}
\end{esempio}

\begin{regola}[Cifratura del Ritardo]
	\label{reg:cifratura_ritardo}
	Nel basso figurato, il ritardo 4-3 si indica scrivendo in sequenza \texttt{4 3}. Il numero 4 rappresenta la dissonanza (la quarta rispetto al basso), mentre il 3 indica la sua risoluzione (la terza dell'accordo).
\end{regola}

Gli esempi mostrano come il ritardo 4-3 possa arricchire la progressione cadenzale, creando uno slancio melodico nelle voci superiori mentre il basso mantiene la sua funzione armonica di sostegno. La scelta tra forma completa e incompleta dipende dal contesto musicale e dalla necessità di preservare o meno tutte le note dell'accordo di arrivo.  % I ritardi
	\chapter{Le note estranee all'armonia}
\label{chap:noteestranee}

Questo capitolo tratta tutte le note che appartengono al tessuto melodico ma non all'accordo sottostante. Queste note, dette "estranee" o "non armoniche", arricchiscono la trama musicale rendendo le melodie più fluide e cantabili.

\begin{definizione}[Nota estranea all'armonia]
    \label{def:nota_estranea}
    Una nota si definisce estranea (o non armonica) quando non fa parte dell'accordo in quel momento presente nel basso. Esse creano dissonanze temporanee che risolvono su note reali dell'accordo.
\end{definizione}

\section{Nota di passaggio (NdP)}
\label{sec:nota_passaggio}

\begin{definizione}[Nota di passaggio]
    È una nota estranea che occupa per grado congiunto l'intervallo tra due note reali dell'accordo, consecutive o appartenenti ad accordi diversi.
\end{definizione}

\begin{regola}[Uso della NdP]
    La nota di passaggio cade tipicamente sul tempo debole (o sulla parte debole del tempo). Se posta sul tempo forte, assume un carattere più teso simile all'appoggiatura. Si distingue tra NdP \textbf{diatonica} (nella scala) e \textbf{cromatica} (alterazione).
\end{regola}

\begin{esempio}[Nota di passaggio in Do maggiore]
    Melodia Do-Re-Mi sopra un accordo di Do maggiore: Re è la nota di passaggio.
    \musica{
        \new Staff {
            \key do \major \time 4/4
            \relative do'' {
                do4 \passaggio re mi2 |
            }
        }
    }
\end{esempio}

\section{Nota di volta (o cambiata)}
\label{sec:nota_volta}

\begin{definizione}[Nota di volta]
    È una nota che si allontana dalla nota reale per grado congiunto (superiore o inferiore) e vi ritorna immediatamente.
\end{definizione}

\begin{regola}
    La nota di volta cade solitamente sul tempo debole. La nota di ritorno deve essere identica alla nota reale di partenza.
\end{regola}

\section{Anticipazione}
\label{sec:anticipazione}

\begin{definizione}[Anticipazione]
    È una nota che appartiene all'accordo \textit{successivo} ma appare prima del cambio armonico, sul tempo debole.
\end{definizione}

\begin{regola}
    A differenza del ritardo, l'anticipazione non viene preparata né risolve: essa \textit{è} già la risoluzione anticipata.
\end{regola}

\begin{esempio}[Anticipazione nella cadenza V-I]
    La tonica al soprano appare sull'ultima croma del V7.
    \musica{
        \new Staff {
            \key do \major \time 4/4
            \relative do'' {
                \voiceOne
                <si fa'>4. do8 do2 |
            }
        }
    }
\end{esempio}

\section{Appoggiatura}
\label{sec:appoggiatura}

\begin{definizione}[Appoggiatura]
    È una nota estranea posta sul tempo \textbf{forte}, raggiunta per salto e risolta per grado congiunto discendente (più raramente ascendente).
\end{definizione}

\begin{regola}[Appoggiatura vs Ritardo]
    A differenza del ritardo (\ref{chap:ritardi}), l'appoggiatura non è preparata da una nota tenuta. Storicamente era spesso scritta come abbellimento, ma ha il peso metrico di una nota reale.
\end{regola}

\section{Nota di sfuggita (Échappée)}
\label{sec:sfuggita}

\begin{definizione}[Nota di sfuggita]
    È una nota estranea raggiunta per grado congiunto ma abbandonata per salto in direzione opposta. È comune nel soprano in movimenti melodici veloci.
\end{definizione}

\section{Bordone e nota pedale}
\label{sec:pedale}

\begin{definizione}[Nota pedale]
    È una nota tenuta (solitamente al basso o al soprano) mentre l'armonia cambia sopra o sotto di essa.
\end{definizione}

\begin{regola}
    La nota pedale deve essere consonante con l'accordo all'inizio e alla fine del passaggio. Le dissonanze centrali sono ammesse per l'effetto di accumulo di tensione.
\end{regola}

\section{Sintesi delle note estranee}
\label{sec:sintesi_note_estranee}

\begin{center}
\begin{tabular}{|l|l|l|l|l|}
\hline
Nome & Tempo metrico & Raggiunta per... & Abbandonata per... & Preparata? \\
\hline
Passaggio & Debole & Grado congiunto & Grado congiunto & No \\
Volta & Debole & Grado congiunto & Ritorno & No \\
Anticipazione & Debole & Qualsiasi & Nota reale & No (è già risoluz.) \\
Appoggiatura & Forte & Salto & Grado congiunto & No \\
Sfuggita & Debole & Grado congiunto & Salto opposto & No \\
Pedale & Entrambi & Tenuta & Tenuta & Sì \\
\hline
\end{tabular}
\end{center}
 % Le note estranee all'armonia
	\chapter{VII7 nel modo maggiore e II7 nel modo minore: stessi suoni, funzioni diverse}
\label{chap:vii7_ii7}

Uno degli aspetti più affascinanti dell'armonia tonale è la capacità di uno stesso insieme di note di assumere significati e funzioni completamente diverse a seconda del contesto tonale. In questo capitolo esploriamo due accordi di settima che, pur essendo costituiti dagli stessi suoni (in tonalità relative), svolgono ruoli armonici distinti e complementari: la \textbf{settima di sensibile} (VII7) nel modo maggiore e la \textbf{settima sul secondo grado} (II7) nel modo minore.

\section{La settima di sensibile nel modo maggiore (VII7)}
\label{sec:vii7_maggiore}

Nel modo maggiore, costruendo un accordo di settima sul settimo grado della scala, otteniamo una quadriade di grande interesse armonico.

\begin{definizione}[Settima di sensibile (VII7)]
	\label{def:vii7}
	La \textbf{settima di sensibile} è un accordo di settima costruito sul VII grado della scala maggiore. È composto da una triade diminuita (fondamentale, terza minore, quinta diminuita) con l'aggiunta di una settima minore. In Do maggiore: Si-Re-Fa-La.
	
	Questo accordo è anche chiamato \textbf{semidiminuito} o \textbf{settima di seconda specie}, indicato con il simbolo VII\textsuperscript{\o 7} o VII7 con la cifratura {\o}.
\end{definizione}

La caratteristica fondamentale di questo accordo è che contiene tre note della settima di dominante, del quale manca la fondamentale. È naturale dunque che questo accordo abbia funzione di dominante e risolva sulla tonica.

\begin{regola}[Funzione e risoluzione della VII7]
	\label{reg:funzione_vii7}
	La VII7 ha funzione di \textbf{dominante}. La sua risoluzione naturale è sulla tonica (I). Secondo le regole generali di risoluzione della dissonanza (Capitolo \ref{chap:teoria_dissonanza}), la settima deve scendere di grado. Le due note che formano il tritono (Si-Fa in Do maggiore) seguono il loro movimento obbligato:
	\begin{itemize}
		\item La sensibile (Si) sale alla tonica (Do)
		\item La quinta diminuita (Fa) scende alla terza della tonica (Mi)
		\item La settima (La) scende anch'essa di grado (Sol)
	\end{itemize}
\end{regola}

\begin{esempio}[VII7 in Do Maggiore e sua risoluzione]
	\label{ex:vii7_maggiore}
	L'accordo di settima di sensibile in stato fondamentale risolve sulla tonica.
	\musica{
		\partilate
		{ \key do \major \time 4/4 }
		{ la'2 sol'2 }
		{ fa'2 mi'2 }
		{ re'2 do'2 }
		{ si,2 do2 }
		\lyricmode { "VIIø7"2 "I"2 }
		\figuremode { <7>2 s2 }
	}
\end{esempio}

\subsection{I rivolti della VII7 e la Regola dell'Ottava}
\label{subsec:rivolti_vii7}

Come tutti gli accordi di settima, anche la VII7 può essere rivoltata. Il suo uso più comune è nel \textbf{primo rivolto} (VII6/5), che appare nella Regola dell'Ottava.

\begin{regola}[VII6/5 nella Regola dell'Ottava discendente]
	\label{reg:vii65_regola}
	Nella scala discendente del modo maggiore, quando il basso raggiunge il VII grado, si utilizza il primo rivolto della settima di sensibile (VII6/5). Questo crea una linea del basso fluida e permette alla settima (che si trova ora al soprano o in una voce intermedia) di risolvere regolarmente scendendo di grado.
\end{regola}

\begin{esempio}[VII6/5 nella Regola dell'Ottava - Do Maggiore]
	\label{ex:vii65_regola_maggiore}
	Frammento della scala discendente con evidenziazione del VII6/5.
	\musica{
		\partilate
		{ \key do \major \time 4/4 }
		{ do''4 si'4 la'4 sol'4 }
		{ mi'4 re'4 mi'4 re'4 }
		{ sol4 sol4 do'4 si4 }
		{ do4 si,4 la,4 sol,4 }
		\lyricmode { "I"4 "VII"4 "VI"4 "V"4 }
		\figuremode { \unset FiguredBass.useBassFigureExtenders
			s4 <6 5>4 <6>4 s4 }
	}
\end{esempio}

\section{La settima sul secondo grado nel modo minore (II7)}
\label{sec:ii7_minore}

Nel modo minore armonico, l'accordo di settima costruito sul secondo grado della scala ha la stessa struttura della VII7 del modo maggiore: una triade diminuita con settima minore (accordo semidiminuito).

\begin{definizione}[Settima sul secondo grado (II7)]
	\label{def:ii7_minore}
	La \textbf{settima sul secondo grado} nel modo minore armonico è un accordo semidiminuito che svolge funzione di \textbf{sottodominante}. In La minore: Si-Re-Fa-La (gli stessi suoni della VII7 di Do maggiore!).
	
	Si indica come II\textsuperscript{\o 7} o semplicemente II7.
\end{definizione}

Tale accordo, come visto nel paragrafo \ref{def:vii7}, contiene tre di quattro note della settima di dominante del modo relativo maggiore e porta alla sua tonica; nel contesto del modo minore, la tonica del modo maggiore corrisponde alla dominante del modo minore: l'accordo II7 nel modo minore, dunque, risolve sulla dominante.

\subsection{La progressione I - II7 - V - I nel modo minore}
\label{subsec:progressione_ii7}

L'uso più efficace della II7 nel modo minore è all'interno della progressione armonica fondamentale che esprime il ciclo Tonica-Sottodominante-Dominante-Tonica.

\begin{regola}[Preparazione della settima sul II grado]
	\label{reg:preparazione_ii7}
	La settima dell'accordo II7, secondo i principi generali di preparazione della dissonanza (cfr. Capitolo \ref{chap:teoria_dissonanza}), deve essere preparata come consonanza nell'accordo precedente (I). La condotta delle voci più fluida si ottiene utilizzando il \textbf{primo rivolto} dell'accordo II7, ponendo la terza al basso (Re).
\end{regola}

\begin{esempio}[Progressione I - II6/5 - V - I in La minore]
	\label{ex:progressione_ii7_minore}
	La settima (La) viene preparata nell'accordo di tonica e risolve scendendo di grado (Sol) nell'accordo di dominante.
	\musica{
		\partilate
		{ \key la \minor \time 4/4 }
		{ do'2 si2 | si1 }
		{ la2 la2 | sold1 }
		{ mi2 fa2 | mi1 }
		{ la,2 re2 | mi2 mi,2 | la,1 }
		\lyricmode { "I"2 "IIø7"2 | "V"1 }
		\figuremode { s2 <6 5>2 | s1 }
	}
\end{esempio}

\subsection{II6/5 nella Regola dell'Ottava (modo minore)}
\label{subsec:ii7_regola_minore}

Analogamente al modo maggiore, anche nel modo minore il secondo grado della scala ascendente può essere armonizzato con il primo rivolto della settima (II6/5).

\begin{regola}[II6/5 nella Regola dell'Ottava ascendente - Modo Minore]
	\label{reg:ii65_regola_minore}
	Nella scala ascendente del modo minore, il II grado al basso viene armonizzato con l'accordo II6/5. La settima deve essere preparata e risolve regolarmente scendendo di grado nell'accordo successivo (III o V).
\end{regola}

\begin{esempio}[II6/5 nella Regola dell'Ottava - La minore]
	\label{ex:ii65_regola_minore}
	Frammento della scala ascendente in La minore.
	\musica{
		\partilate
		{ \key la \minor \time 4/4 }
		{ do'4 re'4 mi'4 fa'4 }
		{ la4 la4 si4 la4 }
		{ mi4 fa4 sold4 do'4 }
		{ la,4 si,4 do4 re4 }
		\lyricmode { "I"4 "II"4 "III"4 "IV"4 }
		\figuremode { s4 <6 5>4 <6>4 s4 }
	}
\end{esempio}

\section{Sintesi: VII7 (Magg) vs II7 (Min)}
\label{sec:sintesi_vii_ii}

\begin{center}
	\begin{tabular}{|l|c|c|}
		\hline
		\textbf{Proprietà} & \textbf{VII7 (Modo Maggiore)} & \textbf{II7 (Modo Minore)} \\
		\hline
		Struttura & Semidiminuita (m-m-M) & Semidiminuita (m-m-M) \\
		Funzione & Dominante & Sottodominante \\
		Risoluzione naturale & Tonica (I) & Dominante (V) \\
		Preparazione settima & Preferibile & Obbligatoria \\
		Note (es. Do M / La m) & Si-Re-Fa-La & Si-Re-Fa-La \\
		\hline
	\end{tabular}
\end{center}
  % VII7 nel modo maggiore e II7 nel modo minore
	\chapter{Le sette specie di settima}
\label{chap:sette_specie}

Dopo aver studiato la settima di dominante e i suoi derivati diretti, è necessario esplorare l'intera famiglia degli accordi di quadriade. In questo capitolo sistematizzeremo le diverse "specie" di settima, ovvero le diverse combinazioni di terze sovrapposte che possono nascere sui gradi della scala.

Per quanto riguarda il comportamento generale delle dissonanze, la loro preparazione e la loro risoluzione obbligatoria, si rimanda integralmente al \textbf{Capitolo \ref{chap:teoria_dissonanza}}.

\section{Classificazione generale}
\label{sec:classificazione_specie}

Tradizionalmente si distinguono sette specie di accordi di settima, determinate dalla natura della triade di base e dalla qualità dell'intervallo di settima rispetto alla fondamentale.

\begin{enumerate}
    \item \textbf{Prima specie:} Triade Maggiore + settima minore (Settima di dominante).
    \item \textbf{Seconda specie:} Triade minore + settima minore.
    \item \textbf{Terza specie:} Triade diminuita + settima minore (Settima semidiminuita).
    \item \textbf{Quarta specie:} Triade Maggiore + settima Maggiore.
    \item \textbf{Quinta specie:} Triade diminuita + settima diminuita (Settima diminuita).
    \item \textbf{Sesta specie:} Triade minore + settima Maggiore.
    \item \textbf{Settima specie:} Triade aumentata + settima Maggiore.
\end{enumerate}

In un corso di armonia complementare, ci concentreremo solo sulle casistiche più comuni e rilevanti per l'analisi e la composizione standard, ignorando le specie più rare che non trovano applicazione pratica negli esami.

\section{Prima e Terza specie: richiami}
\label{sec:richiami_specie}

Le specie più importanti sono già state trattate nei capitoli precedenti:

\begin{itemize}
    \item \textbf{Prima Specie (Settima di dominante):} Si rimanda al \textbf{Capitolo \ref{chap:settima_dominante}}. È l'unica specie che non richiede preparazione obbligatoria della settima.
    \item \textbf{Terza Specie (Settima di sensibile / semidiminuita):} Si rimanda al \textbf{Capitolo \ref{chap:vii7_ii7}}. Ricordiamo che questa specie richiede sempre la \textbf{preparazione} della settima (fase di consonanza nell'accordo precedente) e risolve scendendo di grado. La sua funzione è di dominante (sul VII grado del maggiore) o di sottodominante (sul II grado del minore).
\end{itemize}

\section{Seconda Specie}
\label{sec:seconda_specie_dettaglio}

La seconda specie è formata da una \textbf{triade minore} e una \textbf{settima minore}. L'accordo di gran lunga più importante di questa specie è quello costruito sul \textbf{secondo grado (II)} del modo maggiore.

\subsection{Funzione e risoluzione}
La sua funzione è tipicamente di \textbf{sottodominante}. Essendo una dissonanza di seconda specie, la settima deve essere rigorosamente \textbf{preparata}. Risolve naturalmente sulla dominante (V o V7).

\begin{regola}[Collegamento I-II7-V e raddoppio della terza]
    \label{reg:ii7_raddoppio}
    Nel collegamento tra il I grado e il II7 (e la successiva risoluzione su V), è frequente e consigliabile \textbf{raddoppiare la terza} dell'accordo di II7 (che corrisponde al IV grado della scala). Questo espediente tecnico è fondamentale per evitare il rischio di \textbf{quinte parallele} tra la fondamentale e la quinta del II grado nel movimento verso l'accordo di dominante.
\end{regola}

\begin{esempio}[II7 in Do Maggiore]
    \musica{
        \partilate
        { \key do \major \time 4/4 }
        { do''2 do''2 | si'1 }
        { sol'2 fa'2 | re'1 }
        { mi'2 fa'2 | fa'1 }
        { do2 re2 | sol,1 }
        \lyricmode { "I"2 "II7"2 | "V"1 }
        \figuremode { s2 <7>2 | s1 }
    }
    Si noti il raddoppio del Fa (terza del II7) tra contralto e tenore.
\end{esempio}

\section{Quarta Specie}
\label{sec:quarta_specie_dettaglio}

La quarta specie è formata da una \textbf{triade Maggiore} e una \textbf{settima Maggiore}. Nel corso di armonia studieremo solo l'accordo costruito sul \textbf{primo grado (I)} del modo maggiore.

\subsection{Uso e risoluzione}
Questo accordo ha una tensione molto particolare e risolve generalmente sulla \textbf{sottodominante (IV)}. La settima (che è la sensibile della scala) deve essere preparata e risolve scendendo di grado verso la terza del IV grado.

\begin{esempio}[I7 della quarta specie]
    \musica{
        \partilate
        { \key do \major \time 4/4 }
        { sol'2 sol'2 | la'1 }
        { mi'2 mi'2 | fa'1 }
        { do'2 si2 | do'1 }
        { do2 do2 | fa,1 }
        \lyricmode { "I"2 "I7 (4a sp.)"2 | "IV"1 }
        \figuremode { s2 <7>2 | s1 }
    }
\end{esempio}

\section{Quinta Specie}
\label{sec:quinta_specie_dettaglio}

La quinta specie è la \textbf{settima diminuita} (triade diminuita + settima diminuita). È l'accordo più dissonante e simmetrico.

\subsection{Uso sul VII grado del modo minore}
Si studia principalmente sul \textbf{VII grado del modo minore} (sulla sensibile), dove svolge funzione di dominante estrema. Come la settima di dominante, può talvolta omettere la preparazione, ma la risoluzione è obbligatoria e tutte le voci tendono a muoversi verso la tonica.

Oltre che sul VII grado, questa specie trova un uso importante come \textbf{accordo alterato sul IV grado} di entrambi i modi (maggiore e minore), agendo come una sorta di "dominante della dominante" o per scopi coloristici nelle cadenze.

\section{Sesta e Settima Specie}
\label{sec:sesta_settima_specie}

Queste due specie sono estremamente rare e non trovano uso pratico in un corso di armonia complementare. È sufficiente sapere dove si collocano teoricamente:

\begin{itemize}
    \item \textbf{Sesta specie (minore-Maggiore):} si trova sul I grado della scala minore armonica e melodica; risolve generalmente sul IV grado.
    \item \textbf{Settima specie (aumentata-Maggiore):} si trova sul III grado della scala minore armonica e melodica; risolve sul VI grado.
\end{itemize}

Nonostante la loro esistenza teorica, non verranno approfondite ulteriormente in quanto non richieste nelle prove d'esame.
  % Le sette specie di settima
	% ============================================================
% NUOVO CAPITOLO DA AGGIUNGERE - REGOLA DELL'OTTAVA COMPLETA
% Questo capitolo andrebbe inserito dopo il Capitolo J
% ============================================================

\chapter{La Regola dell'Ottava: versione completa}
\label{chap:regola_ottava_completa}

Dopo aver esplorato la teoria generale delle settime nel Capitolo \ref{chap:sette_specie}, siamo finalmente pronti a presentare la \textbf{versione completa e definitiva} della Regola dell'Ottava. Questa versione integra l'intero arsenale delle settime diatoniche, offrendo una tavolozza armonica ricchissima per l'armonizzazione della scala al basso.

\section{Le tre versioni a confronto}
\label{sec:tre_versioni_confronto}

Facciamo un breve riepilogo del percorso seguito:

\begin{definizione}[Le tre versioni della Regola dell'Ottava]
	\label{def:tre_versioni_complete}
	
	\begin{enumerate}
		\item \textbf{Versione preliminare (Capitolo \ref{chap:funzioni_armoniche}):}
		\begin{itemize}
			\item Utilizza solo \textbf{triadi} in stato fondamentale e rivolti
			\item Armonizzazione "sicura" e consonante
			\item Ideale per l'apprendimento iniziale della condotta delle voci
		\end{itemize}
		
		\item \textbf{Versione intermedia (Capitolo \ref{chap:settima_dominante}):}
		\begin{itemize}
			\item Aggiunge la \textbf{settima di dominante (V7)} e i suoi tre rivolti
			\item Introduce la dissonanza controllata
			\item Maggiore tensione e direzionalità armonica
		\end{itemize}
		
		\item \textbf{Versione completa (questo capitolo):}
		\begin{itemize}
			\item Include \textbf{tutte le settime diatoniche}: II7, IV7, V7, VII7
			\item Massima ricchezza armonica e varietà di colori
			\item Riflette la pratica compositiva del tardo Barocco e primo Classicismo
		\end{itemize}
	\end{enumerate}
\end{definizione}

\begin{matematica}[Crescita combinatoria delle possibilità armoniche]
	Analizziamo quantitativamente come cresce la complessità passando da una versione all'altra:
	
	\textbf{Versione preliminare:}
	\begin{itemize}
		\item 7 triadi diatoniche
		\item 3 stati per ogni triade (fondamentale + 2 rivolti)
		\item $7 \times 3 = 21$ possibili configurazioni di base
	\end{itemize}
	
	\textbf{Versione intermedia:}
	\begin{itemize}
		\item 21 configurazioni precedenti
		\item +1 accordo di settima (V7) con 4 stati (fondamentale + 3 rivolti)
		\item $21 + 4 = 25$ configurazioni disponibili
	\end{itemize}
	
	\textbf{Versione completa:}
	\begin{itemize}
		\item Configurazioni precedenti
		\item +6 nuove settime diatoniche (I7, II7, III7, IV7, VI7, VII7) × 4 stati ciascuna
		\item Totale teorico: $21 + (7 \times 4) = 49$ configurazioni
	\end{itemize}
	
	La versione completa offre più del doppio delle possibilità armoniche rispetto alla versione preliminare!
\end{matematica}

\section{La Regola dell'Ottava completa nel modo maggiore}
\label{sec:regola_completa_maggiore}

Presentiamo ora lo schema completo per il modo maggiore, integrando tutte le varianti con settime.

\begin{regola}[Regola dell'Ottava Completa - Modo Maggiore]
	\label{reg:regola_ottava_completa_maggiore}
	
	\textbf{Scala ascendente (esempio in Do maggiore):}
	
	\begin{tabular}{|c|c|c|c|}
		\hline
		\textbf{Grado} & \textbf{Basso} & \textbf{Numerica} & \textbf{Funzione armonica} \\
		\hline
		I & Do & 5 & I (Tonica) \\
		II & Re & 6 \textit{oppure} 4/3 & VII6 \textit{oppure} V4/3 \\
		III & Mi & 6 & I6 (Tonica) \\
		IV & Fa & 6 \textit{oppure} 6/5 & II6 \textit{oppure} II6/5 (Sopratonica 7) \\
		V & Sol & 5 \textit{oppure} 7 & V \textit{oppure} V7 (Dominante) \\
		VI & La & 6 & IV6 (Sottodominante) \\
		VII & Si & 6 \textit{oppure} 6/5 & V6 \textit{oppure} V6/5 (Dominante 7) \\
		VIII & Do & 5 & I (Tonica) \\
		\hline
	\end{tabular}
	
	\vspace{0.5cm}
	
	\textbf{Scala discendente (esempio in Do maggiore):}
	
	\begin{tabular}{|c|c|c|c|}
		\hline
		\textbf{Grado} & \textbf{Basso} & \textbf{Numerica} & \textbf{Funzione armonica} \\
		\hline
		VIII & Do & 5 & I (Tonica) \\
		VII & Si & 6 \textit{oppure} 6/5 & V6 \textit{oppure} V6/5 (Dominante 7) \\
		VI & La & 4/3\# & V4/3\# (Sesta Eccedente, \ref{sec:sesta_eccedente}) \\
		V & Sol & 5 \textit{oppure} 7 & V \textit{oppure} V7 (Dominante) \\
		IV & Fa & 2 & V2 (Dominante terza inversione) \\
		III & Mi & 6 & I6 (Tonica) \\
		II & Re & 4/3 & V4/3 (Dominante seconda inversione) \\
		I & Do & 5 & I (Tonica) \\
		\hline
	\end{tabular}
\end{regola}

\begin{esempio}[Regola dell'Ottava Completa in Do Maggiore - scala discendente con tutte le settime]
	\label{ex:regola_completa_disc_magg}
	\musica{es_cap14_01}
	Nota: Sul VI grado discendente (La), la cifratura 4/3\# indica l'uso del II4/3 con la quarta aumentata (caratteristica della dominante secondaria "V del V"), che può essere interpretata anche come un accordo di sesta eccedente.
\end{esempio}

\section{La Regola dell'Ottava completa nel modo minore}
\label{sec:regola_completa_minore}

Nel modo minore, la versione completa presenta alcune particolarità interessanti legate alla natura della scala minore armonica.

\begin{regola}[Regola dell'Ottava Completa - Modo Minore]
	\label{reg:regola_ottava_completa_minore}
	
	\textbf{Scala ascendente (esempio in La minore armonico):}
	
	\begin{tabular}{|c|c|c|c|}
		\hline
		\textbf{Grado} & \textbf{Basso} & \textbf{Numerica} & \textbf{Funzione armonica} \\
		\hline
		I & La & 5 & i (Tonica) \\
		II & Si & 4/3 & V4/3 (Dominante) \\
		III & Do & 6 & i6 (Tonica) \\
		IV & Re & 6 \textit{oppure} 6/5 & ii°6 \textit{oppure} ii°6/5 (Sopratonica 7) \\
		V & Mi & 5 \textit{oppure} 7 & V \textit{oppure} V7 (Dominante) \\
		VI & Fa & 6 & iv6 (Sottodominante) \\
		VII & \Soldiesis  & 6/5 & V6/5 (Dominante 7) \\
		VIII & La & 5 & i (Tonica) \\
		\hline
	\end{tabular}
	
	\vspace{0.5cm}
	
	\textbf{Scala discendente (esempio in La minore):}
	
	\begin{tabular}{|c|c|c|c|}
		\hline
		\textbf{Grado} & \textbf{Basso} & \textbf{Numerica} & \textbf{Funzione armonica} \\
		\hline
		VIII & La & 5 & i (Tonica) \\
		VII & \Soldiesis  & 6 & V6 (Dominante) \\
		VI & Fa & \#6 & iv\#6 (Sesta Eccedente, \ref{sec:sesta_eccedente}) \\
		V & Mi & 5 \textit{oppure} 7 & V \textit{oppure} V7 (Dominante) \\
		IV & Re & 2 & V2 (Dominante terza inversione) \\
		III & Do & 6 & i6 (Tonica) \\
		II & Si & 4/3 & V4/3 (Dominante seconda inversione) \\
		I & La & 5 & i (Tonica) \\
		\hline
	\end{tabular}
\end{regola}

\begin{esempio}[Regola dell'Ottava Completa in La minore - scala ascendente]
	\label{ex:regola_completa_asc_min}
	\musica{es_cap14_02}
\end{esempio}

\begin{esempio}[Regola dell'Ottava Completa in La minore - scala discendente]
	\label{ex:regola_completa_disc_min}
	\musica{es_cap14_03}
	Nota: Sul VI grado discendente (Fa), la cifratura \#6 (sesta eccedente) indica la caratteristica alterazione ascendente della sesta (\Rediesis ) che risolve espandendosi verso l'ottava (Mi).
\end{esempio}

\section{La sesta eccedente nel contesto della Regola dell'Ottava}
\label{sec:sesta_eccedente}

Nelle tabelle della Regola dell'Ottava completa, sul VI grado discendente appare spesso la cifratura $\sharp 6$ o $4/3\sharp$ con l'etichetta "sesta eccedente". Questo accordo è una delle risorse cromatiche più potenti dell'armonia tonale.

\begin{definizione}[Accordo di sesta eccedente]
    \label{def:sesta_eccedente}
    È un accordo cromatico caratterizzato dall'intervallo di sesta eccedente tra il basso e una voce superiore. Si forma tipicamente tra il VI grado abbassato al basso e il IV grado alzato cromaticamente in una voce superiore.
    La sua funzione principale è preparare la dominante con una spinta cromatica convergente verso la nota della dominante stessa.
\end{definizione}

\subsection{Le tre specie di sesta eccedente}

\begin{itemize}
    \item \textbf{Sesta italiana (It+6):} Formata da VI$\flat$ (basso), tonica e IV$\sharp$ (sesta eccedente). Solo tre note.
    \musica{es_cap14_04}
    \item \textbf{Sesta francese (Fr+6):} Aggiunge il II grado alla sesta italiana. In Do m: \Labemolle -Do-Re-\Fadiesis .
    \musica{es_cap14_05}
    \item \textbf{Sesta tedesca (Ger+6):} Aggiunge il III$\flat$ alla sesta italiana. In Do m: \Labemolle -Do-\Mibemolle -\Fadiesis . Produce quinte parallele se risolve direttamente su V; si usa spesso il I6/4 come ponte (\ref{chap:quarta_sesta}).
    \musica{es_cap14_06}
\end{itemize}

\begin{regola}[Risoluzione della sesta eccedente]
    \label{reg:risoluzione_sesta_eccedente}
    \begin{itemize}
        \item Il basso (VI$\flat$) scende di semitono alla dominante.
        \item La voce con la sesta eccedente (IV$\sharp$) sale di semitono alla dominante.
        \item Le altre voci si muovono per grado congiunto verso il V o verso il I6/4 preparatorio.
    \end{itemize}
\end{regola}

La sesta eccedente appare naturalmente nella Regola dell'Ottava sul VI grado discendente, potenziando la spinta verso la dominante (V grado).

\section{Usi pratici e stilistici}
\label{sec:usi_pratici_completi}

La versione completa della Regola dell'Ottava non è uno schema rigido da applicare meccanicamente, ma una \textbf{tavolozza di possibilità} da cui il compositore attinge in base al contesto musicale.

\begin{regola}[Criteri di scelta tra le versioni]
	\label{reg:criteri_scelta_versioni}
	
	\begin{itemize}
		\item \textbf{Periodo storico:}
		\begin{itemize}
			\item Primo Barocco: preferire la versione preliminare
			\item Tardo Barocco (Bach, Händel): versione intermedia o completa
			\item Classicismo (Mozart, Haydn): versione completa con libertà di scelta
		\end{itemize}
		
		\item \textbf{Contesto espressivo:}
		\begin{itemize}
			\item Passaggi leggeri, danzanti: versione preliminare
			\item Momenti drammatici, intensi: versione intermedia o completa
			\item Sequenze modulanti: versione completa (le settime facilitano la modulazione)
		\end{itemize}
		
		\item \textbf{Velocità del movimento del basso:}
		\begin{itemize}
			\item Basso veloce (semicrome, terzine): preferire triadi semplici
			\item Basso moderato (crome, semiminime): versione intermedia
			\item Basso lento (minime, semibrevi): versione completa con massima ricchezza
		\end{itemize}
		
		\item \textbf{Texture:}
		\begin{itemize}
			\item Contrappunto denso: versione preliminare (evitare troppe dissonanze)
			\item Accompagnamento omofono: versione completa (maggiore libertà armonica)
		\end{itemize}
	\end{itemize}
\end{regola}

\begin{fisica}[Perché le settime "colorano" l'armonia?]
	Dal punto di vista acustico, l'aggiunta di settime agli accordi produce effetti percettivi molto specifici:
	
	\begin{enumerate}
		\item \textbf{Rugosità timbrica:} Gli intervalli di settima (maggiore o minore) creano battimenti tra armonici parziali, producendo una texture sonora più "densa" e complessa
		
		\item \textbf{Fusione incompleta:} A differenza delle triadi, dove le note tendono a "fondersi" in un'unica sensazione sonora, le settime mantengono una certa separazione percettiva, creando un senso di "movimento interno"
		
		\item \textbf{Tensione direzionale:} Le dissonanze di settima attivano aspettative percettive di risoluzione, guidando l'orecchio verso specifiche destinazioni armoniche
	\end{enumerate}
	
	Studi di psicoacustica hanno dimostrato che ascoltatori di culture diverse concordano nel percepire le settime come "instabili" o "in movimento", suggerendo che questa percezione ha basi in parte universali legate alla fisica del suono.
\end{fisica}

\section{Esercizi di armonizzazione completa}
\label{sec:esercizi_armonizzazione}

Per padroneggiare la versione completa della Regola dell'Ottava, è essenziale esercitarsi con bassi dati che richiedono scelte armoniche consapevoli.

\begin{esempio}[Basso dato con armonizzazione mista]
	\label{ex:basso_dato_misto}
	
	Armonizzare il seguente basso scegliendo strategicamente tra le tre versioni della Regola dell'Ottava:
	
	\musica{es_cap14_07}
	
	\textbf{Soluzione suggerita:}
	\begin{itemize}
		\item Battute 1-2 (scala ascendente e discendente): usare versione completa per massima varietà
		\item Battuta 3 (salto Do-Sol): usare stati fondamentali (I-V)
		\item Battuta 4 (cadenza finale): ritorno alla semplicità della versione preliminare
	\end{itemize}
\end{esempio}

% ============================================================
% FINE NUOVO CAPITOLO REGOLA DELL'OTTAVA COMPLETA
% ============================================================


% ============================================================
% NOTE PER L'IMPLEMENTAZIONE
% ============================================================

% Questo nuovo capitolo dovrebbe essere inserito nella struttura del manuale
% dopo il Capitolo J (Le sette specie di settima) e prima del Capitolo K 
% (Modulazione ai toni vicini).

% Nel file manuale.tex, aggiungere la riga:
% \input{Capitoli/capitoloL}  % Regola dell'Ottava completa

% Inoltre, nei capitoli esistenti che citano "la versione definitiva" o "completa"
% della regola dell'ottava, aggiornare i riferimenti con:
% \ref{chap:regola_ottava_completa}

% ============================================================
% TABELLA RIASSUNTIVA DELLA REGOLA DELL'OTTAVA COMPLETA
% Da inserire come appendice o tabella di riferimento rapido
% ============================================================

\begin{table}[h]
	\centering
	\caption{Regola dell'Ottava Completa - Tabella di riferimento}
	\label{tab:regola_ottava_completa_riepilogo}
	\small
	\begin{tabular}{|c|c|c|c|c|c|}
		\hline
		\textbf{Dir.} & \textbf{Grado} & \textbf{Basso} & \textbf{MAGGIORE} & \textbf{MAGGIORE} & \textbf{MINORE} \\
		& \textbf{Scala} & \textbf{(Es. Do/La)} & \textbf{Numerica} & \textbf{Funzione} & \textbf{Numerica} \\
		\hline
		Asc & I & Do / La & 5 & I & 5 \\
		Asc & II & Re / Si & 6 (o 4/3) & VII6 (o V4/3) & 4/3 \\
		Asc & III & Mi / Do & 6 & I6 & 6 \\
		Asc & IV & Fa / Re & 6 (o 6/5) & II6 (o II6/5) & 6 (o 6/5) \\
		Asc & V & Sol / Mi & 5 (o 7) & V (o V7) & 7 \\
		Asc & VI & La / Fa & 6 & IV6 & 6 \\
		Asc & VII & Si / \Soldiesis  & 6 (o 6/5) & V6 (o V6/5) & 6/5 \\
		Asc & VIII & Do / La & 5 & I & 5 \\
		\hline
		Disc & VIII & Do / La & 5 & I & 5 \\
		Disc & VII & Si / Sol & 6 (o 6/5) & V6 (o V6/5) & 6 \\
		Disc & VI & La / Fa & 4/3\# & II4/3\# (V del V) & \#6 \\
		Disc & V & Sol / Mi & 5 (o 7) & V (o V7) & 5 (o 7) \\
		Disc & IV & Fa / Re & 2 & V2 & 2 \\
		Disc & III & Mi / Do & 6 & I6 & 6 \\
		Disc & II & Re / Si & 4/3 & V4/3 & 4/3 \\
		Disc & I & Do / La & 5 & I & 5 \\
		\hline
	\end{tabular}
\end{table}
  % La Regola dell'Ottava: versione completa
	\chapter{Tonicizzazione e dominanti secondarie}
\label{chap:tonicizzazione}

In questo capitolo esploriamo il fenomeno della tonicizzazione, che rappresenta il primo passo verso la comprensione della modulazione.

\section{Modulazione e tonicizzazione: due fenomeni distinti}
\label{sec:mod_vs_ton}

\begin{definizione}[Modulazione]
    È il cambio \textbf{stabile} di centro tonale, confermato da una cadenza nella nuova tonalità. Rimanda a \ref{chap:modulazione}.
\end{definizione}

\begin{definizione}[Tonicizzazione]
    È il trattamento \textbf{momentaneo} di un grado non tonicale come fosse una tonica locale, senza abbandonare la tonalità principale. Il ritorno alla tonica originale avviene immediatamente dopo uno o due accordi.
\end{definizione}

\begin{regola}
    Se l'accordo tonicizzato è seguito da una cadenza nella nuova tonalità locale → modulazione. Se è seguito da un ritorno alla tonalità originale → tonicizzazione.
\end{regola}

\section{Il meccanismo: la dominante secondaria}
\label{sec:meccanismo_dominante}

\begin{definizione}[Dominante secondaria]
    È un accordo costruito e funzionante come dominante di un grado diverso dalla tonica. Si indica come V/II ("dominante del secondo"), V/V ("dominante del quinto"), ecc. È sempre una triade maggiore o una settima di dominante.
\end{definizione}

\begin{regola}
    Qualsiasi grado può essere tonicizzato tranne il VII (la triade diminuita non può essere tonica). I casi più comuni sono II, IV e V.
\end{regola}

\section{I casi più frequenti nel basso dato}
\label{sec:casi_frequenti}

\subsection{V/V (dominante della dominante)}
\begin{itemize}
    \item Struttura: accordo sul II grado alzato cromaticamente (\Fadiesis  in Do M).
    \item Risoluzione: su V o V7.
    \item Condotta: \Fadiesis  (sensibile della dominante) sale a Sol.
\end{itemize}
\musica{
    \partilate
    { \key do \major \time 4/4 }
    { sol'2 sol'2 | fad'2 sol'2 }
    { mi'2 re'2 | re'2 re'2 }
    { do'2 si2 | la2 si2 }
    { do2 sol,2 | re2 sol,2 }
    \lyricmode { "I"2 "V" | "V/V"2 "V" }
    \figuremode { s2 s | <7>2 s }
}

\subsection{V/IV (dominante della sottodominante)}
\begin{itemize}
    \item Struttura: accordo sul I grado con settima minore (\Sibemolle  in Do M).
    \item Risoluzione: su IV.
\end{itemize}
\musica{
    \partilate
    { \key do \major \time 4/4 }
    { do'2 sib | la2 r }
    { sol2 sol | fa2 r }
    { mi2 mi | do2 r }
    { do2 do | fa,2 r }
    \lyricmode { "I"2 "V/IV" | "IV"2 }
    \figuremode { s2 <7>2 | s2 }
}

\section{Riconoscimento nel basso dato}
\label{sec:riconoscimento}

\begin{regola}
    Se nel basso appare un'alterazione non prevista dalla chiave, chiedersi: "questa nota è la sensibile (terza) di una dominante secondaria?". I segnali tipici sono il \Fadiesis  sul II grado (V/V) o il \Sibemolle  sul I grado (V/IV).
\end{regola}

\section{Collegamento alla modulazione}
\label{sec:collegamento_mod}

Tonicizzazione e modulazione sono due punti di un \textit{continuum}. Una tonicizzazione prolungata e confermata da una cadenza diventa a tutti gli effetti una modulazione (\ref{chap:modulazione}).
 % Tonicizzazione e dominanti secondarie
	% ============================================================
% CAPITOLO K - LA MODULAZIONE AI TONI VICINI (VERSIONE ESTESA)
% ============================================================
\chapter{La modulazione ai toni vicini}
\label{chap:modulazione}

Nei capitoli precedenti abbiamo costruito il nostro edificio armonico all'interno di un'unica tonalità, imparando a muoverci tra le funzioni di tonica, sottodominante e dominante con fluidità e coerenza. Abbiamo classificato le sette specie di accordi di settima (Capitolo \ref{chap:sette_specie}), ognuna con la propria struttura e funzione privilegiata. Ora è il momento di imparare a viaggiare: la \textbf{modulazione} è l'arte di spostarsi da una tonalità a un'altra, aprendo nuove prospettive espressive e creando architetture musicali più ampie e complesse.

La modulazione non è un salto improvviso nel vuoto, ma un processo organico e logico, che sfrutta le relazioni naturali tra le tonalità e la capacità di certi accordi di avere un doppio significato armonico. In questo capitolo ci concentreremo sulle modulazioni più naturali e immediate: quelle verso i \textbf{toni vicini}, ovvero le tonalità che condividono la maggior parte delle note con la tonalità di partenza.

\section{Che cos'è una modulazione}
\label{sec:definizione_modulazione}

Prima di addentrarci nelle tecniche, è necessario definire con precisione cosa intendiamo per ``modulazione'' e distinguerla da altri fenomeni armonici.

\begin{definizione}[Modulazione]
	\label{def:modulazione}
	Una \textbf{modulazione} è il processo mediante il quale una composizione musicale abbandona stabilmente una tonalità di partenza (tonalità di impianto) per stabilirne una nuova. La modulazione comporta:
	\begin{enumerate}
		\item L'affermazione chiara della tonalità di partenza
		\item Un percorso armonico di transizione
		\item La conferma della nuova tonalità attraverso una cadenza
	\end{enumerate}
	
	Si distingue dal semplice \textbf{accenno} o \textbf{tonicizzazione}, in cui si crea momentaneamente l'illusione di una nuova tonalità senza abbandonare stabilmente quella di partenza.
\end{definizione}

La differenza tra modulazione vera e propria e tonicizzazione è sottile ma importante: nella tonicizzazione, la nuova tonalità viene ``visitata'' brevemente ma non confermata da una cadenza autentica; nella modulazione, invece, il cambiamento è strutturale e la nuova tonalità diventa il nuovo centro di gravità armonico.

\section{Le tre fasi del processo modulante}
\label{sec:fasi_modulazione}

Ogni modulazione ben costruita si articola in tre momenti distinti, ciascuno con una funzione specifica.

\begin{regola}[Le tre fasi della modulazione]
	\label{reg:tre_fasi_modulazione}
	Il processo modulante si sviluppa attraverso:
	
	\begin{enumerate}
		\item \textbf{Affermazione (o consolidamento):} La tonalità di partenza viene stabilita chiaramente attraverso progressioni caratteristiche e, preferibilmente, una cadenza. L'ascoltatore deve percepire senza ambiguità qual è il ``centro tonale'' da cui si parte.
		
		\item \textbf{Transizione (o modulazione propriamente detta):} Si introduce un \textbf{accordo pivot} (o accordo cardine), ovvero un accordo che appartiene simultaneamente a entrambe le tonalità. Questo accordo funge da ``ponte'' armonico, permettendo il passaggio dalla tonalità vecchia a quella nuova senza discontinuità percettive.
		
		\item \textbf{Conferma (o cadenza):} La nuova tonalità viene confermata attraverso una progressione cadenzale chiara, tipicamente una cadenza perfetta (V--I) o composta (II--V--I o IV--V--I). Questa fase è essenziale: senza una cadenza convincente nella nuova tonalità, abbiamo solo una tonicizzazione passeggera, non una modulazione stabile.
	\end{enumerate}
\end{regola}

Questo schema tripartito non è una regola rigida ma un modello ideale. Nelle composizioni reali, specialmente in quelle più brevi, la fase di affermazione può essere implicita (se la tonalità di partenza è già ben stabilita) e la transizione può essere molto rapida. Ciò che non può mai mancare è la \textbf{conferma} della nuova tonalità.

\section{I toni vicini: una geografia armonica}
\label{sec:toni_vicini}

Non tutte le tonalità sono ugualmente ``lontane'' l'una dall'altra. Alcune condividono molte note, altre poche o nessuna. Questa ``distanza'' armonica è visualizzabile attraverso il circolo delle quinte (cfr. Figura \ref{fig:circolo_quinte_migliorato_compatto}).

\begin{definizione}[Tonalità vicine]
	\label{def:tonalita_vicine}
	Due tonalità si dicono \textbf{vicine} (o adiacenti) se le loro armature di chiave differiscono per non più di \textbf{un'alterazione} (un diesis o un bemolle). Questa vicinanza si traduce nella condivisione di sei note su sette della scala diatonica.
	
	Dal punto di vista del circolo delle quinte, le tonalità vicine sono quelle immediatamente adiacenti (distanza di una quinta) o legate dal rapporto di relativa maggiore/minore.
\end{definizione}

\subsection{I cinque toni vicini di una tonalità maggiore}
\label{subsec:toni_vicini_maggiore}

\begin{regola}[Tonalità vicine al modo maggiore]
	\label{reg:toni_vicini_maggiore}
	Data una tonalità maggiore, i suoi cinque toni vicini sono:
	
	\begin{enumerate}
		\item La \textbf{relativa minore} (stessa armatura di chiave)
		\item La tonalità della \textbf{dominante} (una quinta sopra, un diesis in più o un bemolle in meno)
		\item La \textbf{relativa minore della dominante} (stessa armatura della dominante)
		\item La tonalità della \textbf{sottodominante} (una quinta sotto, un bemolle in più o un diesis in meno)
		\item La \textbf{relativa minore della sottodominante} (stessa armatura della sottodominante)
	\end{enumerate}
\end{regola}

\begin{esempio}[I cinque toni vicini di Do Maggiore]
	\label{ex:toni_vicini_do}
	Partendo da Do Maggiore (nessuna alterazione), i toni vicini sono:
	
	\begin{itemize}
		\item \textbf{La minore} (relativa, 0 alterazioni)
		\item \textbf{Sol Maggiore} (dominante, 1\#)
		\item \textbf{Mi minore} (relativa di Sol, 1\#)
		\item \textbf{Fa Maggiore} (sottodominante, 1$\flat$)
		\item \textbf{Re minore} (relativa di Fa, 1$\flat$)
	\end{itemize}
	
	Tutte queste tonalità condividono sei note su sette con Do Maggiore.
\end{esempio}

\subsection{I cinque toni vicini di una tonalità minore}
\label{subsec:toni_vicini_minore}

\begin{regola}[Tonalità vicine al modo minore]
	\label{reg:toni_vicini_minore}
	Data una tonalità minore, i suoi cinque toni vicini sono:
	
	\begin{enumerate}
		\item La \textbf{relativa maggiore} (stessa armatura di chiave della minore naturale)
		\item La tonalità del \textbf{III grado} (relativa maggiore, ma costruita sul III della scala minore naturale)
		\item La tonalità del \textbf{IV grado} (sottodominante)
		\item La tonalità del \textbf{V grado} (dominante)
		\item La tonalità del \textbf{VI grado} (sopradominante)
	\end{enumerate}
	
	Nota: Nel modo minore, la situazione è leggermente più complessa a causa delle varianti armoniche e melodiche della scala, ma il principio della vicinanza per alterazioni rimane valido.
\end{regola}

\begin{esempio}[I cinque toni vicini di La minore]
	\label{ex:toni_vicini_la_m}
	Partendo da La minore (nessuna alterazione in chiave), i toni vicini sono:
	
	\begin{itemize}
		\item \textbf{Do Maggiore} (relativa, 0 alterazioni)
		\item \textbf{Mi minore} (dominante, 1\#)
		\item \textbf{Re minore} (sottodominante, 1$\flat$)
		\item \textbf{Sol Maggiore} (VII della minore = sottodominante della relativa, 1\#)
		\item \textbf{Fa Maggiore} (VI della minore = sottodominante della relativa, 1$\flat$)
	\end{itemize}
\end{esempio}

\section{L'accordo pivot: il ponte tra due mondi}
\label{sec:accordo_pivot}

Il cuore della tecnica modulante ai toni vicini è l'uso intelligente degli \textbf{accordi pivot} (detti anche accordi cardine o di perno).

\begin{definizione}[Accordo pivot]
	\label{def:accordo_pivot}
	Un \textbf{accordo pivot} è un accordo che appartiene simultaneamente a due tonalità diverse, avendo una funzione armonica specifica in ciascuna di esse. Nella modulazione, l'accordo pivot viene dapprima percepito come parte della tonalità di partenza, poi ``reinterpretato'' come parte della tonalità di arrivo.
	
	La doppia identità dell'accordo pivot permette una transizione fluida e naturale, senza discontinuità percettive.
\end{definizione}

Non tutti gli accordi sono ugualmente efficaci come pivot. La scelta dipende da due fattori:

\begin{enumerate}
	\item \textbf{Condivisione delle note:} L'accordo deve essere costruibile con le note di entrambe le scale
	\item \textbf{Forza funzionale nella nuova tonalità:} Un accordo pivot che nella nuova tonalità ha funzione di II, IV o V grado è più efficace di uno che ha funzione di III o VI
\end{enumerate}

\subsection{Accordi pivot tra tonalità relative}
\label{subsec:pivot_relative}

La modulazione tra tonalità relative (Do Maggiore $\leftrightarrow$ La minore) è la più semplice, poiché le due tonalità condividono tutte le note della scala diatonica.

\begin{regola}[Accordi pivot tra relative]
	\label{reg:pivot_relative}
	Tra Do Maggiore e La minore, gli accordi pivot più efficaci sono:
	
	\begin{itemize}
		\item \textbf{Re minore:} II grado in Do M, IV grado in La m (funzione di sottodominante in entrambe)
		\item \textbf{Mi minore:} III grado in Do M, V grado in La m (da tonica a dominante)
		\item \textbf{Fa Maggiore:} IV grado in Do M, VI grado in La m
	\end{itemize}
	
	Il più efficace è tipicamente il II di Do M / IV di La m, perché ha funzione di sottodominante forte in entrambe le tonalità.
\end{regola}

\begin{esempio}[Modulazione Do Maggiore $\rightarrow$ La minore tramite IV di La m]
	\label{ex:mod_do_lam}
	Progressione che modula usando Re minore come pivot.

	\musica{
		\partilate
		{ \key do \major \time 4/4 }
		{ mi'2 fa'2 | sold2 la2 | la1 }
		{ do'2 do'2 | mi2 mi2 | mi1 }
		{ sol2 la2 | si,2 do2 | do1 }
		{ do2 fa,2 | mi,2 la,2 | la,1 }
		\lyricmode { 
			\markup{"Do M:"} "I"2 "IV (pivot)"2 | 
			\markup{"La m:"} "V"2 "I"2 | 
			"I"1 
		}
		\figuremode { s2 s2 | <\markup{\sharp}>2 s2 | s1 }
	}
	Il Fa maggiore (IV di Do M) viene reinterpretato come IV di La minore, preparando la cadenza V--I in La minore.
\end{esempio}

\subsection{Accordi pivot verso la dominante}
\label{subsec:pivot_dominante}

La modulazione dalla tonica alla sua dominante (Do M $\rightarrow$ Sol M) è altrettanto naturale.

\begin{regola}[Accordi pivot verso la dominante maggiore]
	\label{reg:pivot_dominante}
	Tra Do Maggiore e Sol Maggiore, gli accordi pivot più comuni sono:
	
	\begin{itemize}
		\item \textbf{Do Maggiore stesso:} I grado in Do M, IV grado in Sol M
		\item \textbf{Mi minore:} III grado in Do M, VI grado in Sol M
		\item \textbf{Sol Maggiore (con settima):} V$^7$ in Do M, I$^7$ in Sol M (meno comune come pivot stabile)
	\end{itemize}
	
	Il più efficace è Do Maggiore reinterpretato come IV di Sol, perché permette di precedere immediatamente la cadenza V--I nella nuova tonalità.
\end{regola}

\begin{esempio}[Modulazione Do Maggiore $\rightarrow$ Sol Maggiore tramite pivot I = IV]
	\label{ex:mod_do_sol}

	\musica{
		\partilate
		{ \key do \major \time 4/4 }
		{ mi'2 mi'2 | fad'2 sol'2 | sol'1 }
		{ do'2 do'2 | do'2 si2 | si1 }
		{ sol2 sol2 | la2 sol2 | sol1 }
		{ do2 do2 | re2 sol,2 | sol,1 }
		\lyricmode { 
			\markup{"Do M:"} "I (pivot)"2 "I"2 | 
			\markup{"Sol M:"} "IV"2 "V"2 | 
			"I"1 
		}
		\figuremode { s2 s2 | s2 s2 | s1 }
	}
	L'accordo di Do Maggiore, percepito inizialmente come tonica, viene reinterpretato come sottodominante di Sol Maggiore.
\end{esempio}

\subsection{Accordi pivot verso la sottodominante}
\label{subsec:pivot_sottodominante}

\begin{regola}[Accordi pivot verso la sottodominante maggiore]
	\label{reg:pivot_sottodominante}
	Tra Do Maggiore e Fa Maggiore:
	
	\begin{itemize}
		\item \textbf{Fa Maggiore stesso:} IV in Do M, I in Fa M
		\item \textbf{Re minore:} II in Do M, VI in Fa M
		\item \textbf{Do Maggiore:} I in Do M, V in Fa M
	\end{itemize}
	
	Il più efficace è Do Maggiore reinterpretato come dominante di Fa, che porta direttamente alla cadenza.
\end{regola}

\begin{esempio}[Modulazione Do Maggiore $\rightarrow$ Fa Maggiore tramite pivot I = V]
	\label{ex:mod_do_fa}
	
	\musica{
		\partilate
		{ \key do \major \time 4/4 }
		{ mi'2 mi'2 | fa'2 mi'2 | fa'1 }
		{ do'2 sib2 | la2 do'2 | do'1 }
		{ sol2 sol2 | do2 sib2 | la1 }
		{ do2 do2 | fa,2 fa,2 | fa,1 }
		\lyricmode { 
			\markup{"Do M:"} "I"2 "I (pivot)"2 | 
			\markup{"Fa M:"} "V"2 "V"2 | 
			"I"1 
		}
		\figuremode { s2 s2 | s2 s2 | s1 }
	}
\end{esempio}

\section{Le sette specie come strumenti modulanti}
\label{sec:specie_modulazione}

La conoscenza delle sette specie di accordi di settima (Capitolo \ref{chap:sette_specie}) apre possibilità modulanti molto più ricche e sofisticate rispetto all'uso delle sole triadi.

\subsection{La seconda specie: il pivot per eccellenza}
\label{subsec:seconda_specie_pivot}

Come abbiamo visto nel Capitolo \ref{chap:vii7_ii7}, la settima semidiminuita (seconda specie) ha una natura doppia: può essere VII$\oslash^7$ nel modo maggiore o II$\oslash^7$ nel modo minore.

\begin{regola}[La seconda specie come accordo pivot]
	\label{reg:seconda_specie_pivot}
	L'accordo di settima semidiminuita è il pivot ideale tra tonalità relative:
	
	\begin{itemize}
		\item \textbf{VII$\oslash^7$ di Do Maggiore = II$\oslash^7$ di La minore}
		\item Stesso accordo (Si--Re--Fa--La), due funzioni: dominante in Do M, sottodominante in La m
	\end{itemize}
	
	Questa ambiguità funzionale permette modulazioni estremamente fluide tra tonalità relative.
\end{regola}

\begin{esempio}[Modulazione con VII$\oslash^7$ = II$\oslash^7$ come pivot]
	\label{ex:mod_seconda_specie}

	\musica{
		\partilate
		{ \key do \major \time 4/4 }
		{ mi'2 mi'2 | la'2 sold'2 | la'1 }
		{ do'2 re'2 | fa'2 mi'2 | mi'1 }
		{ sol2 si2 | do'2 si2 | do'1 }
		{ do2 si,2 | re2 mi2 | la,1 }
		\lyricmode { 
			\markup{"Do M:"} "I"2 "VII$\oslash^7$ (pivot)"2 | 
			\markup{"La m:"} "II$\oslash^7$"2 "V"2 | 
			"I"1 
		}
		\figuremode { s2 <6 5>2 | <6 5>2 s2 | s1 }
	}
	L'accordo Si--Re--Fa--La funge da cerniera: viene percepito come VII$\oslash^7$ (funzione di dominante) in Do Maggiore, poi reinterpretato come II$\oslash^7$ (funzione di sottodominante) in La minore.
\end{esempio}

\subsection{La terza specie: versatilità modulante}
\label{subsec:terza_specie_pivot}

La settima minore (terza specie) appare su molti gradi di diverse scale, rendendola estremamente versatile.

\begin{regola}[La terza specie come pivot tra tonalità maggiori]
	\label{reg:terza_specie_pivot}
	Esempi di accordi di terza specie usabili come pivot:
	
	\begin{itemize}
		\item \textbf{Re--Fa--La--Do:} II$^7$ in Do M = VI$^7$ in Fa M = VII$^7$ in Mi m
		\item \textbf{Mi--Sol--Si--Re:} III$^7$ in Do M = VI$^7$ in Sol M = II$^7$ in Re M
	\end{itemize}
	
	La terza specie permette di modulare non solo ai toni vicini immediati, ma anche di concatenare più modulazioni successive.
\end{regola}

\begin{esempio}[Modulazione Do M $\rightarrow$ Fa M tramite II$^7$ = VI$^7$]
	\label{ex:mod_terza_specie}
	
	\musica{
		\partilate
		{ \key do \major \time 4/4 }
		{ mi'2 do'2 | sib2 la2 | la1 }
		{ do'2 la2 | fa2 fa2 | fa1 }
		{ sol2 re2 | re2 do2 | do1 }
		{ do2 re,2 | sol,2 fa,2 | fa,1 }
		\lyricmode { 
			\markup{"Do M:"} "I"2 "II$^7$ (pivot)"2 | 
			\markup{"Fa M:"} "VI$^7$"2 "V"2 | 
			"I"1 
		}
		\figuremode { s2 <7>2 | <7>2 s2 | s1 }
	}
\end{esempio}

\subsection{La quinta specie: modulazioni audaci}
\label{subsec:quinta_specie_modulazione}

Come discusso nella Sezione \ref{sec:quinta_specie_dettaglio}, la settima diminuita ha una struttura simmetrica che le permette di risolvere su quattro toniche diverse.

\begin{regola}[La quinta specie per modulazioni multiple]
	\label{reg:quinta_specie_modulazioni}
	Un singolo accordo di settima diminuita può essere reinterpretato come sensibile di quattro tonalità diverse (distanti una terza minore l'una dall'altra). Questo apre la possibilità di modulazioni:
	
	\begin{itemize}
		\item Ai toni vicini (quando la nuova tonica dista una terza minore)
		\item Ai toni lontani (argomento che svilupperemo in un capitolo successivo)
	\end{itemize}
	
	Attenzione: L'uso della quinta specie per modulazioni ai toni lontani esce dall'ambito di questo capitolo e verrà trattato separatamente.
\end{regola}

Per ora, ci limitiamo a osservare che la quinta specie è uno strumento potentissimo, ma richiede grande attenzione nella preparazione e nella conferma della nuova tonalità per non creare ambiguità percettive.

\section{La tecnica della modulazione: consigli pratici}
\label{sec:tecnica_modulazione}

\subsection{Scegliere il pivot giusto}
\label{subsec:scelta_pivot}

\begin{regola}[Criteri per la scelta del pivot]
	\label{reg:scelta_pivot}
	Un buon accordo pivot deve soddisfare questi requisiti:
	
	\begin{enumerate}
		\item \textbf{Esistere in entrambe le tonalità} con le stesse note (o con alterazioni compatibili)
		\item \textbf{Avere una funzione forte nella nuova tonalità:} preferire II, IV o V grado della nuova tonalità piuttosto che III o VI
		\item \textbf{Permettere una condotta delle voci fluida} dalla tonalità vecchia a quella nuova
		\item \textbf{Non creare ambiguità premature:} il pivot deve essere chiaramente percepito come parte della tonalità vecchia finché non viene reinterpretato
	\end{enumerate}
\end{regola}

\subsection{Confermare la nuova tonalità}
\label{subsec:conferma_tonalita}

\begin{regola}[Importanza della cadenza di conferma]
	\label{reg:conferma_cadenza}
	Una modulazione non è completa senza una cadenza convincente nella nuova tonalità. Le opzioni più efficaci sono:
	
	\begin{itemize}
		\item \textbf{Cadenza perfetta V--I:} La più forte e inequivocabile
		\item \textbf{Cadenza composta II--V--I o IV--V--I:} Rafforza ulteriormente la nuova tonalità
		\item \textbf{Ripetizione della cadenza:} In caso di dubbio, ripetere la progressione cadenzale elimina ogni ambiguità
	\end{itemize}
	
	Evitare di interrompere la nuova tonalità subito dopo il pivot, senza dare tempo alla nuova tonica di stabilizzarsi.
\end{regola}

\subsection{Gestire le alterazioni}
\label{subsec:gestione_alterazioni}

\begin{regola}[Introduzione delle alterazioni]
	\label{reg:introduzione_alterazioni}
	Quando si modula a una tonalità con un diesis o un bemolle in più:
	
	\begin{itemize}
		\item \textbf{Introdurre l'alterazione gradualmente:} La nuova alterazione appare tipicamente nella sensibile dell'accordo di dominante della nuova tonalità
		\item \textbf{Evitare salti cromatici bruschi:} Se possibile, far apparire l'alterazione in una voce interna piuttosto che nel soprano
		\item \textbf{Preparare l'orecchio:} L'accordo pivot dovrebbe già contenere alcune delle note ``nuove'' della scala di arrivo
	\end{itemize}
\end{regola}

\begin{esempio}[Gestione del \Fadiesis  nella modulazione Do M $\rightarrow$ Sol M]
	\label{ex:gestione_fad}

	\musica{
		\partilate
		{ \key do \major \time 4/4 }
		{ do'2 do'2 | fad'2 sol'2 | sol'1 }
		{ sol2 sol2 | do'2 si2 | si1 }
		{ mi2 mi2 | la2 sol2 | sol1 }
		{ do2 do2 | re2 sol,2 | sol,1 }
		\lyricmode { 
			\markup{"Do M:"} "I (pivot)"2 "I"2 | 
			\markup{"Sol M:"} "IV"2 "V"2 | 
			"I"1 
		}
		\figuremode { s2 s2 | s2 s2 | s1 }
	}
	Il \Fadiesis  (sensibile di Sol Maggiore) appare solo nell'accordo di dominante (Re--\Fadiesis --La--Do), rendendo la modulazione fluida e naturale.
\end{esempio}

\section{Esempi completi di modulazione}
\label{sec:esempi_completi}

Presentiamo ora alcuni esempi completi che illustrano le tre fasi del processo modulante in diversi contesti.

\begin{esempio}[Modulazione completa: La minore $\rightarrow$ Do Maggiore]
	\label{ex:mod_completa_lam_dom}
	Affermazione in La m, pivot attraverso Fa M (VI di La m = IV di Do M), conferma in Do M.

	\musica{
		\partilate
		{ \key la \minor \time 4/4 }
		{ la'2 do''2 | do''2 re''2 | mi''2 do''2 | do''1 }
		{ mi'2 mi'2 | fa'2 fa'2 | do'2 si2 | do'1 }
		{ do'2 do'2 | la2 la2 | sol2 sol2 | mi1 }
		{ la,2 la,2 | fa,2 fa,2 | do2 sol,2 | do1 }
		\lyricmode { 
			\markup{"La m:"} "I"2 "III"2 | 
			"VI (pivot)"2 "VI"2 | 
			\markup{"Do M:"} "IV"2 "V"2 | 
			"I"1 
		}
		\figuremode { s2 s2 | s2 s2 | s2 s2 | s1 }
	}
	La tonalità di partenza (La minore) viene affermata con la progressione I--III. L'accordo di Fa Maggiore funge da pivot, inizialmente percepito come VI di La m, poi reinterpretato come IV di Do M. La cadenza IV--V--I conferma solidamente Do Maggiore.
\end{esempio}

\begin{esempio}[Modulazione con settima: Do M $\rightarrow$ Mi minore]
	\label{ex:mod_settima_do_mim}
	Uso del II$\oslash^7$ di Mi minore come pivot.

	\musica{
		\partilate
		{ \key do \major \time 4/4 }
		{ do'2 do'2 | do'2 si2 | do'2 si2 | si1 }
		{ sol2 la2 | la2 fad2 | sold2 fad2 | mi1 }
		{ mi2 fad2 | re2 re2 | mi2 re2 | si,1 }
		{ do2 re,2 | fad,2 si,,2 | mi,2 si,,2 | mi,,1 }
		\lyricmode { 
			\markup{"Do M:"} "I"2 "II"2 | 
			\markup{"Mi m:"} "II$\oslash^7$"2 "V"2 | 
			"I"2 "V"2 | 
			"I"1 
		}
		\figuremode { s2 s2 | <6 5>2 s2 | s2 s2 | s1 }
	}
\end{esempio}

\section{Regole pratiche per la modulazione rapida}
\label{sec:regole_pratiche_modulazione}

Esistono alcuni procedimenti cromatici ricorrenti che permettono di modulare rapidamente verso le tonalità vicine partendo da una tonalità maggiore.

\begin{regola}[Modulazione al IV grado tramite la sensibile abbassata]
	\label{reg:mod_iv_sensibile_abbassata}
	In una tonalità maggiore, abbassando di un semitono la \textbf{sensibile} (VII grado), si ottiene una nota che funge da \textbf{settima} dell'accordo costruito sulla tonica. Tale accordo (\textbf{I7}) assume funzione di dominante della sottodominante, permettendo una modulazione immediata al \textbf{IV grado}.
	
	\textit{Esempio in Do Maggiore: il Si viene abbassato a \Sibemolle , trasformando l'accordo di Do maggiore in Do7, che risolve su Fa maggiore.}
\end{regola}

\begin{regola}[Modulazione al V grado tramite il IV grado innalzato]
	\label{reg:mod_v_quarto_innalzato}
	Per modulare al \textbf{V grado} (dominante), si agisce sul \textbf{IV grado} della tonalità di partenza innalzandolo di un semitono. Questa nota diventa la \textbf{nuova sensibile} del V grado, trasformando l'accordo sul II grado (o sul IV) in un accordo di dominante della dominante (\textbf{V/V}).
	
	\textit{Esempio in Do Maggiore: il Fa viene innalzato a \Fadiesis , diventando la sensibile di Sol maggiore e preparando l'arrivo della nuova tonica.}
\end{regola}

\begin{regola}[Modulazione al VI grado tramite il V grado innalzato]
	\label{reg:mod_vi_quinto_innalzato}
	Per modulare al \textbf{VI grado} (relativo minore della tonica), si innalza di un semitono il \textbf{V grado} della tonalità di partenza. Questa nota funge da \textbf{nuova sensibile} del VI grado, creando un accordo di dominante (\textbf{V/VI}) che risolve sulla nuova tonica minore.
	
	\textit{Esempio in Do Maggiore: il Sol viene innalzato a \Soldiesis , diventando la sensibile di La minore.}
\end{regola}

\begin{regola}[Modulazione ai relativi minori degli altri gradi]
	\label{reg:mod_altri_minori}
	Analogamente, è possibile modulare ai relativi minori degli altri gradi principali alterando i gradi della scala per ottenere la nuova sensibile:
	
	\begin{itemize}
		\item \textbf{Al III grado (relativo del V):} Si innalzano il \textbf{II} e il \textbf{IV grado} (es. in Do M: \Fadiesis  e \Rediesis  per modulare a Mi minore).
		\item \textbf{Al II grado (relativo del IV):} Si innalza il \textbf{I grado} (es. in Do M: \Dodiesis  per modulare a Re minore).
	\end{itemize}
\end{regola}  % Modulazione ai toni vicini
	\chapter{Progressioni armoniche per l'armonizzazione del basso}
\label{chap:progressioni_armoniche}

Dopo aver studiato le funzioni armoniche, le cadenze e la Regola dell'Ottava, siamo pronti ad affrontare uno dei compiti più creativi e allo stesso tempo più tecnici dell'armonia: l'\textbf{armonizzazione di un basso dato}. Questo capitolo è dedicato alle \textbf{progressioni armoniche}, ovvero successioni di accordi che si ripetono trasposte, costituendo uno strumento essenziale per armonizzare efficacemente passaggi scalari, sequenze modulanti e movimenti melodici complessi.

\section{Che cos'è una progressione armonica}
\label{sec:definizione_progressione}

\begin{definizione}[Progressione armonica]
	\label{def:progressione_armonica}
	Una \textbf{progressione armonica} (o \textbf{sequenza armonica}) è una successione di accordi che viene ripetuta trasportata su diversi gradi della scala. Questa ripetizione crea un movimento armonico coerente e prevedibile che può essere utilizzato per costruire frasi musicali più lunghe o per effettuare modulazioni.
	
	Elementi caratteristici:
	\begin{itemize}
		\item \textbf{Schema di base (modello):} Una breve successione di accordi (tipicamente 2-4 accordi)
		\item \textbf{Trasposizione:} Il modello viene ripetuto su gradi diversi della scala
		\item \textbf{Direzione:} La progressione può muoversi verso l'acuto (ascendente) o verso il grave (discendente)
		\item \textbf{Numero di ripetizioni:} Tipicamente 2-4 volte, raramente di più per evitare monotonia
	\end{itemize}
\end{definizione}

Le progressioni armoniche sono particolarmente utili nell'armonizzazione del basso per diversi motivi:

\begin{itemize}
	\item Forniscono una soluzione sistematica per passaggi scalari lunghi
	\item Creano movimento e dinamismo armonico
	\item Facilitano la modulazione verso tonalità vicine
	\item Rendono il discorso musicale coerente e memorabile
\end{itemize}

\begin{fisica}[Perché le progressioni sono così efficaci?]
	Dal punto di vista percettivo, le progressioni sfruttano uno dei principi fondamentali della cognizione musicale: il \textbf{riconoscimento di pattern}. 
	
	Il cervello umano è straordinariamente abile nell'identificare schemi ricorrenti. Quando ascoltiamo una progressione, dopo la prima o seconda ripetizione del modello, il nostro sistema cognitivo:
	\begin{enumerate}
		\item \textbf{Anticipa} il movimento successivo
		\item \textbf{Categorizza} il materiale come "coerente"
		\item \textbf{Elabora} più facilmente la struttura complessiva
	\end{enumerate}
	
	Questo fenomeno, studiato dalla psicologia della Gestalt, spiega perché le progressioni risultano così naturali all'ascolto: non solo creano varietà armonica, ma lo fanno in modo che il nostro cervello possa seguire agevolmente.
	
	Interessante notare che studi di neuroimaging hanno dimostrato che, quando una progressione viene interrotta inaspettatamente, si attiva una risposta cerebrale specifica (chiamata N400) associata alla violazione delle aspettative, confermando che il cervello costruisce attivamente modelli predittivi durante l'ascolto musicale.
\end{fisica}

\section{Progressioni fondamentali}
\label{sec:progressioni_fondamentali}

Esistono diverse famiglie di progressioni, classificabili in base alla direzione del movimento armonico. Le più comuni sono le progressioni per \textbf{quinte discendenti} (o quarte ascendenti) e le progressioni per \textbf{terze} (ascendenti o discendenti).

\subsection{Progressione per quinte discendenti}
\label{subsec:progressione_quinte}

La progressione per quinte discendenti è una delle più potenti e naturali nel sistema tonale, poiché replica il movimento cadenzale V-I su ogni grado della scala.

\begin{regola}[Progressione per quinte discendenti]
	\label{reg:progressione_quinte_disc}
	Nella \textbf{progressione per quinte discendenti}, ogni accordo viene seguito da un altro la cui fondamentale è una quinta sotto (o una quarta sopra). La successione tipica è:
	
	\textbf{Nel modo maggiore:}
	I - IV - VII - III - VI - II - V - I
	
	\textbf{Nel modo minore:}
	i - iv - VII - III - VI - ii° - V - i
	
	Caratteristiche:
	\begin{itemize}
		\item Il movimento del basso procede per \textbf{quarte ascendenti} (equivalente intervallare della quinta discendente)
		\item Gli accordi alternano generalmente stati fondamentali e primi rivolti per mantenere il movimento del basso fluido
		\item La sensazione armonica è di continua "risoluzione" e direzione verso la tonica
	\end{itemize}
\end{regola}

\begin{esempio}[Progressione per quinte discendenti in Do Maggiore]
	\label{ex:progressione_quinte_disc}
	Progressione completa con alternanza di stati fondamentali e rivolti.
	\musica{
		\partilate
		{ \key do \major \time 4/4 }
		{ mi''2 fa'' | re'' mi'' | do'' re'' | do'' si' | do''1 }
		{ do''2 do'' | si' do'' | la' si' | sol' sol' | mi'1 }
		{ sol'2 la' | sol' sol' | fa' fa' | mi' re' | sol1 }
		{ do2 fa | sol do' | la re | mi sol | do1 }
		\lyricmode { "I"2 "IV" | "VII" "III" | "VI" "II" | "III" "V" | "I"1 }
		\figuremode { \unset FiguredBass.useBassFigureExtenders 
			s2 s | s <6> | s s | <6> s | s1 }
	}
	
	Notare come l'alternanza tra stato fondamentale (5) e primo rivolto (6) mantiene il basso in movimento per grado congiunto o salti piccoli, evitando salti eccessivi.
\end{esempio}

\begin{matematica}[Il cerchio delle quinte e la progressione]
	La progressione per quinte discendenti non è altro che una "passeggiata" lungo il \textbf{cerchio delle quinte}. Questo cerchio può essere rappresentato come un grafo ciclico dove ogni nodo rappresenta una tonalità e gli archi connettono tonalità adiacenti per quinte.
	
	Se rappresentiamo le 7 note diatoniche di Do Maggiore sul cerchio delle quinte, otteniamo:
	
	\begin{center}
	\begin{tikzpicture}[scale=1.2]
		\foreach \note/\angle in {Fa/0, Do/51.43, Sol/102.86, Re/154.29, La/205.71, Mi/257.14, Si/308.57} {
			\node[circle, draw, minimum size=0.8cm] (\note) at (\angle:2cm) {\note};
		}
		\draw[->, thick] (Do) to[bend left=15] (Fa);
		\draw[->, thick] (Fa) to[bend left=15] (Si);
		\draw[->, thick] (Si) to[bend left=15] (Mi);
		\draw[->, thick] (Mi) to[bend left=15] (La);
		\draw[->, thick] (La) to[bend left=15] (Re);
		\draw[->, thick] (Re) to[bend left=15] (Sol);
		\draw[->, thick] (Sol) to[bend left=15] (Do);
	\end{tikzpicture}
	\end{center}
	
	La progressione per quinte discendenti percorre questo cerchio in senso antiorario. Dal punto di vista della teoria dei grafi, questo è un \textbf{cammino hamiltoniano} sul sottografo ciclico delle note diatoniche.
	
	Interessante notare che in temperamento equabile, dopo 12 quinte giuste si torna alla nota di partenza (a meno di una piccola discrepanza chiamata \textit{comma pitagorico}), ma nel sistema diatonico, dopo 7 passi si ritorna alla tonica originale.
\end{matematica}

\subsection{Progressione per terze discendenti}
\label{subsec:progressione_terze_disc}

La progressione per terze discendenti crea un movimento armonico più morbido e cromatico rispetto a quella per quinte.

\begin{regola}[Progressione per terze discendenti]
	\label{reg:progressione_terze_disc}
	Nella \textbf{progressione per terze discendenti}, ogni accordo scende di terza (maggiore o minore) rispetto al precedente. Una forma comune è:
	
	\textbf{Nel modo maggiore:}
	I - VI - IV - II - VII - V - I
	
	Caratteristiche:
	\begin{itemize}
		\item Il movimento è più graduale e "dolce" rispetto alla progressione per quinte
		\item Spesso utilizzata in passaggi espressivi o melanconici
		\item Il basso può muoversi per terze o, più frequentemente, per grado congiunto usando rivolti appropriati
	\end{itemize}
\end{regola}

\begin{esempio}[Progressione per terze discendenti in Do Maggiore]
	\label{ex:progressione_terze_disc}
	Sequenza che scende per terze con movimento del basso per grado congiunto.
	\musica{
		\partilate
		{ \key do \major \time 3/4 }
		{ mi''2. | do'' | la' | fa' | re' | si' | do''2. }
		{ do''2. | la' | fa' | re' | si' | sol' | mi'2. }
		{ sol'2. | mi' | do' | la' | sol' | re' | sol'2. }
		{ do'2. | la | fa | re | sol | sol | do2. }
		\lyricmode { "I"2. | "VI" | "IV" | "II" | "VII" | "V" | "I"2. }
		\figuremode { \unset FiguredBass.useBassFigureExtenders 
			s2. | <6> | s | <6> | <6> | s | s2. }
	}
	Questa progressione scende cromaticamente attraverso la scala, creando una sensazione di dolce malinconia.
\end{esempio}

\section{Applicazione pratica: armonizzare un basso dato}
\label{sec:applicazione_basso_dato}

L'abilità di riconoscere quando e come applicare le progressioni è essenziale per l'armonizzazione del basso. Vediamo un metodo sistematico.

\begin{regola}[Strategia per l'armonizzazione del basso]
	\label{reg:strategia-armonizzazione}
	Quando ci troviamo di fronte a un basso da armonizzare, seguire questi passi:
	
	\begin{enumerate}
		\item \textbf{Identificare la struttura:}
		\begin{itemize}
			\item Cadenze (dove si trovano i punti di riposo?)
			\item Passaggi scalari (ascendenti o discendenti)
			\item Salti significativi (quarte, quinte, ottave)
		\end{itemize}
		
		\item \textbf{Scegliere la tecnica appropriata:}
		\begin{itemize}
			\item Per \textbf{movimenti scalari brevi} (2-4 note): applicare la Regola dell'Ottava
			\item Per \textbf{movimenti scalari lunghi} (5+ note): considerare progressioni
			\item Per \textbf{salti}: usare prevalentemente stati fondamentali
		\end{itemize}
		
		\item \textbf{Applicare le progressioni:}
		\begin{itemize}
			\item Identificare il pattern che il basso suggerisce (quinte, terze, seconde?)
			\item Applicare la progressione appropriata
			\item Verificare che la condotta delle voci sia corretta
		\end{itemize}
		
		\item \textbf{Raffinare:}
		\begin{itemize}
			\item Correggere parallelismi
			\item Migliorare il canto (linea del soprano)
			\item Aggiungere varietà ritmica se necessario
		\end{itemize}
	\end{enumerate}
\end{regola}

\begin{esempio}[Basso dato con progressione per quinte]
	\label{ex:basso_dato_progressione}
	
	\textbf{Basso da armonizzare:}
	\musica{
		\new Staff {
			\clef bass
			\key do \major
			\time 4/4
			do4 fa sol do' | la re sol do
		}
	}
	
	\textbf{Analisi:}
	\begin{itemize}
		\item Battuta 1: salto Do-Fa (quarta) suggerisce I-IV
		\item Fa-Sol-Do': parte di una progressione per quinte (IV-VII-III o IV-V-I)
		\item Battuta 2: La-Re-Sol-Do continua la discesa per quinte
	\end{itemize}
	
	\textbf{Soluzione armonizzata:}
	\musica{
		\partilate
		{ \key do \major \time 4/4 }
		{ mi'4 fa' sol' mi' | do' re' re' do' }
		{ do'4 do' re' do' | la si si mi' }
		{ sol4 la si sol | mi sol sol sol }
		{ do4 fa sol do' | la re sol do }
		\lyricmode { "I"4 "IV" "V" "III" | "VI" "II" "V" "I" }
		\figuremode { \unset FiguredBass.useBassFigureExtenders 
			s4 s s <6> | s s s s }
	}
	
	La progressione per quinte (IV-V-III-VI-II-V-I) permette di creare un movimento armonico fluido e convincente.
\end{esempio}

\section{Progressioni con settime}
\label{sec:progressioni_settime}

L'uso delle settime nelle progressioni aumenta considerevolmente la ricchezza armonica e la tensione direzionale.

\begin{regola}[Progressioni di settime]
	\label{reg:progressioni_settime}
	Nelle progressioni, le settime possono essere impiegate in due modi:
	
	\begin{enumerate}
		\item \textbf{Settime alternate:} Solo alcuni accordi della progressione hanno la settima
		\item \textbf{Settime consecutive:} Ogni accordo della progressione è un accordo di settima
	\end{enumerate}
	
	Nel caso di settime consecutive, la condotta delle voci segue una regola precisa:
	\begin{itemize}
		\item La settima di ciascun accordo \textbf{scende di grado} nell'accordo successivo
		\item Questa nota diventa spesso la \textbf{terza} dell'accordo successivo
		\item Il pattern crea una "catena" di risoluzioni che lega gli accordi
	\end{itemize}
\end{regola}

\begin{esempio}[Progressione per quinte con settime consecutive]
	\label{ex:progressione_settime}
	Ogni accordo è una settima che risolve nel successivo.
	\musica{
		\partilate
		{ \key do \major \time 4/4 }
		{ fa''2 mi'' | re'' do'' | si'1 }
		{ re''2 do'' | si' la' | sol'1 }
		{ la'2 sol' | sol' fa' | re'1 }
		{ re2 do | sol, la, | sol,1 }
		\lyricmode { "II7"2 "V7" | "I7" "IV7" | "VII°7"1 }
		\figuremode { \unset FiguredBass.useBassFigureExtenders 
			<7>2 <7> | <7> <7> | <\markup{\flat}7>1 }
	}
	
	Notare come ogni settima (Fa nel II7, Fa nel V7, Si nel I7, Mi nel IV7) scende di grado per diventare la terza dell'accordo successivo.
\end{esempio}

\begin{errore}[Progressioni troppo lunghe]
	\label{err:progressioni_lunghe}
	Un errore comune è quello di applicare meccanicamente una progressione per troppe ripetizioni, creando monotonia. 
	
	\textbf{Problema:}
	\musica{
		\new Staff {
			\clef bass
			\key do \major
			\time 4/4
			do4 fa si, mi | la, re sol, do | fa, si,, mi, la,, | re,1
		}
	}
	
	Questa progressione per quinte discendenti si estende per tre battute complete, risultando meccanica e prevedibile.
	
	\textbf{Soluzione:}
	Interrompere la progressione dopo 2-3 cicli e introdurre varietà:
	\begin{itemize}
		\item Usare una cadenza per chiudere la frase
		\item Cambiare tipo di progressione
		\item Inserire un accordo di passaggio o nota pedale
	\end{itemize}
	
	In generale, raramente si dovrebbero superare le 3-4 ripetizioni consecutive dello stesso modello.
\end{errore}

\section{Progressioni modulanti}
\label{sec:progressioni_modulanti}

Una delle applicazioni più affascinanti delle progressioni è la loro capacità di facilitare la modulazione verso tonalità vicine.

\begin{definizione}[Progressione modulante]
	\label{def:progressione_modulante}
	Una \textbf{progressione modulante} è una sequenza armonica che, attraverso le sue ripetizioni trasposte, conduce naturalmente da una tonalità a un'altra. La modulazione avviene spesso in modo graduale e quasi impercettibile, sfruttando la coerenza del pattern ripetitivo.
\end{definizione}

\begin{regola}[Modulazione attraverso progressioni]
	\label{reg:modulazione_progressioni}
	Per modulare usando una progressione:
	
	\begin{enumerate}
		\item Iniziare la progressione nella tonalità di partenza
		\item Continuare il pattern introducendo gradualmente alterazioni che appartengono alla nuova tonalità
		\item Confermare la nuova tonalità con una cadenza autentica
	\end{enumerate}
	
	Le progressioni per quinte discendenti sono particolarmente efficaci perché:
	\begin{itemize}
		\item Ogni "quinta" può potenzialmente diventare una dominante della tonalità successiva
		\item L'introduzione di alterazioni cromati (es. \# per creare sensibili) avviene naturalmente
		\item Il movimento è direzionale e porta organicamente verso la nuova tonica
	\end{itemize}
\end{regola}

\begin{esempio}[Progressione modulante: da Do Maggiore a Sol Maggiore]
	\label{ex:progressione_modulante}
	Una progressione per quinte che modula al tono della dominante.
	\musica{
		\partilate
		{ \key do \major \time 4/4 }
		{ mi''4 fa'' fad'' sol'' | sol'' fad'' sol''2 }
		{ do''4 do'' re'' si' | re'' re'' si'2 }
		{ sol'4 la' la' re'' | si' la' re'2 }
		{ do4 fa si, mi | la, re sol,2 }
		\lyricmode { "I (Do)"4 "IV" "V/V" "V" | "II (Sol)" "V" "I"2 }
		\figuremode { \unset FiguredBass.useBassFigureExtenders 
			s4 s <6 \markup{\sharp}> s | s s s2 }
	}
	
	Il \Fadiesis  introdotto nel secondo accordo trasforma il VII° di Do in V di Sol, preparando la modulazione. La conferma avviene con la cadenza II-V-I in Sol Maggiore.
\end{esempio}

\section{Esercizi di armonizzazione}
\label{sec:esercizi_progressioni}

Per padroneggiare le progressioni, è essenziale esercitarsi con bassi dati di complessità crescente.

\begin{esempio}[Esercizio 1: Basso con progressione per quinte]
	\label{ex:esercizio_1}
	
	\textbf{Basso da armonizzare:}
	\musica{
		\new Staff {
			\clef bass
			\key sol \major
			\time 3/4
			sol,4 do4 si, | la, re sol, | do2.
		}
	}
	
	\textbf{Suggerimenti:}
	\begin{itemize}
		\item Identificare la progressione per quinte: Sol-Do (I-IV), Do-Si (IV-III), Si-La (III-II), La-Re-Sol (II-V-I)
		\item Usare rivolti per mantenere il movimento fluido del basso
		\item Fare attenzione alla condotta della sensibile (\Fadiesis )
	\end{itemize}
\end{esempio}

\begin{esempio}[Esercizio 2: Basso con Monte Romanesca]
	\label{ex:esercizio_2}
	
	\textbf{Basso da armonizzare:}
	\musica{
		\new Staff {
			\clef bass
			\key fa \major
			\time 4/4
			fa,4 sol, la, sib, | do re mib do | fa,1
		}
	}
	
	\textbf{Suggerimenti:}
	\begin{itemize}
		\item Riconoscere il Monte (salita) Fa-Sol-La-\Sibemolle -Do
		\item Il \Mibemolle  introduce una breve colorazione cromatica
		\item Applicare l'alternanza 6-5 caratteristica del Monte
		\item Concludere con cadenza plagale o autentica
	\end{itemize}
\end{esempio}

\begin{esempio}[Esercizio 3: Basso misto (Regola dell'Ottava + Progressione)]
	\label{ex:esercizio_3}
	
	\textbf{Basso da armonizzare:}
	\musica{
		\new Staff {
			\clef bass
			\key re \minor
			\time 4/4
			re4 mi fa sol | la sib do' la | re1
		}
	}
	
	\textbf{Suggerimenti:}
	\begin{itemize}
		\item Battuta 1 (Re-Mi-Fa-Sol): applicare Regola dell'Ottava ascendente
		\item Battuta 2 (La-\Sibemolle -Do-La): breve progressione (IV-V-VI-IV) o cadenza d'inganno
		\item Combinare le due tecniche per creare un'armonizzazione coerente
	\end{itemize}
\end{esempio}

\section{Progressioni nella letteratura musicale}
\label{sec:progressioni_letteratura}

Le progressioni armoniche non sono semplici esercizi scolastici, ma costituiscono una delle tecniche compositive più utilizzate nella storia della musica.

\begin{esempio}[Progressioni in Bach]
	\label{ex:bach_progressioni}
	Johann Sebastian Bach utilizza magistralmente le progressioni in molte opere:
	
	\begin{itemize}
		\item \textbf{Preludio in Do Maggiore BWV 846} (Clavicembalo ben temperato I): Progressione per quinte discendenti che attraversa l'intera gamma tonale
		
		\item \textbf{Ciaccona dalla Partita II per violino solo BWV 1004}: Progressione per terze discendenti ripetuta con variazioni
		
		\item \textbf{Fuga in \Dodiesis  minore BWV 849}: Monte Romanesca utilizzato per costruire tensione nel contrappunto
	\end{itemize}
	
	In Bach, le progressioni non sono mai meccaniche ma sono integrate in un discorso contrappuntistico ricco, dove ogni voce mantiene interesse melodico indipendente.
\end{esempio}

\begin{fisica}[La "fisica" delle progressioni: energia armonica]
	Possiamo pensare alle progressioni armoniche in termini di \textbf{energia potenziale armonica}. Ogni accordo in una tonalità possiede un certo livello di "tensione" rispetto alla tonica:
	
	\begin{itemize}
		\item \textbf{Tonica (I):} Energia minima (riposo)
		\item \textbf{Dominante (V):} Energia massima (tensione)
		\item \textbf{Altri gradi:} Livelli intermedi
	\end{itemize}
	
	Le progressioni possono essere viste come "percorsi" attraverso questo paesaggio energetico:
	
	\begin{enumerate}
		\item \textbf{Progressione per quinte discendenti}: Simile a una "cascata" energetica dove ogni accordo "cade" verso il successivo (come dominante locale)
		
		\item \textbf{Monte Romanesca}: Come una "salita" graduale che accumula tensione fino a un punto di rilascio
		
		\item \textbf{Progressione per terze}: Movimento più "orizzontale", mantenendo livelli energetici simili
	\end{enumerate}
	
	Questa metafora non è puramente astratta: studi sulla percezione armonica hanno dimostrato che gli ascoltatori percepiscono effettivamente diverse successioni di accordi come aventi diversi gradi di "stabilità" o "tensione", e che queste percezioni sono sorprendentemente coerenti tra ascoltatori di culture diverse.
\end{fisica}

\section{Progressioni con salti di quinta e di quarta}
\label{sec:progressioni_salti_quinta}

Oltre alle progressioni scalari che abbiamo finora esaminato, esistono progressioni particolarmente efficaci che si basano su \textbf{salti di quinta} (o quarta, che è l'inversione intervallare della quinta) nel basso. Queste progressioni creano un movimento armonico robusto e diretto, ideale per passaggi virtuosistici o per stabilire con forza il centro tonale.

\subsection{Progressione di quinte con numera 5}
\label{subsec:progressione_quinte_5}

La progressione più semplice e diretta con salti di quinta prevede l'uso di \textbf{accordi allo stato fondamentale} (numerati tutti 5, o senza numerazione esplicita). Il basso si muove per salti di quinta discendente (o quarta ascendente), e ogni accordo è presentato nella sua forma più stabile.

\begin{regola}[Progressione per quinte discendenti (tutti 5)]
	\label{reg:progressione_5}
	Nella progressione per salti di quinta con accordi allo stato fondamentale:
	
	\begin{itemize}
		\item Il \textbf{basso} si muove per salti di quinta discendente (o quarta ascendente)
		\item Tutti gli accordi sono allo \textbf{stato fondamentale} (cifra 5 o non numerati)
		\item Schema tipico: I - IV - VII - III - VI - II - V - I
		\item La condotta delle voci richiede particolare attenzione per evitare quinte e ottave parallele
	\end{itemize}
	
	\textbf{Tecnica di condotta delle voci:}
	Nel collegamento tra due accordi allo stato fondamentale con basso che salta di quinta, una tecnica efficace è muovere le tre voci superiori in \textbf{moto contrario} rispetto al basso, oppure mantenere una o due note comuni quando possibile.
\end{regola}

\begin{esempio}[Progressione per quinte discendenti (tutti 5) in Sol Maggiore]
	\label{ex:progressione_5_sol}
	Una sequenza di accordi tutti allo stato fondamentale.
	\musica{
		\partilate
		{ \key sol \major \time 4/4 }
		{ si'1 | la' | sol' | fad' | sol'1 }
		{ sol'1 | fad' | mi' | re' | si1 }
		{ re'1 | re' | si | la | sol1 }
		{ sol1 | do' | mi | la, | sol,1 }
		\lyricmode { "I"1 | "IV" | "VII" | "III" | "I"1 }
		\figuremode { s1 | s | <6> | s | s1 }
	}
	
	In questo esempio, il basso salta per quinte discendenti (Sol-Do-Mi-La-Sol), creando una forte progressione cadenzale. Notare che per evitare problemi di condotta si può usare occasionalmente un rivolto (come il VII6 nell'esempio).
\end{esempio}

\subsection{Progressione di quinte con numera 7 (catena di settime)}
\label{subsec:progressione_quinte_7}

Una variante più sofisticata della progressione per quinte prevede l'uso sistematico di \textbf{accordi di settima}, tutti allo stato fondamentale. Questa tecnica, chiamata anche \textbf{catena di settime}, crea un movimento armonico particolarmente fluido e ricco.

\begin{regola}[Progressione per quinte con accordi di settima (tutti 7)]
	\label{reg:progressione_7}
	Nella \textbf{catena di settime}:
	
	\begin{itemize}
		\item Ogni accordo è una \textbf{settima} allo stato fondamentale (numerato 7)
		\item Il basso si muove per salti di quinta discendente (o quarta ascendente)
		\item La settima di ogni accordo \textbf{deve essere preparata} (apparire come consonanza nell'accordo precedente)
		\item La settima \textbf{risolve scendendo di grado} nell'accordo successivo
		\item Questa tecnica crea un movimento quasi automatico delle voci
	\end{itemize}
	
	\textbf{Meccanismo della risoluzione:}
	La settima di un accordo diventa automaticamente la terza dell'accordo successivo (che è una quinta sotto), creando una catena di preparazioni e risoluzioni perfettamente interconnesse.
\end{regola}

\begin{esempio}[Catena di settime in Re minore]
	\label{ex:catena_settime}
	Una progressione con accordi di settima che si susseguono.
	\musica{
		\partilate
		{ \key re \minor \time 4/4 }
		{ fa'2 mi' | re' do' | si' la' | sol'1 }
		{ re'2 do' | si' la' | sol' fa' | mi'1 }
		{ la2 sol | fa mi | re do | si1 }
		{ re2 la | si, fa | sol re | mi1 }
		\lyricmode { "I7"2 "IV7" | "VII7" "III7" | "VI7" "II7" | "V"1 }
		\figuremode { <7>2 <7> | <7> <7> | <7> <7> | s1 }
	}
	\begin{itemize}
		\item Il Do del soprano (settima di I7) scende al Si (terza di IV7)
		\item Il La del contralto (settima di IV7) scende al Sol (terza di VII7)
		\item Questo schema si ripete in tutta la progressione
		\item L'ultima settima (II7) risolve sulla dominante pura (V senza settima)
	\end{itemize}
\end{esempio}

\begin{matematica}[La struttura ciclica delle progressioni per quinte]
	Le progressioni per quinte discendenti hanno una bellissima struttura matematica legata al \textbf{gruppo ciclico} $\mathbb{Z}_7$ (i sette gradi della scala diatonica).
	
	Se rappresentiamo i gradi come numeri modulo 7:
	\[
	\text{Gradi: } I=0, II=1, III=2, IV=3, V=4, VI=5, VII=6
	\]
	
	Un salto di quinta discendente corrisponde a \textbf{sottrarre 3 modulo 7} (oppure aggiungere 4, che è equivalente). La progressione completa può essere scritta come:
	
	\[
	0 \xrightarrow{-3} 4 \xrightarrow{-3} 1 \xrightarrow{-3} 5 \xrightarrow{-3} 2 \xrightarrow{-3} 6 \xrightarrow{-3} 3 \xrightarrow{-3} 0
	\]
	
	Ovvero (traducendo in gradi):
	\[
	I \rightarrow V \rightarrow II \rightarrow VI \rightarrow III \rightarrow VII \rightarrow IV \rightarrow I
	\]
	
	Questa è esattamente la progressione per quinte discendenti! Notare che:
	\begin{itemize}
		\item Dopo 7 passi si torna alla tonica (il gruppo è ciclico di ordine 7)
		\item Ogni grado viene visitato esattamente una volta (la funzione $x \mapsto x-3 \mod 7$ è una permutazione)
		\item L'operazione $-3$ è l'inversa di $+3$ (salto di quarta ascendente), quindi le due progressioni sono simmetriche
	\end{itemize}
	
	\textbf{Proprietà del gruppo:} Poiché 3 e 7 sono coprimi ($\gcd(3,7) = 1$), l'operazione $x \mapsto x-3 \mod 7$ genera l'intero gruppo $\mathbb{Z}_7$. Questo garantisce che la progressione percorra tutti i gradi della scala prima di tornare alla tonica.
	
	Se invece usassimo un intervallo di terza ($-2 \mod 7$), otterremmo cicli più corti, che non percorrono tutti i gradi. Ad esempio:
	\[
	0 \xrightarrow{-2} 5 \xrightarrow{-2} 3 \xrightarrow{-2} 1 \xrightarrow{-2} 6 \xrightarrow{-2} 4 \xrightarrow{-2} 2 \xrightarrow{-2} 0
	\]
	
	Ancora una volta, la musica rivela la sua profonda struttura algebrica!
\end{matematica}

\section{Progressioni cromatiche}
\label{sec:progressioni_cromatiche}

Le progressioni cromatiche introducono note alterate che non appartengono alla scala diatonica di partenza. Queste progressioni sono particolarmente efficaci per creare movimento ascendente o discendente per gradi congiunti cromatici, e sono state ampiamente utilizzate dai compositori barocchi per costruire tensione o per modulare verso tonalità lontane.

\subsection{Progressione cromatica ascendente (5, 6/5)}
\label{subsec:progressione_cromatica_ascendente}

La progressione cromatica ascendente è caratterizzata da un basso che sale per semitoni, armonizzato alternando accordi allo stato fondamentale con accordi di settima al primo rivolto.

\begin{regola}[Progressione cromatica ascendente (numera 5, 6/5)]
	\label{reg:progressione_cromatica_asc}
	Nella \textbf{progressione cromatica ascendente}:
	
	\begin{itemize}
		\item Il basso sale per \textbf{semitoni} (movimento cromatico)
		\item Gli accordi alternano \textbf{triadi in stato fondamentale} (cifra 5) e \textbf{settime al primo rivolto} (cifra 6/5)
		\item Schema tipico: 5 - 6/5 - 5 - 6/5 - 5...
		\item Le alterazioni cromatiche creano sensibili "artificiali" che tendono a salire
		\item Particolarmente efficace per costruire tensione verso un climax
	\end{itemize}
	
	\textbf{Meccanismo armonico:}
	Ogni accordo di settima (6/5) contiene una dissonanza che tende a risolvere sull'accordo successivo. Le alterazioni cromatiche (di solito \#) creano delle "piccole dominanti" che spingono il discorso armonico verso l'alto.
\end{regola}

\begin{esempio}[Progressione cromatica ascendente in Do minore]
	\label{ex:progressione_cromatica_asc}
	Una salita cromatica dal VI al I grado.
	\musica{
		\partilate
		{ \key do \minor \time 4/4 }
		{ mib'4 mib' mi' mi' | fa' fa' fad' fad' | sol'1 }
		{ do'4 sib lab' do' | do' re' la' do' | si'1 }
		{ lab4 sol fa' mi' | mib' la si' la' | re'1 }
		{ lab,4 sol, lab, do | la, re si, re | sol,1 }
		\lyricmode { "VI"4 "V7/VI" "VI" "V7" | "I" "V7/II" "II" "V7/V" | "V"1 }
		\figuremode { s4 <6 5> s <6 5> | s <6 5> s <6 5> | s1 }
	}
	
	In questo esempio:
	\begin{itemize}
		\item Il basso sale cromaticamente: \Labemolle  - Sol - \Labemolle  - Do - La - Re - Si - Re - Sol
		\item L'alterazione del basso (La$\natural$ invece di \Labemolle ) crea una dominante secondaria di Re (II grado)
		\item Il Si$\natural$ crea una dominante della dominante (V/V)
		\item L'intera progressione spinge energicamente verso la dominante finale (Sol)
	\end{itemize}
\end{esempio}

\begin{fisica}[L'energia cinetica delle progressioni ascendenti]
	Dal punto di vista percettivo, le progressioni cromatiche ascendenti hanno un effetto psicoacustico molto specifico: creano una sensazione di \textbf{accumulo di energia}.
	
	Questo fenomeno può essere compreso attraverso due principi:
	
	\begin{enumerate}
		\item \textbf{Il pitch ascendente e l'eccitazione}: Studi di psicoacustica hanno dimostrato che l'ascolto di altezze crescenti attiva nel cervello meccanismi simili a quelli della risposta allo stress o all'eccitazione. Il movimento verso l'acuto è percepito come "energizzante" e "tensivo".
		
		\item \textbf{La densità delle risoluzioni}: In una progressione cromatica, la frequenza delle risoluzioni dissonanza-consonanza è molto alta (praticamente ad ogni accordo). Ogni micro-risoluzione crea un piccolo rilascio di tensione, ma poiché subito dopo si presenta una nuova dissonanza, l'effetto netto è di \textbf{tensione crescente}.
	\end{enumerate}
	
	Il risultato è simile a una salita in montagna: ogni passo richiede energia, e anche se ci sono piccoli momenti di "riposo", la fatica complessiva aumenta fino al raggiungimento della vetta (il climax armonico).
	
	I compositori hanno utilizzato questo effetto per secoli. Ad esempio, nelle fughe di Bach, le progressioni cromatiche ascendenti sono spesso usate negli \textit{stretti} finali per creare la massima tensione prima della risoluzione conclusiva.
\end{fisica}

\subsection{Progressione cromatica discendente (2 7)}
\label{subsec:progressione_cromatica_discendente}

La progressione cromatica discendente è in un certo senso speculare alla precedente: il basso scende per semitoni, armonizzato con una particolare sequenza di accordi.

\begin{regola}[Progressione cromatica discendente (numera 2 7)]
	\label{reg:progressione_cromatica_disc}
	Nella \textbf{progressione cromatica discendente}:
	
	\begin{itemize}
		\item Il basso scende per \textbf{semitoni} (movimento cromatico discendente)
		\item Gli accordi alternano \textbf{settime al terzo rivolto} (cifra 2) e \textbf{settime allo stato fondamentale} (cifra 7)
		\item Schema tipico: 2 - 7 - 2 - 7 - 2...
		\item Le alterazioni cromati (spesso $\flat$) facilitano la discesa cromatica
		\item Crea un effetto di discesa controllata, spesso con carattere malinconico o patetico
	\end{itemize}
	
	\textbf{Meccanismo armonico:}
	Il terzo rivolto (2) ha il basso che è la settima dell'accordo. Questa settima deve risolvere scendendo di grado, creando naturalmente un movimento discendente. La successione 2-7 permette di mantenere questo movimento cromatico nel basso mentre le altre voci si muovono più liberamente.
\end{regola}

\begin{esempio}[Progressione cromatica discendente in La minore]
	\label{ex:progressione_cromatica_disc}
	Una discesa cromatica con alternanza 2-7.
	\musica{
		\partilate
		{ \key la \minor \time 3/4 }
		{ la'2. | sol' | fa' | mi' | re' | do' | si' | la'2. }
		{ fa'2. | mi' | re' | do' | si' | la' | sol' | mi'2. }
		{ re'2. | do' | si' | la | sold | mi | re | do'2. }
		{ re2. | do | si, | la, | sold, | sol, | fa, | mi2. }
		\lyricmode { "II7"2. | "I7" | "VII7" | "VI7" | "V7" | "IV7" | "III7" | "I"2. }
		\figuremode { <2>2. | <7> | <2> | <7> | <2> | <7> | <2> | s2. }
	}
	In questo esempio:
	\begin{itemize}
		\item Il basso scende cromaticamente da Re a Mi (passando per \Dodiesis , Do$\natural$, Si, La, \Soldiesis , Sol$\natural$, \Fadiesis )
		\item L'alternanza tra terzo rivolto (2) e stato fondamentale (7) mantiene la continuità del movimento
		\item Le settime risolvono regolarmente scendendo di grado
		\item L'effetto complessivo è di una discesa controllata e espressiva
	\end{itemize}
\end{esempio}

\begin{esempio}[Uso della progressione 2-7 in Bach]
	\label{ex:bach_27}
	J.S. Bach utilizza frequentemente la progressione cromatica discendente (2-7) in passaggi particolarmente espressivi:
	
	\begin{itemize}
		\item \textbf{Clavicembalo ben temperato I, Fuga n. 4 in \Dodiesis  minore}: Una memorabile discesa cromatica con alternanza 2-7 al culmine dello sviluppo
		
		\item \textbf{Passione secondo Matteo, Aria "Erbarme dich"}: L'accompagnamento del violino obbligato include discese cromatiche armonizzate con questa tecnica, sottolineando il testo patetico
		
		\item \textbf{Messa in Si minore, Crucifixus}: L'intero movimento è costruito su un basso ostinato cromatico discendente, simboleggiando la discesa di Cristo dalla croce
	\end{itemize}
	
	In tutti questi casi, la progressione 2-7 non è solo una tecnica compositiva, ma un mezzo espressivo per comunicare dolore, compassione o tragedia. Il movimento discendente cromatico è stato tradizionalmente associato nella retorica musicale barocca al \textit{passus duriusculus} (passo difficile), una figura retorica che rappresenta sofferenza.
\end{esempio}

\begin{matematica}[Cromatismo e temperamento equabile]
	Le progressioni cromatiche sollevano una questione interessante legata al \textbf{temperamento equabile}.
	
	Nel sistema diatonico puro, i semitoni non sono tutti uguali:
	\begin{itemize}
		\item \textbf{Semitono cromatico}: Es. Do-\Dodiesis  (rapporto $\approx 25:24 \approx 1.042$)
		\item \textbf{Semitono diatonico}: Es. Mi-Fa (rapporto $\approx 16:15 \approx 1.067$)
	\end{itemize}
	
	In temperamento equabile, invece, tutti i semitoni sono identici:
	\[
	\text{Rapporto di semitono} = 2^{1/12} \approx 1.059
	\]
	
	Questo valore è una media geometrica che rende possibile la modulazione libera e le progressioni cromatiche che percorrono tutte le dodici note della scala cromatica.
	
	Possiamo vedere la progressione cromatica come una \textbf{passeggiata sul gruppo ciclico} $\mathbb{Z}_{12}$:
	
	Se partiamo da Do (0) e saliamo cromaticamente:
	\[
	0 \rightarrow 1 \rightarrow 2 \rightarrow 3 \rightarrow 4 \rightarrow 5 \rightarrow 6 \rightarrow 7 \rightarrow 8 \rightarrow 9 \rightarrow 10 \rightarrow 11 \rightarrow 0
	\]
	
	Dopo 12 semitoni torniamo esattamente all'ottava superiore della nota di partenza. Questa proprietà è garantita solo dal temperamento equabile.
	
	Nel sistema di intonazione naturale (basato su rapporti semplici di frequenze), una sequenza di 12 semitoni non tornerebbe esattamente all'ottava ma avrebbe una discrepanza chiamata \textbf{comma pitagorico}:
	\[
	\text{Discrepanza} = \left(\frac{3}{2}\right)^{12} \div 2^7 \approx 1.013
	\]
	
	Il temperamento equabile "distribuisce" questo errore uniformemente su tutti i semitoni, permettendo la libertà modulatoria e cromatica che caratterizza la musica dal Barocco in poi.
\end{matematica}

\section{Sintesi: quando usare quale progressione}
\label{sec:sintesi_progressioni}

\begin{regola}[Guida alla scelta delle progressioni]
	\label{reg:scelta_progressioni}
	
	\begin{tabular}{|p{4.5cm}|p{4.5cm}|p{4.5cm}|}
		\hline
		\textbf{Tipo di basso} & \textbf{Progressione consigliata} & \textbf{Effetto espressivo} \\
		\hline
		Salita scalare lunga (5+ note) & Progressione scalare ascendente per gradi (con alternanza 5, 6/5 secondo la Regola dell'Ottava, \ref{chap:regola_ottava_completa}) \footnotemark & Costruzione di tensione progressiva verso la cadenza \\
		\hline

		Discesa scalare & Progressione per terze o Regola dell'Ottava & Dolcezza, malinconia \\
		\hline
		Movimento circolare & Progressione per quinte discendenti & Forza, direzione, conclusività \\
		\hline
		Passaggio modulante & Progressione per quinte con alterazioni & Transizione fluida \\
		\hline
		Basso statico (note pedale) & Variazione armonica sopra nota fissa & Stabilità con varietà \\
		\hline
		Salti di quarta/quinta & Stati fondamentali (5) o Catena di settime (7) & Chiarezza strutturale, robustezza \\
		\hline
		Salita cromatica & Progressione 5, 6/5 & Accumulo di tensione, virtuosismo \\
		\hline
		Discesa cromatica & Progressione 2, 7 & Pathos, espressività, discesa controllata \\
		\hline
	\end{tabular}
\end{regola}
\footnotetext{Alcune fonti citano la "Romanesca" come schema per armonizzare salite scalari. La Romanesca è in realtà una progressione di basso ostinato specifica del repertorio rinascimentale e barocco — non una tecnica generica di armonizzazione scalare. La sua trattazione esula dall'ambito di questo manuale introduttivo.}

\vspace{1cm}

Le progressioni armoniche rappresentano uno strumento compositivo straordinariamente versatile. La loro padronanza richiede pratica costante: armonizzare bassi dati quotidianamente, analizzare partiture di grandi compositori per identificare le progressioni utilizzate, e sperimentare creative combinazioni di pattern diversi. Con il tempo, l'uso delle progressioni diventerà intuitivo, permettendo di affrontare con sicurezza anche i bassi più complessi.

% ============================================================
% FINE CAPITOLO M - PROGRESSIONI ARMONICHE
% ============================================================
  % Progressioni armoniche
	% ============================================================
% CAPITOLO - METODO DI LAVORO PER IL BASSO NUMERATO
% (Versione espansa di 99_basso_dato.tex)
% ============================================================
\chapter{Come affrontare il basso numerato}
\label{chap:metodo_lavoro}

L'armonizzazione di un basso numerato è la prova fondamentale di ogni esame di armonia. In questo capitolo presentiamo un \textbf{metodo di lavoro sistematico}, una procedura in fasi sequenziali pensata per trasformare un esercizio potenzialmente disorientante in un percorso logico e controllato.

Il basso numerato (o \textbf{basso figurato}) è un basso con le cifre che indicano gli intervalli da costruire sopra di esso. Rispetto al basso senza numeri (Capitolo \ref{chap:dal_numerato_al_senza_numeri}), la numerica ci fornisce già la ``numerica implicita'': il nostro compito è tradurre quei numeri in accordi reali, rispettando la condotta delle voci.

\section{Lettura della numerica: convenzioni fondamentali}
\label{sec:lettura_numerica}

\begin{definizione}[La numerica del basso figurato]
	\label{def:numerica}
	La \textbf{numerica} (o \textbf{cifratura}) è un sistema di numeri scritti sotto le note del basso che indicano gli intervalli delle voci superiori rispetto al basso stesso. Le convenzioni principali sono:
	\begin{itemize}
		\item \textbf{Nessun numero} (o \textbf{5}): accordo in stato fondamentale (quinta e terza sopra il basso).
		\item \textbf{6} (o 6/3): primo rivolto (sesta e terza sopra il basso).
		\item \textbf{6/4}: secondo rivolto (sesta e quarta sopra il basso) --- usare solo nei contesti ammessi (cfr. Capitolo \ref{chap:quarta_sesta}).
		\item \textbf{7}: accordo di settima allo stato fondamentale.
		\item \textbf{6/5}: primo rivolto di settima (sesta, quinta e terza sopra il basso).
		\item \textbf{4/3}: secondo rivolto di settima.
		\item \textbf{2} (o 4/2): terzo rivolto di settima (la settima è al basso).
		\item Un \textbf{numero preceduto da \#, \flat{} o $\natural$}: l'alterazione si applica all'intervallo indicato.
		\item Un \textbf{\# o + da solo} sotto una nota: alzare di un semitono la terza dell'accordo (comune in Do minore per alzare la sensibile).
	\end{itemize}
\end{definizione}

\begin{osservazione}[La numerica è relativa al basso, non alla scala]
	\label{oss:numerica_relativa}
	Gli intervalli della numerica sono \textit{relativi al basso reale}, non al grado della scala. Se il basso è Re e la numerica è 6, si deve costruire una sesta sopra Re (Si), indipendentemente dalla tonalità. Le alterazioni esplicite nella numerica hanno sempre la precedenza sulle note ``naturali'' della scala.
\end{osservazione}

\section{Fase 1 — Analisi preliminare}
\label{sec:fase1_numerato}

Anche con la numerica disponibile, la prima operazione non è scrivere accordi, ma \textbf{comprendere il terreno}.

\begin{itemize}
	\item \textbf{Determinare la tonalità:} Osservare l'armatura di chiave e verificare la nota iniziale e finale. Prestare attenzione alle alterazioni cromatiche presenti nel basso: spesso indicano la sensibile del modo minore (\ref{chap:tonicizzazione}) o modulazioni (\ref{chap:modulazione}).
	\item \textbf{Individuare il modo:} Cercare la presenza della sensibile (VII grado alzato nel minore) e le note tipiche del minore armonico.
	\item \textbf{Leggere la numerica:} I numeri sono \textit{suggerimenti} preziosi. Indicano la struttura degli accordi desiderata dal compositore/esaminatore, ma vanno sempre interpretati nel contesto della tonalità e della condotta delle voci.
	\item \textbf{Riconoscere le note non armoniche:} Identificare eventuali note di passaggio, note di volta o bordoni nel basso stesso (\ref{chap:noteestranee}). Questi movimenti non portano accordi propri.
\end{itemize}

\section{Fase 2 — Segmentazione fraseologica}
\label{sec:fase2_numerato}

Dividere il basso in unità di senso compiuto permette di pianificare i punti di riposo.

\begin{itemize}
	\item \textbf{Individuare le cadenze:} Cercare i movimenti tipici (V-I, IV-I, formule cadenzali) che segnano la fine di una frase o di un periodo. Spesso la numerica stessa rivela la cadenza: una successione di numeri \texttt{<4> <3>} indica un ritardo 4-3 cadenzale (\ref{chap:ritardi}); \texttt{<7>} sulla dominante indica la settima di dominante in cadenza.
	\item \textbf{Struttura antecedente-conseguente:} Verificare se il basso segue la tipica struttura di 8 misure (\ref{chap:struttura_frase}). Pianificare in anticipo la cadenza di metà (aperta) e quella finale (perfetta).
\end{itemize}

\begin{regola}[Pianificazione cadenzale nel basso numerato]
	\label{reg:pianificazione_cadenzale_numerato}
	\begin{itemize}
		\item \textbf{Metà (mis. 4 di 8):} Deve presentare una cadenza sospesa (V) o imperfetta (I non fondamentale o soprano non in ottava).
		\item \textbf{Finale (mis. 8):} Deve presentare una cadenza perfetta (V-I in stato fondamentale, soprano all'ottava o alla quinta).
		\item Se la numerica non specifica risoluzione della dominante: usare la forma completa (V7 completa) o incompleta secondo il contesto (\ref{def:v7_completa_incompleta}).
	\end{itemize}
\end{regola}

\section{Fase 3 — Classificazione di ogni nota di basso}
\label{sec:fase3_numerato}

Stabilire la funzione e lo stato di ogni accordo, \textit{leggendo la numerica e interpretandola nel contesto}.

\begin{itemize}
	\item \textbf{Assegnazione degli accordi:} Decidere quale accordo assegnare a ogni nota in base alla funzione armonica (\ref{chap:funzioni_armoniche}) e alla numerica disponibile. Verificare che l'accordo indicato dalla numerica sia coerente con la tonalità.
	\item \textbf{Tonicizzazioni e cromatismi:} Identificare se un'alterazione nella numerica indica una dominante secondaria (\ref{chap:tonicizzazione}) o una modulazione (\ref{chap:modulazione}). Un \# sulla terza dell'accordo spesso indica la sensibile di una tonalità locale.
	\item \textbf{Scelta dello stato:} La numerica specifica lo stato, ma va verificata la coerenza con:
		\begin{enumerate}
			\item La posizione metrica (stati fondamentali sui tempi forti di norma).
			\item L'uso corretto del secondo rivolto (\ref{chap:quarta_sesta}).
			\item Le regole di risoluzione delle dissonanze (\ref{chap:teoria_dissonanza}).
		\end{enumerate}
\end{itemize}

\section{Fase 4 — Costruzione delle voci superiori}
\label{sec:fase4_numerato}

Una volta stabilita l'ossatura armonica, si procede alla scrittura SATB.

\begin{regola}[Priorità nella scrittura delle voci superiori]
	\label{reg:priorita_scrittura_numerato}
	\begin{enumerate}
		\item \textbf{Vettori obbligati:} Risolvere sempre correttamente la sensibile (sale alla tonica) e la settima (scende di grado). Questi movimenti hanno \textit{priorità assoluta} sugli altri.
		\item \textbf{Note comuni:} Tenere le note comuni tra accordi successivi nella stessa voce.
		\item \textbf{Grado congiunto:} Muovere le altre voci preferibilmente per grado congiunto.
		\item \textbf{Moto contrario:} Preferire il moto contrario tra basso e soprano per evitare quinte e ottave nascoste.
	\end{enumerate}
\end{regola}

\begin{itemize}
	\item \textbf{Scelta dei raddoppi:} Raddoppiare di norma la fondamentale. \textbf{Mai} raddoppiare la sensibile o la settima.
	\item \textbf{Posizione al soprano:} Scegliere la posizione (ottava, terza, quinta) cercando di creare una linea melodica cantabile. La nota indicata dalla numerica deve comparire nelle voci superiori, ma non necessariamente al soprano.
\end{itemize}

\subsection{La gestione delle dissonanze nel basso numerato}

Quando la numerica indica una dissonanza (settima, ritardo, sesta eccedente):

\begin{regola}[Trattamento delle dissonanze nella realizzazione]
	\label{reg:dissonanze_realizzazione}
	\begin{enumerate}
		\item \textbf{Identificare la nota dissonante}: qual è la nota che crea la dissonanza rispetto al basso?
		\item \textbf{Verificare la preparazione}: la nota dissonante era consonante nell'accordo precedente nella stessa voce? (Obbligatorio per II7, IV7, ritardi; non necessario per V7.)
		\item \textbf{Pianificare la risoluzione}: la nota dissonante deve scendere di grado nell'accordo successivo. Quale nota dell'accordo successivo è la sua destinazione?
		\item \textbf{Evitare conflitti di risoluzione}: se due note devono risolvere contemporaneamente, assicurarsi che le loro risoluzioni non producano parallelismi o ottave.
	\end{enumerate}
\end{regola}

\section{Fase 5 — Revisione dell'elaborato}
\label{sec:fase5_numerato}

L'ultima fase è la più critica: il controllo degli errori.

\begin{regola}[Protocollo di revisione del basso numerato]
	\label{reg:revisione_numerato}
	Utilizzare sistematicamente la checklist degli errori (\ref{reg:checklist_errori}). Suggerire di rileggere ogni coppia di voci separatamente (S+B, A+T, S+A, T+B) prima di valutare l'insieme.
	
	Attenzione particolare a:
	\begin{itemize}
		\item La sensibile al soprano ha sempre risolto verso l'alto?
		\item Ogni settima si è preparata (dove necessario) e risolta in discesa?
		\item Ogni ritardo indicato nella numerica è stato realizzato nel tempo forte corretto?
		\item Il secondo rivolto è stato usato solo nei contesti ammessi (\ref{chap:quarta_sesta})?
		\item I raddoppi rispettano il principio (mai sensibile, mai settima)?
	\end{itemize}
\end{regola}

Il segreto di un buon basso non è solo l'assenza di errori, ma la \textbf{coerenza del discorso musicale} che riusciamo a costruire sopra le note date. La numerica ci vincola nella scelta degli accordi, ma ci lascia liberi nella condotta delle voci e nella qualità melodica del soprano.

\section{Esempi di realizzazione commentati}
\label{sec:esempi_realizzazione}

\subsection{Basso numerato con settima di dominante e ritardo}

\musica{
	\new Staff {
		\clef bass \key do \major \time 4/4
		\relative do, {
			do2 fa | sol2 do
		}
	}
}
Numerica: \texttt{s2 s | <7>2 <4 3>}

\textbf{Lettura:}
\begin{itemize}
	\item Battuta 1: Do (I, 5) -- Fa (IV, 5). Nessun numero = stati fondamentali.
	\item Battuta 2: Sol con 7 (V7) -- Do con 4-3 (I con ritardo 4-3).
\end{itemize}

\textbf{Decisioni di condotta:}
\begin{itemize}
	\item V7: basso Sol; sensibile (Si) al tenore o soprano; settima (Fa) deve scendere a Mi.
	\item Ritardo 4-3: il Fa (quarta sopra Do) è già presente nel V7 come settima. Viene tenuto e risolve su Mi (terza di I). Questo è il ritardo preparato dalla settima.
	\item Il soprano può avere il Do (ritardo 4, poi Mi, risoluzione 3) oppure un'altra voce porta il ritardo.
\end{itemize}

\musica{
	\partilate
	{ \key do \major \time 4/4 }
	{ mi'2 fa' | \freccia fa'4 \glissando mi'4 mi'2 }
	{ do'2 do' | re'2 do'2 }
	{ sol2 la | si2 sol2 }
	{ do2 fa, | sol,2 do2 }
	\lyricmode { "I"2 "IV" | "V7"2 "I"2 }
	\figuremode { s2 s | <7>2 <4 3>2 }
}

\subsection{Basso numerato con modulazione implicita}

\musica{
	\new Staff {
		\clef bass \key re \major \time 4/4
		\relative re, {
			re2 mi | fad2 sol | la2 si | dod1
		}
	}
}
Numerica: \texttt{s2 s | <6>2 s | s2 s | <6>1}

\textbf{Analisi:}
\begin{itemize}
	\item Tonalità: Re maggiore (\Fadiesis{}, \Dodiesis{}).
	\item Battuta 4: \Dodiesis{} con 6 $\rightarrow$ accordo di \Dodiesis{} con la terza al basso. \Dodiesis{} è la sensibile di Re. L'accordo sarà VII°6 di Re maggiore $\rightarrow$ rimaniamo in Re maggiore.
	\item Il brano sale per gradi congiunti (I-II-III-IV-V-VI-VII): Regola dell'Ottava completa ascendente in Re maggiore.
\end{itemize}

\musica{
	\partilate
	{ \key re \major \time 4/4 }
	{ fad'2 fad' | fad' sol' | la'2 si' | do''1 }
	{ re'2 re' | re' re' | fad'2 sol' | la'1 }
	{ la2 si | la si | do'2 re' | fad'1 }
	{ re,2 mi, | fad, sol, | la,2 si, | dod1 }
	\lyricmode { "I"2 "V" | "I6" "II6" | "V"2 "III" | "VII°6"1 }
	\figuremode { s2 s | <6>2 <6> | s2 s | <6>1 }
}
 % Come affrontare il basso numerato
	% ============================================================
% CAPITOLO N - DAL BASSO NUMERATO AL BASSO SENZA NUMERI
% (Capitolo ponte - NUOVO)
% ============================================================
\chapter{Dal basso numerato al basso senza numeri}
\label{chap:dal_numerato_al_senza_numeri}

Fino a questo punto del manuale, ogni esercizio di armonizzazione ci è stato offerto con un prezioso aiuto: la \textbf{numerica}, ovvero le cifre scritte sotto le note del basso che ci dicono quale accordo usare in ogni momento. La numerica è la bussola del basso figurato: indica lo stato degli accordi, segnala le dissonanze, guida la condotta delle voci.

Arrivare al basso senza numeri non significa buttare via quella bussola. Significa \textbf{averla interiorizzata} al punto da non averne più bisogno: l'obiettivo di questo capitolo è costruire la capacità di dedurre la numerica da sé, prima ancora di scrivere una nota.

\section{Che cosa cambia senza la numerica}
\label{sec:cosa_cambia}

Nel basso \textbf{figurato} (con numeri), il percorso operativo è:
\begin{center}
	Numerica $\rightarrow$ Accordo $\rightarrow$ Condotta delle voci
\end{center}

Nel basso \textbf{senza numeri}, il percorso diventa:
\begin{center}
	Analisi del basso $\rightarrow$ Numerica implicita $\rightarrow$ Accordo $\rightarrow$ Condotta delle voci
\end{center}

Sembra una piccola differenza. In realtà, il primo passo --- ``analisi del basso'' --- richiede la mobilitazione di \textit{tutte} le conoscenze acquisite: tonalità, funzioni armoniche, cadenze, progressioni, modulazioni, ritmo armonico. È un'operazione di \textbf{sintesi} che presuppone la padronanza di tutti i nuclei analitici precedenti.

\begin{osservazione}[Il basso senza numeri è un'operazione di scelta responsabile]
	\label{oss:scelta_responsabile}
	Nel basso numerato, lo studente applica regole note a una situazione determinata. Nel basso senza numeri, lo studente \textbf{sceglie} quali regole applicare e in quale contesto. Questa differenza è profonda: non si tratta di sapere di più, ma di usare quello che si sa in modo \textit{autonomo} e \textit{argomentato}.
\end{osservazione}

\section{La lettura del basso come testo musicale}
\label{sec:lettura_basso}

Il basso senza numeri è una \textbf{linea melodica} che contiene, implicite, informazioni armoniche. Saperlo leggere significa saper estrarre quelle informazioni prima ancora di sedersi al pianoforte o al foglio di carta.

\subsection{Informazioni immediate: tonalità e modo}

La prima operazione è sempre l'identificazione della \textbf{tonalità e del modo}.

\begin{regola}[Identificazione della tonalità nel basso senza numeri]
	\label{reg:identificazione_tonalita}
	Per determinare la tonalità e il modo di un basso senza numeri, si seguono questi passi:
	\begin{enumerate}
		\item \textbf{Armatura di chiave:} è il punto di partenza. Stabilisce il pool di note ``naturali'' del brano.
		\item \textbf{Nota iniziale e finale:} tipicamente la nota iniziale è la tonica o la dominante; la nota finale è quasi sempre la tonica. Questo conferma (o smentisce) l'ipotesi suggerita dall'armatura.
		\item \textbf{Alterazioni cromatiche:} la presenza di una sensibile alzata (VII grado) conferma il \textit{modo minore armonico}. Alterazioni che non appartengono alla scala di riferimento segnalano tonicizzazioni o modulazioni (cfr. \ref{chap:tonicizzazione} e \ref{chap:modulazione}).
		\item \textbf{Salto di quarta o quinta iniziale:} se il basso inizia con un salto I-V o V-I, la tonalità è subito chiara.
	\end{enumerate}
\end{regola}

\begin{esempio}[Identificazione della tonalità da un basso senza numeri]
	\label{ex:identificazione_tonalita}
	Basso da analizzare in Do minore:
	\musica{
		\new Staff {
			\clef bass \key do \minor \time 4/4
			\relative do, {
				do4 re mib fa | sol2 sold4 la | sib2 la4 sol | do1
			}
		}
	}
	Osservazioni:
	\begin{itemize}
		\item Armatura: 3 bemolli $\rightarrow$ Do minore (o \Mibemolle  maggiore).
		\item Nota iniziale e finale: Do $\rightarrow$ conferma Do minore.
		\item \Soldiesis : sensibile alzata $\rightarrow$ modo minore armonico confermato.
		\item Il \Ladiesis  in realtà è un La$\natural$ (la sensibile alzata è il Sol\diesis{}) --- in questo basso il La è la sopradominante naturale.
	\end{itemize}
\end{esempio}

\subsection{Informazioni strutturali: cadenze e frasi}

Una volta stabilita la tonalità, il secondo passo è la \textbf{segmentazione fraseologica}: trovare dove si trovano le cadenze e quindi dove sono i confini tra le frasi.

\begin{regola}[Riconoscimento delle cadenze nel basso senza numeri]
	\label{reg:riconoscimento_cadenze_basso}
	Le cadenze si riconoscono cercando le seguenti \textbf{formule armoniche tipiche} alla fine delle frasi:
	\begin{itemize}
		\item \textbf{Cadenza perfetta} (V-I): il basso sale di quarta (o scende di quinta) da V a I, entrambi in stato fondamentale.
		\item \textbf{Cadenza sospesa} (--V): il basso termina sulla nota della dominante, spesso preceduta da IV o II.
		\item \textbf{Cadenza plagale} (IV-I): il basso scende di quarta da IV a I.
		\item \textbf{Cadenza composta} (IV-V-I o II-V-I): tre note del basso che formano la formula conclusiva più completa.
		\item \textbf{Cadenza d'inganno} (V-VI): il basso sale di grado congiunto da V a VI dopo aver fatto aspettare la tonica.
	\end{itemize}
	Le cadenze si trovano quasi sempre ai bordi delle frasi di 4 misure (struttura antecedente-conseguente, cfr. \ref{chap:struttura_frase}).
\end{regola}

\subsection{Informazioni dinamiche: il ritmo armonico}

Il \textbf{ritmo armonico} è la frequenza con cui gli accordi cambiano all'interno di una battuta. È un'informazione cruciale che spesso il basso ci fornisce implicitamente.

\begin{definizione}[Ritmo armonico]
	\label{def:ritmo_armonico}
	Il \textbf{ritmo armonico} è la durata media di ogni accordo all'interno del brano. Un ritmo armonico \textit{lento} significa accordi lunghi (semibrevi o minime); un ritmo armonico \textit{veloce} significa accordi corti (semiminime o crome).
\end{definizione}

\begin{regola}[Ritmo armonico e numero di accordi per battuta]
	\label{reg:ritmo_armonico_battuta}
	In un basso senza numeri in stile corale:
	\begin{itemize}
		\item In \textbf{4/4} e \textbf{C} (alla breve): tipicamente uno o due accordi per battuta (uno per minima).
		\item In \textbf{3/4}: tipicamente uno o tre accordi per battuta.
		\item In \textbf{2/2} o \textbf{3/2}: tipicamente un accordo per minima.
		\item Ogni nota del basso che vale più di un tempo \textit{può} essere armonizzata con uno o più accordi diversi --- ma il cambio armonico deve essere argomentato.
		\item Una nota di basso che vale un solo movimento (semiminima in 4/4) di solito porta un accordo solo.
	\end{itemize}
	\textbf{Regola pratica di default:} in un esercizio scolastico, si assegna un accordo per ogni \textit{accento metrico principale} della misura, a meno che il contesto non indichi chiaramente il contrario.
\end{regola}

\section{Il processo deduttivo: da nota a funzione}
\label{sec:processo_deduttivo}

Il cuore del lavoro sul basso senza numeri è il \textbf{ragionamento funzionale}: partendo dalla nota del basso, dedurre quale funzione armonica essa esercita nel contesto in cui si trova.

\subsection{La nota del basso non è una funzione isolata}

Questo è il principio fondamentale, evidenziato anche dalle fonti metodologiche di riferimento: \textbf{una nota del basso, presa isolatamente, non esprime alcuna funzione tonale}. Il suo significato diventa chiaro solo nel confronto con le note vicine, nel contesto delle cadenze e delle progressioni riconoscibili.

\begin{esempio}[Ambiguità funzionale di un Sol nel basso]
	\label{ex:ambiguita_sol}
	La nota Sol al basso può essere:
	\begin{itemize}
		\item \textbf{V grado in Do maggiore} (dominante) $\rightarrow$ accordo di Sol maggiore (V)
		\item \textbf{I grado in Sol maggiore} (tonica) $\rightarrow$ accordo di Sol maggiore (I)
		\item \textbf{IV grado in Re maggiore} $\rightarrow$ accordo di Sol maggiore (IV) o Sol settima
		\item \textbf{VII grado in La minore} (sensibile alzata) $\rightarrow$ accordo di \Soldiesis  diminuito (VII°)
		\item \textbf{Nota di passaggio} tra La e Fa $\rightarrow$ nessun cambio armonico
	\end{itemize}
	Solo il contesto --- le note precedenti e seguenti, le alterazioni, la posizione metrica --- ci dice quale delle interpretazioni è corretta.
\end{esempio}

\subsection{Gerarchia delle informazioni contestuali}

Quando si analizza una nota del basso, le informazioni utili vanno cercate in questo ordine di priorità:

\begin{regola}[Gerarchia delle informazioni per la deduzione armonica]
	\label{reg:gerarchia_informazioni}
	\begin{enumerate}
		\item \textbf{Posizione metrica:} le note sui tempi forti portano quasi sempre accordi in stato fondamentale (cfr. \ref{reg:stato_basso}). Le note sui tempi deboli possono portare rivolti o note di passaggio.
		\item \textbf{Presenza di alterazioni:} un'alterazione cromatica segnala quasi sempre una funzione speciale: sensibile del modo minore, dominante secondaria, o cambio di tonalità.
		\item \textbf{Movimento rispetto alla nota precedente:} grado congiunto suggerisce Regola dell'Ottava o progressione; salto di quarta/quinta suggerisce cambio di pilastro armonico; salto di ottava può indicare cambio di rivolto.
		\item \textbf{Contesto cadenzale:} se la nota si trova in prossimità di un punto di cadenza già individuato, la sua funzione è determinata dalla formula cadenzale.
		\item \textbf{Progressioni riconoscibili:} se le note del basso formano un pattern di progressione (cfr. \ref{chap:progressioni_armoniche}), la funzione è determinata dalla struttura sequenziale.
	\end{enumerate}
\end{regola}

\section{Applicazione della Regola dell'Ottava al basso senza numeri}
\label{sec:regola_ottava_basso_senza_numeri}

La \textbf{Regola dell'Ottava} (cfr. Capitolo \ref{chap:funzioni_armoniche} e \ref{chap:regola_ottava_completa}) è lo strumento principale per armonizzare i movimenti scalari del basso. Nel basso senza numeri, si applica la stessa logica, ma la scelta del livello di complessità (versione preliminare, intermedia o completa) dipende dal contesto.

\begin{regola}[Applicazione della Regola dell'Ottava senza numeri]
	\label{reg:applicazione_regola_ottava_basso}
	Quando il basso procede per gradi congiunti (scala ascendente o discendente), si applicano le regole seguenti:
	\begin{itemize}
		\item \textbf{I e VIII grado}: accordo di tonica in stato fondamentale (5).
		\item \textbf{II grado}: accordo sul VII in primo rivolto (6) oppure V4/3 --- a seconda del contesto e del livello.
		\item \textbf{III grado}: accordo di I in primo rivolto (6).
		\item \textbf{IV grado}: accordo di II in primo rivolto (6) oppure II6/5.
		\item \textbf{V grado}: accordo di V in stato fondamentale (5) oppure V7. In alternativa, I6/4 cadenzale o di passaggio.
		\item \textbf{VI grado}: accordo di IV in primo rivolto (6).
		\item \textbf{VII grado}: accordo di V in primo rivolto (6) oppure V6/5.
	\end{itemize}
	Per i dettagli e le varianti, si rimanda alla tabella completa in \ref{reg:regola_ottava_completa_maggiore}.
\end{regola}

\section{Riconoscimento delle progressioni nel basso senza numeri}
\label{sec:progressioni_basso_senza_numeri}

Le \textbf{progressioni} (Capitolo \ref{chap:progressioni_armoniche}) sono strumenti particolarmente utili nel basso senza numeri, perché permettono di riconoscere un pattern e applicare un sistema di armonizzazione coerente a tutto il passaggio.

\begin{regola}[Come riconoscere una progressione nel basso]
	\label{reg:riconoscimento_progressione}
	Una progressione si riconosce quando le note del basso formano un pattern che si ripete a distanza regolare:
	\begin{itemize}
		\item \textbf{Salti ripetuti di quarta o quinta}: indicano una progressione per quinte (ascendente o discendente).
		\item \textbf{Alternanza grado congiunto-salto}: può indicare una progressione per terze con rivolti alternati.
		\item \textbf{Discesa o salita cromatica}: indica una progressione cromatica (cfr. \ref{sec:progressioni_cromatiche}).
		\item \textbf{Scalare con ritmo regolare}: può indicare una progressione scalare secondo la Regola dell'Ottava.
	\end{itemize}
	Quando si riconosce una progressione, si applica lo stesso tipo di armonizzazione a tutte le ripetizioni del modello.
\end{regola}

\section{Alterazioni cromatiche: segnali di tonicizzazione e modulazione}
\label{sec:cromatismi_basso}

Nel basso senza numeri, ogni alterazione cromatica è un \textbf{segnale} che richiede un'interpretazione attenta. Le possibilità sono quattro, in ordine di complessità crescente:

\begin{enumerate}
	\item \textbf{Sensibile del modo minore}: l'alterazione è il VII grado alzato nella tonalità in corso. Non cambia nulla di strutturale; si usa semplicemente la dominante maggiore (V o V7) con la sensibile corretta.
	\item \textbf{Cromatismo ornamentale}: l'alterazione non ha funzione armonica, ma abbellisce la linea melodica del basso. Raro negli esercizi scolastici, ma da riconoscere.
	\item \textbf{Tonicizzazione} (cfr. \ref{chap:tonicizzazione}): l'alterazione crea una dominante secondaria che ``punta'' verso un grado non-tonicale come fosse una tonica locale. L'effetto è momentaneo e la tonalità principale si ristabilisce subito dopo.
	\item \textbf{Modulazione} (cfr. \ref{chap:modulazione}): l'alterazione segna l'ingresso in una nuova tonalità, confermata da una cadenza nella nuova chiave. Il cambiamento è strutturale e duraturo.
\end{enumerate}

\begin{regola}[Come interpretare un'alterazione nel basso senza numeri]
	\label{reg:interpretazione_alterazione}
	Di fronte a un'alterazione cromatica nel basso, chiedersi nell'ordine:
	\begin{enumerate}
		\item È il VII grado alzato del modo minore in corso? $\rightarrow$ sensibile, nessun problema.
		\item Questa nota è la \textit{terza} di una triade maggiore che risolve per quinta discendente sull'accordo successivo? $\rightarrow$ tonicizzazione o dominante secondaria.
		\item La tonalità sembra cambiare e viene confermata da una cadenza? $\rightarrow$ modulazione.
		\item Se nessuna delle precedenti: cromatismo ornamentale (usare la Regola dell'Ottava con la nota alterata come accordo di passaggio).
	\end{enumerate}
\end{regola}

\section{Conferma della tonalità: strumento di analisi e di composizione}
\label{sec:conferma_tonalita}

Sia nell'analisi che nella realizzazione del basso senza numeri, la \textbf{conferma della tonalità} è un concetto operativo fondamentale. Essa indica il processo con cui si stabilisce, senza possibilità di doppia lettura, la tonalità e il modo di un brano o di una sezione.

\begin{definizione}[Conferma della tonalità]
	\label{def:conferma_tonalita}
	La \textbf{conferma della tonalità} è la presenza di quegli accordi che, insieme, permettono di identificare inequivocabilmente la tonalità e il modo del brano. Il metodo più diretto è la \textbf{cadenza composta} (IV-V-I): i tre accordi costruiti su questi gradi contengono \textit{tutte} le note della scala, rendendo impossibile qualsiasi ambiguità tonale.
\end{definizione}

\begin{regola}[Quando e dove la tonalità viene confermata]
	\label{reg:quando_conferma_tonalita}
	\begin{itemize}
		\item \textbf{All'inizio del brano}: i primi accordi stabiliscono la tonalità. Anche solo I-V o I-IV sono sufficienti per orientare l'orecchio.
		\item \textbf{Dopo una modulazione}: la nuova tonalità deve essere confermata da una cadenza autentica (V-I) nella nuova chiave. Senza questa conferma, si parla di semplice tonicizzazione.
		\item \textbf{Alla fine del brano}: la cadenza finale (di solito composta II-V-I o IV-V-I) è la conferma definitiva della tonalità di impianto.
	\end{itemize}
\end{regola}

\section{Nota del basso che vale più tempi}
\label{sec:nota_piu_tempi}

Una delle situazioni più frequenti nel basso senza numeri è quella di una nota che vale più tempi, o addirittura più misure. In questo caso, lo studente deve decidere se assegnare un solo accordo oppure far cambiare l'armonia sopra la stessa nota del basso.

\begin{regola}[Nota del basso che vale più tempi]
	\label{reg:nota_piu_tempi}
	\begin{itemize}
		\item Una nota del basso che vale \textbf{più di un tempo} può appartenere a tonalità diverse o portare più di un accordo. La scelta dipende dal contesto fraseologico: se siamo in zona cadenzale, conviene armonizzarla con l'accordo cadenzale corrispondente.
		\item Una nota lunga in posizione di \textbf{tonica} porta di norma un unico accordo di tonica.
		\item Una nota lunga in posizione di \textbf{dominante} può portare I6/4 (cadenzale) risolto su V, oppure semplicemente V per l'intera durata.
		\item Se la nota lunga non è in posizione cadenzale, si può scegliere di armonizzarla con due accordi diversi, sfruttando ad esempio un cambio di posizione o un accordo vicino (per terza o per quinta).
	\end{itemize}
\end{regola}

\section{La nota del basso che salta}
\label{sec:nota_salta}

\begin{regola}[Nota del basso che salta: principio degli stati fondamentali]
	\label{reg:nota_salta_stati_fondamentali}
	Una nota di basso che è raggiunta per \textbf{salto} (non per grado congiunto) porta quasi sempre un accordo in \textbf{stato fondamentale}. I rivolti sono rari in corrispondenza dei salti, perché il salto segnala un cambio di pilastro armonico che richiede solidità strutturale.
	
	\textbf{Unica eccezione rilevante:} il I grado in primo rivolto (I6) può comparire dopo un salto di ottava ascendente quando il basso realizza un arpeggio della tonica.
\end{regola}

\begin{errore}[Accordo rivoltato dopo un salto nel basso]
	\label{err:rivoltato_dopo_salto}
	Evitare di assegnare accordi in secondo rivolto (6/4) o terzo rivolto (2) dopo un salto nel basso. Questi accordi instabili richiedono un contesto melodico che solo il grado congiunto può fornire.
\end{errore}

\section{Dalla deduzione alla realizzazione: il flusso di lavoro}
\label{sec:flusso_lavoro_basso_senza_numeri}

Il flusso di lavoro completo per il basso senza numeri si articola in queste fasi, che riprenderemo in modo dettagliato nel Capitolo \ref{chap:metodo_basso_senza_numeri}:

\begin{enumerate}
	\item \textbf{Analisi preliminare}: tonalità, modo, alterazioni cromatiche.
	\item \textbf{Segmentazione}: frasi, cadenze, struttura del periodo.
	\item \textbf{Classificazione di ogni nota}: grado della scala, posizione metrica, tipo di movimento.
	\item \textbf{Deduzione della numerica implicita}: applicare Regola dell'Ottava, progressioni, formule cadenzali.
	\item \textbf{Scrittura delle voci superiori}: applicare le regole di condotta delle voci conosciute.
	\item \textbf{Revisione}: checklist degli errori (\ref{reg:checklist_errori}).
\end{enumerate}

Il capitolo seguente (\ref{chap:metodo_basso_senza_numeri}) sviluppa ognuna di queste fasi in modo operativo, con esempi svolti commentati e strategie per la gestione dei casi più complessi.
 % Dal basso numerato al basso senza numeri
	% ============================================================
% CAPITOLO O - METODO COMPLETO PER L'ARMONIZZAZIONE
%              DEL BASSO SENZA NUMERI
% (Capitolo finale - NUOVO)
% ============================================================
\chapter{Metodo operativo per il basso senza numeri}
\label{chap:metodo_basso_senza_numeri}

Questo è il capitolo conclusivo del percorso. Tutto ciò che è stato studiato --- fondamenti acustici, triadi, rivolti, condotta delle voci, funzioni armoniche, cadenze, dissonanze, progressioni, modulazioni --- converge qui in un \textbf{metodo operativo} per l'armonizzazione autonoma di un basso privo di numerica.

Non esiste una formula magica. Esiste però una \textbf{procedura di analisi e scelta} che, se seguita sistematicamente, riduce l'incertezza e permette di produrre realizzazioni musicalmente coerenti. Questo capitolo presenta quella procedura, la illustra su esempi completi, cataloga gli errori tipici e suggerisce strategie di scelta nei casi più ambigui.

\section{Le cinque fasi del metodo}
\label{sec:cinque_fasi}

\begin{regola}[Il metodo in cinque fasi]
	\label{reg:metodo_cinque_fasi}
	L'armonizzazione del basso senza numeri si articola in cinque fasi sequenziali:
	\begin{enumerate}
		\item \textbf{Analisi preliminare} --- Tonalità, modo, struttura metrica.
		\item \textbf{Segmentazione fraseologica} --- Frasi, cadenze, periodo.
		\item \textbf{Classificazione delle note} --- Gradi, funzioni, movimento, alterazioni.
		\item \textbf{Stesura dell'armatura} --- Scrittura delle voci superiori con rispetto delle regole di condotta.
		\item \textbf{Revisione sistematica} --- Controllo degli errori.
	\end{enumerate}
	Queste fasi sono \textit{sequenziali}: non si passa alla fase successiva prima di aver completato la precedente.
\end{regola}

\section{Fase 1 — Analisi preliminare}
\label{sec:fase1_basso}

\subsection{Tonalità e modo}

Prima di toccare carta e penna, dedicare almeno un minuto alla \textbf{lettura silenziosa} dell'intero basso. Porsi le domande:
\begin{itemize}
	\item Qual è l'armatura di chiave?
	\item La nota iniziale è la tonica o la dominante?
	\item La nota finale (di solito lunga) è la tonica?
	\item Ci sono alterazioni cromatiche? Quali e dove?
\end{itemize}

\begin{regola}[Lettura diagnostica del modo]
	\label{reg:lettura_diagnostica}
	\begin{itemize}
		\item Se il basso contiene regolarmente la \textbf{sensibile alzata} (VII grado cromatico), il brano è in modo minore armonico.
		\item Se il basso alterna il VII grado alzato e abbassato, potremmo essere in modo minore melodico: stare attenti al contesto (salita vs.\ discesa).
		\item Se il basso ha solo le note dell'armatura, ma si muove attorno alla tonica di un modo minore, siamo probabilmente in modo minore naturale (raro negli esercizi scolastici).
		\item Un accordo iniziale con la terza maggiore visibile nell'arpeggio del basso conferma il modo maggiore.
	\end{itemize}
\end{regola}

\subsection{Struttura metrica e ritmo armonico previsto}

Leggere il metro (4/4, 3/4, 3/2, ecc.) e stimare il ritmo armonico previsto. In un esercizio scolastico standard:
\begin{itemize}
	\item In \textbf{4/4}: un accordo per ogni semiminima o minima, a seconda delle indicazioni dell'esaminatore.
	\item In \textbf{3/2}: un accordo per minima (tre accordi per misura) o uno per misura intera, a seconda della musica.
	\item La regola generale è: tanti accordi quanti sono gli \textit{accenti principali} del metro.
\end{itemize}

\subsection{Identificazione dei pattern visivi}

Prima ancora di analizzare armonia, cercare \textbf{pattern visibili} nel profilo del basso:
\begin{itemize}
	\item Movimenti scalari lunghi (più di 3-4 note in senso congiunto) $\rightarrow$ Regola dell'Ottava o progressione.
	\item Salti ripetuti in senso discendente $\rightarrow$ progressione per quinte discendenti.
	\item Alternanza semitoni/toni $\rightarrow$ progressione cromatica.
	\item Salti di quarta o quinta seguiti da grado congiunto $\rightarrow$ inizio di progressione per terze.
\end{itemize}

\section{Fase 2 — Segmentazione fraseologica}
\label{sec:fase2_basso}

\subsection{Identificare le cadenze}

Le cadenze si trovano quasi sempre ai \textbf{punti di articolazione ritmica del brano}: fine di una frase di 4 misure, metà del periodo, fine del brano.

\begin{regola}[Segnali di cadenza nel basso senza numeri]
	\label{reg:segnali_cadenza}
	\begin{itemize}
		\item Una nota lunga (semibreve o minima in 4/4) dopo una frase in movimento $\rightarrow$ cadenza.
		\item Il basso che arriva alla tonica dopo un movimento di quarta ascendente o quinta discendente $\rightarrow$ cadenza perfetta.
		\item Il basso che si ferma sulla dominante $\rightarrow$ cadenza sospesa.
		\item Presenza di due note del basso V-I con ritmo regolare prima di una nota lunga $\rightarrow$ cadenza composta.
		\item Una nota alterata seguita immediatamente da un grado congiunto $\rightarrow$ possibile cadenza con sensibile o con tonicizzazione.
	\end{itemize}
\end{regola}

\subsection{Struttura antecedente-conseguente}

In un basso di 8 misure, applicare il modello del periodo:
\begin{itemize}
	\item \textbf{Misure 1-4} (antecedente): la cadenza di metà deve essere \textit{aperta} --- sospesa (V) o imperfetta (I non fondamentale, o soprano non all'ottava).
	\item \textbf{Misure 5-8} (conseguente): la cadenza finale deve essere \textit{chiusa} --- cadenza perfetta (V-I in stato fondamentale, soprano all'ottava o alla quinta).
\end{itemize}

Se il basso è di 4 misure, applicare lo stesso criterio: cadenza di metà aperta, cadenza finale chiusa.

\section{Fase 3 — Classificazione delle note}
\label{sec:fase3_basso}

Questa è la fase più impegnativa. Si passa nota per nota, annotando (idealmente sopra o sotto il basso) la funzione armonica dedotta.

\subsection{Schema di classificazione}

Per ogni nota del basso, rispondere nell'ordine:

\begin{enumerate}
	\item \textbf{Qual è il suo grado nella scala?} (I, II, III\ldots)
	\item \textbf{È su un tempo forte o debole?}
	\item \textbf{Come si muove rispetto alla nota precedente?} (salto, grado congiunto, nota tenuta)
	\item \textbf{Come si muove verso la nota seguente?} (uguale a sopra)
	\item \textbf{Appartiene a un pattern riconoscibile?} (cadenza, progressione, Regola dell'Ottava)
	\item \textbf{Quale accordo è più coerente con il contesto?}
\end{enumerate}

\subsection{Tabella di riferimento rapido}

\begin{center}
\begin{tabular}{|l|l|l|}
\hline
\textbf{Situazione} & \textbf{Soluzione di default} & \textbf{Note} \\
\hline
I grado, tempo forte & I (stato fond.) & Tonica, 5 \\
II grado, congiunto & VII6 o V4/3 & Regola dell'Ottava \\
III grado, congiunto & I6 & Regola dell'Ottava \\
IV grado, congiunto & II6 o II6/5 & Regola dell'Ottava \\
V grado, congiunto & V o V7 & o I6/4 cadenzale/passaggio \\
VI grado, congiunto & IV6 & Regola dell'Ottava \\
VII grado, congiunto & V6 o V6/5 & Regola dell'Ottava \\
Nota con salto & Stato fond. & Evitare rivolti \\
Nota cromatica alt. & Dom. secondaria & Analizzare contesto \\
Nota finale & I (stato fond.) & Cadenza perfetta \\
\hline
\end{tabular}
\end{center}

\subsection{Gestione delle tonicizzazioni nel basso senza numeri}

Quando nel basso appare un'alterazione cromatica non prevista dall'armatura, il primo passo è verificare se si tratta di una \textbf{dominante secondaria} (tonicizzazione).

\begin{regola}[Riconoscimento della dominante secondaria nel basso]
	\label{reg:dom_secondaria_basso}
	\begin{itemize}
		\item Il \Fadiesis{} nel basso di Do maggiore, in corrispondenza del II grado, indica una \textbf{dominante della dominante} (V/V): l'accordo sul Re dovrà essere Re maggiore con \Fadiesis{}.
		\item Il \Sibemolle{} nel basso di Do maggiore, sul I grado, indica una \textbf{dominante della sottodominante} (V/IV): si usa Do con settima minore (\Sibemolle{}).
		\item \Midiesis{} o \Ladiesis{} nel basso in La minore: verificare se sono la sensibile del relativo maggiore (Do) o di un altro grado tonicizzato.
	\end{itemize}
\end{regola}

\section{Fase 4 — Stesura delle voci superiori}
\label{sec:fase4_basso}

Una volta stabilita la ``numerica implicita'' (il piano armonico), si scrivono le voci superiori con le stesse regole usate per il basso figurato.

\begin{regola}[Priorità nella stesura delle voci superiori]
	\label{reg:priorita_stesura}
	\begin{enumerate}
		\item \textbf{Vettori obbligati prima di tutto}: sensibile al soprano sempre verso la tonica; settima sempre verso la risoluzione discendente. Questi movimenti non sono negoziabili.
		\item \textbf{Note comuni}: tenere nella stessa voce le note che i due accordi consecutivi hanno in comune.
		\item \textbf{Moto contrario con il basso}: quando possibile, le voci superiori vanno in direzione opposta al basso.
		\item \textbf{Grado congiunto}: le voci restanti si muovono per la distanza minima.
		\item \textbf{Cantabilità del soprano}: scegliere la posizione dell'accordo (ottava, terza, quinta) che produce la linea melodica più fluida e cantabile al soprano.
	\end{enumerate}
\end{regola}

\subsection{La scelta della posizione iniziale}

La prima decisione da prendere è quale nota del primo accordo mettere al soprano. Questo determina la posizione di partenza e influenza tutta la condotta successiva.

\begin{regola}[Scelta della posizione iniziale]
	\label{reg:posizione_iniziale}
	Se non ci sono indicazioni specifiche, la posizione di ottava (fondamentale al soprano) è la scelta più solida e neutra per iniziare. Essa favorisce la cadenza perfetta finale (dove il soprano deve arrivare alla tonica). Se il brano è vivace o aperto, la posizione di quinta può dare un inizio più arioso.
\end{regola}

\subsection{Gestione del soprano nei passaggi scalari}

Nei passaggi in cui il basso sale o scende per gradi (Regola dell'Ottava), il soprano deve anch'esso procedere per gradi congiunti o restare su note del rispettivo accordo. È un errore far saltare il soprano mentre il basso procede per gradi.

\section{Fase 5 — Revisione sistematica}
\label{sec:fase5_basso}

Dopo aver scritto l'intero elaborato, si esegue la revisione secondo la checklist degli errori (\ref{reg:checklist_errori}). Si raccomanda di effettuare la lettura:
\begin{enumerate}
	\item Soprano + Basso (verifica delle quinte/ottave parallele e nascoste tra le voci estreme).
	\item Contralto + Tenore (spesso trascurato, ma fertile di errori nascosti).
	\item Soprano + Contralto, poi Tenore + Basso (le coppie adiacenti).
	\item Tutte le voci insieme (verifica visiva della coerenza generale e dei raddoppi).
\end{enumerate}

\section{Esempio svolto n. 1: basso semplice in Do maggiore}
\label{sec:esempio1}

\subsection{Basso dato}
\musica{
	\new Staff {
		\clef bass \key do \major \time 4/4
		\relative do, {
			do2 re | mi fa | sol1 | do,1
		}
	}
}

\subsection{Analisi preliminare}

\textbf{Armatura:} nessuna alterazione $\rightarrow$ Do maggiore o La minore.
\textbf{Nota iniziale:} Do $\rightarrow$ probabilmente Do maggiore (la nota finale confermerà).
\textbf{Nota finale:} Do (lunga) $\rightarrow$ confermato Do maggiore.
\textbf{Alterazioni cromatiche:} nessuna.
\textbf{Struttura:} 4 misure, metro 4/4 (due accordi per misura suggeriti dal ritmo in minime).

\subsection{Segmentazione}

Cadenza attesa in battuta 3 (metà) e in battuta 4 (finale).
\begin{itemize}
	\item Battuta 3: Sol = V grado $\rightarrow$ possibile cadenza sospesa.
	\item Battuta 4: Do = I grado $\rightarrow$ cadenza perfetta finale.
\end{itemize}

\subsection{Classificazione}

\begin{center}
\begin{tabular}{|l|l|l|l|}
\hline
\textbf{Batt.} & \textbf{Nota basso} & \textbf{Grado} & \textbf{Accordo scelto} \\
\hline
1 & Do (forte) & I & I (stato fond., 5) \\
1 & Re & II, congiunto & VII6 (primo rivolto) \\
2 & Mi & III, congiunto & I6 (primo rivolto) \\
2 & Fa & IV, congiunto & II6 (primo rivolto) \\
3 & Sol (lunga, forte) & V & V (stato fond.) \\
4 & Do (lunga, forte) & I & I (stato fond.) \\
\hline
\end{tabular}
\end{center}

\subsection{Realizzazione}

\musica{
	\partilate
	{ \key do \major \time 4/4 }
	{ do''2 si'2 | do''2 la'2 | si'1 | do''1 }
	{ mi'2 fa'2 | do'2 re'2 | sol'1 | sol'1 }
	{ sol2 fa2 | sol2 la2 | re'1 | mi'1 }
	{ do2 re2 | mi2 fa2 | sol1 | do1 }
	\lyricmode { "I"2 "VII°6" | "I6" "II6" | "V"1 | "I"1 }
	\figuremode { s2 <6> | <6> <6> | s1 | s1 }
}

\subsection{Commento}

La realizzazione applica sistematicamente la Regola dell'Ottava per i gradi congiunti della prima e della seconda battuta. Il soprano forma una linea cantabile Do-Si-Do-La, che sale poi alla sensibile Si (battuta 3) e arriva al Do. La cadenza sospesa in battuta 3 crea la tensione attesa, risolta dalla cadenza perfetta in battuta 4 con il soprano all'ottava (Do'').

\section{Esempio svolto n. 2: basso con progressione e modulazione}
\label{sec:esempio2}

\subsection{Basso dato}
\musica{
	\new Staff {
		\clef bass \key re \minor \time 4/4
		\relative re, {
			re2 mi | fa2 sol | la1 | re,1 \bar "||"
			re2 mi | fa2 fad | sol1 | re1 \bar "|."
		}
	}
}

\subsection{Analisi preliminare}

\textbf{Armatura:} 1 bemolle (Si\bemolle{}) $\rightarrow$ Fa maggiore o Re minore.
\textbf{Nota iniziale e finale:} Re $\rightarrow$ Re minore confermato.
\textbf{Alterazione cromatica:} \Fadiesis{} alla battuta 6 $\rightarrow$ segnale importante da analizzare.
\textbf{Struttura:} 8 misure con doppia stanghetta a metà.

\subsection{Segmentazione}

\begin{itemize}
	\item Prima parte (bb. 1-4): sale da Re a La, cadenza sospesa su La (V grado di Re minore) in battuta 3, poi torna a Re (cadenza perfetta) in battuta 4.
	\item Seconda parte (bb. 5-8): stessa apertura, ma con \Fadiesis{} in battuta 6 $\rightarrow$ tonicizzazione del IV (Sol minore) o modulazione?
	\item La nota finale è Re (4 battute dopo il punto di sospensione) $\rightarrow$ rimaniamo in Re minore. Si tratta quindi di una \textbf{tonicizzazione di Sol minore} (il IV di Re minore), seguita dalla cadenza plagale finale iv--i.
\end{itemize}

\subsection{Classificazione delle note critiche}

\begin{itemize}
	\item \textbf{Fa (battuta 2 e 6)}: III grado di Re minore, armonizzato con i6; il Sol che segue (battuta 2) è il IV grado e prende II°6 (sottodominante). In battuta 6 il \Fadiesis{} (sensibile di Sol) trasforma l'accordo in \textbf{Re7 in primo rivolto}, ovvero V6/5 del IV grado: la progressione Fa-Fa\diesis{}-Sol è un cromatismo ascendente che porta alla tonicizzazione di Sol minore in battuta 7.
\end{itemize}

\subsection{Realizzazione}

\musica{
	\partilate
	{ \key re \minor \time 4/4 }
	{ fa'2 la'2 | la'2 sib'2 | la'1 | la'1 }
	{ re'2 dod'2 | re'2 mi'2 | dod'1 | re'1 }
	{ la2 sol2 | fa2 mi2 | mi1 | fa1 }
	{ re,2 mi, | fa, sol, | la,1 | re,1 }
	\lyricmode { "i"2 "V4/3" | "i6" "II°6" | "V"1 | "i"1 }
	\figuremode { s2 <4 3> | <6> <6> | s1 | s1 }
}

\musica{
	\partilate
	{ \key re \minor \time 4/4 }
	{ fa'2 la'2 | la'2 la'2 | sib'1 | la'1 }
	{ re'2 dod'2 | re'2 re'2 | re'1 | re'1 }
	{ la2 sol2 | fa2 do'2 | sib1 | fa1 }
	{ re,2 mi, | fa, fad, | sol,1 | re,1 }
	\lyricmode { "i"2 "V4/3" | "i6" "V6/5/iv" | "iv"1 | "i"1 }
	\figuremode { s2 <4 3> | <6> <6 5> | s1 | s1 }
}

\subsection{Commento}

La prima parte è una semplice progressione scalare ascendente armonizzata con la Regola dell'Ottava. La seconda parte è identica fino alla battuta 6, dove il \Fadiesis{} trasforma l'accordo sul Fa in V6/5 del IV (Re7 in primo rivolto), creando una spinta cromatica verso Sol minore. La cadenza finale è plagale (iv--i) in Re minore.

\section{Errori tipici nell'armonizzazione del basso senza numeri}
\label{sec:errori_tipici}

\begin{errore}[Armonizzare nota per nota senza visione d'insieme]
	\label{err:nota_per_nota}
	L'errore più comune dei principianti è quello di analizzare il basso nota per nota senza mai alzare lo sguardo sull'insieme. Il risultato è una successione di accordi corretti singolarmente ma senza logica fraseologica: le cadenze appaiono nei posti sbagliati, le progressioni non vengono riconosciute, le modulazioni sembrano accidentali.
	
	\textbf{Rimedio:} fare sempre prima l'analisi strutturale (fase 2) prima di assegnare accordi.
\end{errore}

\begin{errore}[Usare troppi stati fondamentali in successione scalare]
	\label{err:troppi_fondamentali}
	Usare tutti gli accordi allo stato fondamentale in un basso scalare produce un effetto ``pesante'' e crea facilmente errori di parallelismo.
	
	\textbf{Rimedio:} usare sistematicamente la Regola dell'Ottava con i suoi rivolti (6 per i gradi II, III, IV, VI, VII).
\end{errore}

\begin{errore}[Trascurare la sensibile nel modo minore]
	\label{err:sensibile_minore}
	Nel modo minore, è essenziale alzare il VII grado (sensibile) nella dominante e nei rivolti della dominante. Dimenticarsi di farlo produce una dominante minore senza forza propulsiva.
	
	\textbf{Rimedio:} ogni volta che il basso è sul V grado del modo minore, verificare che la sensibile sia alzata nelle voci superiori.
\end{errore}

\begin{errore}[Confondere tonicizzazione e modulazione]
	\label{err:tonicizzazione_modulazione}
	Trattare ogni alterazione cromatica come una modulazione produce brani che cambiano tonalità ogni misura. Al contrario, ignorare alterazioni che segnalano modulazioni reali produce brani che non hanno senso armonico.
	
	\textbf{Rimedio:} cercare sempre la cadenza di conferma. Se c'è una cadenza autentica nella nuova tonalità, è una modulazione; altrimenti, è una tonicizzazione.
\end{errore}

\begin{errore}[Non risolvere le dissonanze]
	\label{err:dissonanze_non_risolte}
	Introdurre settimi accordi o accordi di sesta eccedente senza poi far risolvere correttamente le note dissonanti.
	
	\textbf{Rimedio:} applicare sempre la regola di risoluzione discendente per grado alla settima, e il movimento obbligato della sensibile.
\end{errore}

\section{Strategie per i casi ambigui}
\label{sec:strategie_ambigue}

Non esiste sempre una sola soluzione corretta. Quando il contesto è ambiguo, usare le seguenti \textbf{strategie di scelta}:

\begin{regola}[Strategie per l'ambiguità armonica]
	\label{reg:strategie_ambiguita}
	\begin{itemize}
		\item \textbf{In caso di dubbio, privilegiare la soluzione più semplice}: accordo di tonica o di dominante invece di dominanti secondarie complesse.
		\item \textbf{Preferire la coerenza alla varietà}: se una progressione funziona, continuarla. Non cambiare strategia a metà senza una ragione armonica chiara.
		\item \textbf{Valutare sempre il soprano}: tra due soluzioni armonicamente equivalenti, scegliere quella che produce la linea melodica più cantabile al soprano.
		\item \textbf{Verificare la struttura delle cadenze}: le cadenze devono trovarsi nei punti giusti. Se un'armonizzazione produce due cadenze perfette consecutive, qualcosa non va nella pianificazione fraseologica.
		\item \textbf{La nota del basso ha funzione tonale nel suo contesto}: ricordare che una nota non ha significato da sola. Se un Sol appare dopo un La-Si-Do-Re in scala ascendente di Mi minore, quasi certamente è il III grado (Sol$\natural$) della scala melodica, non il dominante di Do.
	\end{itemize}
\end{regola}

\section{Esercizi graduati}
\label{sec:esercizi_graduati}

I seguenti bassi sono ordinati per difficoltà crescente. Per ciascuno, seguire le cinque fasi del metodo prima di scrivere le voci superiori.

\subsection{Livello 1: basso scalare in modo maggiore (senza alterazioni)}

\musica{
	\new Staff {
		\clef bass \key sol \major \time 4/4
		\relative sol, {
			sol2 la | si do' | re1 | sol,1
		}
	}
}

\textbf{Suggerimenti:} tonalità Sol maggiore; scala ascendente I-II-III-IV, poi V e tonica; applicare Regola dell'Ottava nella prima metà; cadenza sospesa (V) a metà e perfetta alla fine.

\subsection{Livello 2: basso misto (congiunto e salti) in modo minore}

\musica{
	\new Staff {
		\clef bass \key la \minor \time 4/4
		\relative la, {
			la2 la | re2 mi | fa2 mi | la,1
		}
	}
}

\textbf{Suggerimenti:} La minore; battuta 1, La ripetuta $\rightarrow$ I, possibile cambio posizione; battuta 2, Re-Mi $\rightarrow$ IV-V; battuta 3, Fa-Mi $\rightarrow$ VI-V; battuta 4, cadenza perfetta. Attenzione alla sensibile (\Soldiesis{}) nel V grado.

\subsection{Livello 3: basso con progressione per quinte discendenti}

\musica{
	\new Staff {
		\clef bass \key do \major \time 3/4
		\relative do, {
			do4 fa si, | mi la re | sol do2 | do1.
		}
	}
}

\textbf{Suggerimenti:} Do maggiore; le prime due battute formano una progressione per quinte discendenti; alternare stati fondamentali e rivolti; battuta 3 è cadenza composta IV-V-I (Do-Sol); battuta 4 è la tonica finale.

\subsection{Livello 4: basso con tonicizzazione}

\musica{
	\new Staff {
		\clef bass \key mi \minor \time 4/4
		\relative mi, {
			mi2 fad | sol la | si2 lad | si1
		}
	}
}

\textbf{Suggerimenti:} Mi minore (1 diesis); in battuta 3, il \Ladiesis{} è la sensibile di Si (V di Mi m) $\rightarrow$ tonicizzazione del V; il brano inizia e finisce in Mi minore; cadenza sospesa su Si in battuta 4 è in realtà una cadenza perfetta nella tonicizzazione di Si, ma vista in Mi minore è una cadenza sospesa.

\subsection{Livello 5: basso con modulazione a un tono vicino}

\musica{
	\new Staff {
		\clef bass \key fa \major \time 4/4
		\relative fa, {
			fa2 sol | la sib | do2 re | sib1 \bar "||"
			do2 re | mi fa | sol1 | do1 \bar "|."
		}
	}
}

\textbf{Suggerimenti:} Fa maggiore (prima parte); il Mi$\natural$ nella seconda parte (battuta 5) è la sensibile di Fa? No --- in Fa il Mi è già la sensibile naturale. Allora, il Do della seconda parte potrebbe essere il I di Do maggiore. Verificare: la seconda parte finisce su Do (tonica di Do maggiore) $\rightarrow$ modulazione a Do maggiore (tonalità della dominante di Fa). L'accordo pivot è il Sol (II di Fa M = III di Sol M oppure V di Do M) --- individuarlo e usarlo come cerniera.

\section{Dalla procedura all'intuizione}
\label{sec:dall_procedura_intuizione}

Il metodo descritto in questo capitolo è, per sua natura, una procedura esplicita e sistematica. È così che deve essere per chi si avvicina per la prima volta al basso senza numeri: un insieme di passi chiari, controllabili, verificabili.

L'obiettivo a lungo termine, tuttavia, è la \textbf{fluidità}: la capacità di leggere un basso e sentire immediatamente le funzioni armoniche, senza doverle dedurre passo per passo. Questa fluidità si acquisisce con la pratica ripetuta e riflessiva --- non con la memorizzazione delle regole, ma con la loro \textit{applicazione consapevole} su molti bassi diversi.

\begin{osservazione}[Il ruolo dell'analisi dei brani]
	\label{oss:analisi_brani}
	Un acceleratore potente dello sviluppo dell'intuizione armonica è l'\textbf{analisi dei brani reali}: corali di Bach, sonate classiche, mottetti rinascimentali. Leggere un basso analizzato (dove si sa già qual è l'armonia) e chiedersi ``come avrei dedotto questo accordo?'' è un esercizio che sviluppa la capacità di riconoscimento in modo molto più efficace di qualsiasi schema astratto.
\end{osservazione}

L'armonizzazione del basso senza numeri è la dimostrazione pratica che l'armonia è stata compresa, non solo memorizzata. Non è la fine del percorso, ma l'inizio della vera autonomia musicale.
 % Metodo di lavoro per il basso senza numeri
	
	\backmatter
	
	\chapter*{Indici}
	\addcontentsline{toc}{chapter}{Indici}
	
	\tcblistof[\subsection*]{definizioni}{Definizioni}
	
	\tcblistof[\subsection*]{regole}{Regole}
	
	\tcblistof[\subsection*]{esempi}{Esempi}
	
	\tcblistof[\subsection*]{errori}{Errori comuni}
	
	\tcblistof[\subsection*]{fisica}{Con la fisica}
	
	\tcblistof[\subsection*]{matematica}{Con la matematica}
	
	\listoffigures
	
\end{document}
